\documentclass{msthesis}
\usepackage[parfill]{parskip}    % begin paragraphs with an empty line rather than an indent
\usepackage{graphicx}
\usepackage{amssymb, amsrefs}

% Include other useful packages here, these are just examples
\usepackage{hyperref}
    \hypersetup{
        colorlinks=true,
        linkcolor=black, % Internal References
        urlcolor=blue,   % External References
        citecolor=red,   % Reference References
        }
\usepackage[table]{xcolor}
\usepackage{pgfplotstable}
\pgfplotsset{compat=1.18}

% Useful for code listings, can be ommitted if you
% do not have any
\usepackage{listings}
\usepackage{needspace}
\usepackage{color}
\usepackage{graphicx} % Required for inserting images
\usepackage{geometry}
\usepackage{enumitem}
\usepackage{pdflscape}
\usepackage{amsthm}
\usepackage{amssymb}
\usepackage{cite}
\usepackage{amsmath}
\usepackage{tikz-cd}
\usepackage{simplebnf}
\usepackage{breqn}
\usepackage{multicol}
\usepackage{mathpartir}
\tikzcdset{scale cd/.style={every label/.append style={scale=#1},
    cells={nodes={scale=#1}}}}

\usepackage{stmaryrd}
\usepackage{mathrsfs}

\newtheorem{theorem}{Theorem}[section]

\newtheorem{lemma}[theorem]{Lemma}
\newtheorem{prop}[theorem]{Proposition}
\newtheorem{corollary}[theorem]{Corollary}
\newtheorem{conjecture}{Conjecture}
\theoremstyle{definition}
\newtheorem{definition}[theorem]{Definition}
\newtheorem*{notation}{Notation}
\newtheorem*{remark}{Remark}
\newtheorem{example}[theorem]{Example}

\newcommand{\semof}[1]{\llbracket{#1}\rrbracket}

\newcommand{\hi}[1]{\textcolor{red}{#1}}
\newcommand{\greyout}[1]{\textcolor{gray}{#1}}

\newcommand{\ints}{\mathbb{Z}}
\newcommand{\nat}{\mathbb{N}}
\newcommand{\mbC}{\mathbb{C}}
\newcommand{\mcB}{\mathcal{B}}
\newcommand{\mcA}{\mathcal{A}}
\newcommand{\var}{\mathbb{V}}
\newcommand{\mcC}{\mathcal{C}}
\newcommand{\mcD}{\mathcal{D}}
\newcommand{\mcDo}{\mathcal{D}_0}
\newcommand{\Set}{\ensuremath{\mathrm{Set}}}
\newcommand{\ol}[1]{\overline{#1}}
\newcommand{\colim}{\varinjlim}
\newcommand{\naturalto}{
  \mathrel{\vbox{\offinterlineskip
    \mathsurround=0pt
    \ialign{\hfil##\hfil\cr
      \normalfont\scalebox{1.2}{.}\cr
      $\longrightarrow$\cr}
  }}
}
\newcommand{\fv}[1]{\mathsf{FV}(#1)}
\newcommand{\Hom}{\mathrm{Hom}}

\newcommand{\Lmod}{\mathscr{L}}
\newcommand{\Amod}{\mathscr{A}}
\newcommand{\vars}{\mathbb{V}}
\newcommand{\wars}{\mathbb{W}}

% Define an environment for code listings, modify to match
% the language you are using and your formatting preferences
\lstset{frame=tb,
  language=Go,
  aboveskip=3mm,
  belowskip=3mm,
  showstringspaces=false,
  columns=flexible,
  basicstyle={\small\ttfamily},
  numbers=left,
  breaklines=true,
  breakatwhitespace=true,
  tabsize=3
}

\title{ISOMORPHISMS OF ALGEBRAIC DATA TYPES}
\author{Ben Lenox}
\department{Computer Science}
\degree{Master of Science}
\date{April 1st, 2026}

\advisor{Andrew Polonsky}{Ph.D.}
\member{Patricia Johann, Ph.D.}
\member{Pierre Cagne, Ph.D.}
\departmentchair{Jay Fenwick, Ph.D.}
\gradschool{Ashley Colquitt, Ph.D.}{Associate Vice Provost and Dean}

\thesismonth{May}
\thesisyear{2026}

\credential{B.S.}{Appalachian State University}

\begin{document}

\maketitle
\makesignature
\makecopyright

\begin{abstract}
  Algebraic data types are a central feature of functional programming languages like Haskell, OCaml, Agda, and Lean.  Languages with support for these types let the user define new types from sums, products, and recursion.  In this setting, it is important to determine when two types should be considered equivalent.  Syntactic equality is too restrictive of a definition of type equivalence.  Type isomorphism is a more flexible notion of equivalence that captures when two types contain the same data, but in a different form.  Type isomorphism has many practical benefits: it allows a programmer to work with different but equivalent representations of a type interchangeably, facilitates the transport of programs and proofs across isomorphic types, and enhances efficiency in programming language and library design.  
  
  This thesis investigates type isomorphisms between algebraic data types from both categorical and algebraic perspectives.  We examine these types as fixed points of functors and study the algebraic structure induced by sums, products, and recursive definitions.  Building on the work of Fiore and Leinster, we characterize a class of type isomorphisms by reducing the problem of determining isomorphism to determining equality in an associated algebraic structure.  In addition to this algebraic characterization of type isomorphisms, we study interpretations of algebraic data types in a generic categorical setting rather than restricting attention to any particular category (e.g., $\Set$).  We determine additional type isomorphisms from the structure of the least fixed points that interpret recursive types, rather than from properties of the surrounding category.

  These results provide an account of type isomorphisms that connects algebraic reasoning and categorical semantics, showing how equivalences between types arise and how they can be identified constructively.  This work supports systematic reasoning about equivalence of data representations and contributes to a broader understanding of the mathematical structure underlying functional programming languages.
\end{abstract}

\begin{acknowledgements}
I am deeply grateful to my advisor, Dr. Andrew Polonsky, for his outstanding guidance and support throughout this degree and the writing of this thesis.  He has shaped nearly everything I know about the subfield in which I work.  At every stage of this degree program, he has offered thoughtful advice, whether on technical questions, research direction, or broader academic concerns.  Our conversations, both about research and otherwise, played a large role in sustaining my motivation throughout the challenges of graduate studies.  I am especially thankful for his patience, generosity, and steady support over the years.  I am extremely fortunate to have worked with him.

I would like to thank Dr. Patricia Johann for the time and attention she has invested in my development over the course of this degree.  I am thankful for her careful feedback and high standards, particularly with regard to writing and presentation, which have helped me to better appreciate the level of clarity and precision expected in our field.  I am especially grateful for the opportunity to work with her as a research assistant under her NSF grant this past summer, through which I was able to contribute to a collaborative research project.

I would like to thank the members of my thesis committee, Dr.~Andrew Polonsky, Dr.~Patricia Johann, and Dr.~Pierre Cagne, for their time and careful consideration of this work.

Finally, I would like to thank my parents, who were supportive and encouraging throughout my time at Appalachian State University, especially during my research and the writing of this thesis.
\end{acknowledgements}

\tableofcontents
\listoffigures

% This starts the actual chapters of the thesis
\chapters

\chapter{Introduction}

Functional programming languages like Haskell, OCaml, Agda, and Lean are known for their expressive and mathematically grounded type systems, which are often fundamental to program design.  Type systems' most basic use is to detect errors, but they also guide documentation, program structure, and abstraction.  The study of typed programming languages has also revealed connections between programming languages, logic, category theory, and universal algebra.  In particular, types can be understood as mathematical objects that come equipped with a great deal of structure, including relationships with other types.  Programs are often understood as proofs or morphisms between types.

The type systems of languages like Haskell and OCaml provide built-in constructs for forming product types (tuples or records), sum types (variants), and function types.  Product types represent data that contain components from multiple types at once, e.g., a pair containing an integer and a string.  Sum types represent a choice between several possible types of data, e.g., a value that may be an integer or it may be a string.  Function types describe transformations from one type to another.  These basic constructs allow a functional programmer to build complex data structures and abstractions from these simple components.

Modern functional programming languages have type systems that support many more advanced features than those just mentioned.  For example, parametric polymorphism allows functions to operate uniformly over many types. This allows programmers to reuse abstractions such as lists, trees, and other container types, while also enabling the definition of generic functions whose behavior does not depend on the particular types to which they are applied.  Algebraic data types allow programmers to describe recursive data structures in a concise and precise way.  These types are of particular interest to us in this work.  Pattern matching provides the functional programmer with a mechanism for deconstructing these and similar data types.  In Haskell, type classes enable ``ad-hoc'' polymorphism, allowing types to support certain mechanisms shared between types in a type class.  Collectively, these features make the type systems of functional languages a rich, powerful, and expressive tool for describing the structure and behavior of programs.

Given this richness, it is natural to ask when two types should be considered ``the same''.  One might choose to consider types as equivalent when they share exactly the same syntax.  However, this is often too restrictive of a definition of ``sameness'' to reason about programs: two types may represent the same data but be written in two different ways.  For instance, consider a pair type, each element of which contains a value from type $A$ together with a value from type $B$.  This type intuitively carries the same information as a pair containing $B$ and $A$, but these types are syntactically distinct.  Similarly, a pair $(a, (b, c)) \in A \times (B \times C)$ whose right component is itself a pair contains the same information as $((a, b), c) \in (A \times B) \times C$ up to reassociation.

In order to make this intuition precise, we need a more flexible notion of type equivalence than syntactic equality.  One such notion is \emph{type isomorphism}.  Two types $A$ and $B$ are said to be isomorphic if there exist functions $f : A \to B$ and $g : B \to A$ such that composing them in either direction results in an identity function. Intuitively, this means that values of type $A$ can be converted into values of type $B$ and back again without any information being lost.  From the perspective of functional programming, a type isomorphism means that two types contain exactly the same information, merely structured differently.

Products and coproducts obey many familiar algebraic laws when we consider types up to isomorphism.  For example, the product type $A \times B$ is isomorphic to $B \times A$, reflecting the commutativity of pairs of types.  Similarly, products are associative up to isomorphism.  Similar relationships hold for sum types.  These properties correspond to certain algebraic identities --- algebraic data types form an algebra of types.

These relationships between types can be understood more formally from the perspective of category theory.  In this context, types correspond to objects, and programs and functions correspond to morphisms.  Isomorphisms represent invertible functions or programs whose computation can be ``undone''.

Recognizing that types are isomorphic itself can be useful, but constructive proofs of isomorphisms are even more helpful.  A constructive proof of a type isomorphism is an algorithm for generating the isomorphism which it proves exists: it gives an explicit isomorphism.  From a functional programming perspective, such an isomorphism is a function (and its inverse) that allows the programmer to work with different representations of the same information interchangeably.

As previously stated, one important use of type isomorphisms is to transform the representation of data.  Programs may use different internal representations of data structures for efficiency, modularity, or ease of use for the programmer.  Explicit isomorphisms allow the program to work with a different representation of some data for the above-mentioned reasons while preserving compatibility with existing code.  Isomorphisms are invertible: they can transform data of type $A$ into data of type $B$, do some work with it, then transform resulting data of type $A$ back into data of type $B$.  Explicit isomorphisms may allow a programmer to avoid both duplicating logic and introducing inconsistencies.

Another use of type isomorphisms is in library design for programming languages.  When types are isomorphic, algorithms written for one type can be transported along the isomorphism to algorithms for the isomorphic type.  In this way, algorithms can be reused.  Whatever representation of a type is the most convenient for the programmer to work with (among the set of isomorphic types) is the one the programmer may choose to work with.  

Explicit type isomorphisms are also useful for programming language design and implementation.  The abstract syntax tree of a programming language is itself a data type.  An explicit isomorphism between such types provides a translation of one syntax tree into an equivalent syntax tree.  That is, an explicit type isomorphism can provide translations between different but equivalent syntactic presentations of a programming language.

Lastly, explicit type isomorphisms can support formal reasoning about programs.  Proofs that a type has a particular property can often be transported to proofs that an isomorphic type has the same property.  In interactive theorem provers like Agda and Lean, this technique is essential for reasoning about types.  

In this work, we investigate the type isomorphisms that exist between algebraic data types.  We study these types both categorically, as objects that interpret them in a suitable category, and algebraically, where they form expressions that can be manipulated using equational reasoning.  We bring together the tools from category theory and algebra to analyze when two algebraic data types represent the same underlying structure.  

Our main contribution is an extension of work by Marcelo Fiore and Tom Leinster, providing a characterization of type isomorphisms for a class of algebraic data types arising as fixed points of polynomial functors. In particular, we show that, under some conditions, the question of whether two types are isomorphic can be reduced to equality in an associated algebraic structure.  In this proof, we draw on techniques from universal algebra, of which ring theory and semiring theory are important instances.  This provides a decidable and accessible criterion for type isomorphism, connecting the structure of types to well-understood algebraic reasoning.

We further investigate the type isomorphisms that arise under interpretation in a general categorical setting.  Instead of interpreting ADTs in the category of sets, as is commonly done, we consider interpretations across a range of categories that support interpretations of the relevant type constructors.  Working at this level of generality clarifies which isomorphisms are purely structural (that is, they arise from the form of the types themselves) and which depend on additional properties of the chosen context for interpretation.  The benefit of this approach is that our results hold in many semantic settings.

Much of the material in this document is already present in the literature or is folklore. To keep the exposition self-contained, we include our own proofs.  However, we assume the reader is familiar with basic concepts from category theory, equational logic, and abstract algebra.

The rest of this chapter gives a brief overview of the structure of the remainder of this document.  In Chapter \ref{sec:catprelim}, we go over some categorical preliminaries that are necessary to understand the rest of this work.  In particular, this chapter provides the necessary framework for interpreting ADTs.  We first give some useful lemmas and theorems about isomorphisms arising from colimits with some shared structure. Next, we give the categorical methodology for interpreting recursively defined types: we define the $LFP$ functor.  We prove that the $LFP$ together with a specified cocone morphism is an initial object in a category of $F$-algebras.  Lastly, we prove two theorems that will be useful in establishing type isomorphisms for recursively defined types. First, we prove that iterated applications of $LFP$ can be combined into a single application.  Second, we prove that applying $F$ to the least fixed point of $G \circ F$ yields, up to isomorphism, the least fixed point of $F \circ G$.

In Chapter \ref{sec:eqtheory}, we develop the necessary background in universal algebra.  We introduce first-order signatures, $\Sigma$-algebras, initial and free algebras, quotient algebras, and equational theories together with their models.  We present each notion together with its universal property.  As part of this development, we explain how elements of free algebras are interpreted in arbitrary algebras over the same signature.  Later in this chapter, these abstract constructions are instantiated concretely in the familiar settings of rings and semirings.  We show that the free ring and the free semiring over a set $X$ are $\ints[X]$ and $\nat[X]$, respectively.  This chapter culminates in a proposition relating quotients of free rings and semirings to algebras that model the ring axioms as well as additional equations.

Chapter \ref{sec:ADTs} presents the main contributions of this work.  We begin by introducing a syntax and semantics for algebraic data types, along with the notion of rig isomorphisms, which are isomorphisms induced by the algebraic structure formed by algebraic data types.  To this structure, we adjoin additional isomorphisms that are induced by the semantics of recursive types, which we call fixed point isomorphisms.  We then discuss a paper by Andreas Blass that discovered a connection between algebra, complex numbers, and the isomorphisms that rig isomorphisms induce.  This is followed by a review of a paper by Marcelo Fiore and Tom Leinster that provides a constructive proof of exactly when this connection holds.  We give an example illustrating how Fiore and Leinster's proof can be used to generate such an isomorphism.  Finally, we present the primary results of this thesis: an extension of Fiore and Leinster's work to handle isomorphisms involving multiple recursively defined algebraic data types simultaneously.  To conclude, we show how our proof can be applied to construct an explicit instance of such an isomorphism.

In the last substantive chapter of this paper, Chapter \ref{sec:murules}, we study additional isomorphisms between ADTs that go beyond the ones induced by algebraic structure of types.  Using the results from Chapter \ref{sec:catprelim} we show that the interpretations of certain syntactic algebraic data types are always isomorphic.  In particular, the theorems that state that iterated applications of $LFP$ can be collapsed to a single application, and that applying $F$ to the least fixed point of $G \circ F$ yields the least fixed point of $F \circ G$, apply directly to the interpretations of recursively defined algebraic data types.  We discuss the additional algebraic structure that extending the concept of rig isomorphisms in this way might introduce.

Finally, the paper closes with a discussion of the broader implications of our results and potential directions for future work.

\newpage 

\chapter{Categorical Preliminaries}\label{sec:catprelim}

Category theory provides a useful framework for interpreting programming languages and their types as mathematical objects.  In this document, we make heavy use of this framework to reason about certain recursive user-defined data types called algebraic data types that exist notably in functional languages such as Haskell and OCaml, but also in multi-paradigm programming languages like Rust and Swift.  Category theory allows us to work directly with the structure inherent in these data types rather than with pure syntax.

This chapter covers some categorical foundations that are used in the remainder of this document to prove our main theorems about our data types of interest.  In particular, we lay the groundwork for the initial algebra semantics of algebraic data types.  Also in this chapter, we discuss some theorems that are useful for proving the existence of data type isomorphisms based on the categorical background for data type semantics introduced in this chapter.  We assume the reader has basic knowledge of category theory including the following definitions:
\begin{multicols}{2}
\begin{itemize}[nosep]
    \item category
    \item functor
    \item natural transformation
    \item initial object
    \item terminal object
    \item isomorphism
    \item monomorphism
    \item epimorphism
    \item (co)product
    \item (co)cone
    \item (co)limit
    \item (co)complete
    \item (co)continuous
    \item Yoneda lemma
\end{itemize}
\end{multicols}
We now provide some notation that is used throughout this chapter.

We employ the usual notations from category theory.  For example, we write $f \circ g$ for the composition of morphisms $f$ and $g$ and $id_X$ for the identity morphism on object $X$.  The set of morphisms from object $A$ to object $B$ in category $\mcC$ is denoted $\Hom_\mcC(A,B)$.

We write $\omega$ to mean the first countable ordinal with its reflexive partial ordering, viewed as a category.

The category of functors with domain $\mcC$ and codomain $\mcD$ is denoted $\mcD^\mcC$.  The category of $\omega$-cocontinuous functors with domain $\mcC$ and codomain $\mcD$ is denoted $[\mcC, \mcD]$.

$\eta : F \naturalto G$ denotes a natural transformation from a functor $F$ to a functor $G$.

The set of cocones for a functor $F$ with object component $A$ is denoted $Cocone(F, A)$.

The colimit of a functor $F$ is denoted as $\colim(F)$.  

%If $F : \mcC \times \mcC \to \mcD$ is a bifunctor, we write $\colim_{X \in \mcC}(F(X, -))$ to mean the colimit of $F$'s curried counterpart, when passed the argument $X$, in the functor category $\mcD^\mcC$.

We write $\lambda X.F(X)$ to mean the functor that maps $X$ to $F(X)$.
\newpage

\section{Colimits}

In this section, we introduce some properties of colimits that will be useful in the remainder of this chapter.  First, we introduce the concept of a cofinal diagram, a condition under which a subdiagram and its parent diagram will necessarily share a colimit.  Next, we prove that an iterated colimit is the same up to isomorphism as a colimit in a product category.

\subsection{Colimits of cofinal diagrams}

We now prove a sufficient condition for a composition $D \circ H$ of functors $D$ and $H$ to have the same colimit as $D$ up to isomorphism.

The following definition is taken from ``Locally Presentable and Accessible Categories'', by Adamek and Rosicky \cite{Adamek_Rosicky_1994}. The subsequent facts are folklore, but for completeness we provide our own proofs.

\begin{definition}\label{def:cofinal}
    A functor $H : \mcDo \to \mcD$ is \emph{cofinal} if for every object $d \in \mcD$,

    \begin{enumerate}[label=(\alph*),nosep]
        \item there exists a morphism $f : d \to H d_0$ for some object $d_0 \in \mcDo$
        \item given two such morphisms $f : d \to H d_0$ and $f' : d \to H d_0'$, there exist an object $\bar{d_0} \in \mcDo$ and morphisms $g : d_0 \to \bar{d_0}$ and $g' : d_0' \to \bar{d_0}$ such that
        \[ H(g)\ \circ f = H(g')\ \circ f' \]
    \end{enumerate}
\end{definition}

We will show that given any functor $D : \mcD \to \mcC$ and a cofinal functor $H : \mcD_0 \to \mcD$, the colimit of $D \circ H$ is isomorphic to the colimit of $D$.  We prove this result by making use of several supporting lemmas.  Fix such functors $D$ and $H$ throughout the rest of this section.  

\begin{lemma}\label{lemma:coconeDtoDH}
    Every cocone on $D$ is a cocone on $D \circ H$.
\end{lemma}
\begin{proof}
    Let $(C, c_d : D d \to C)_{d \in D}$ form a cocone for $D$.  Given any morphism $f : d_0 \to d_0'$ in $\mcD_0$, $H(f)$ is a morphism in $D$, and by the cocone property
    \[ c_{Hd_0'} \circ D(H(f)) = c_{Hd_0} \]
    Thus $(C, c_{Hd_0} : H d_0 \to C)_{d_0 \in \mcDo}$ forms a cocone for $D \circ H$.
\end{proof}

Lemma \ref{lemma:coconeDtoDH} describes a map from cocones of $D$ to cocones of $D \circ H$. We henceforth call this function $\Phi_C$, where $\Phi_C : Cocone(D, C) \to Cocone(D \circ H, C)$.

\begin{lemma}\label{lemma:unique_cofinal_morphism}
    Let $(C, c_{d_0} : D(H d_0) \to C)_{d_0 \in \mcDo}$ form a cocone for $D \circ H$. Then given any morphisms $f_d : d \to H d_0$ and $f'_d : d \to H d_0'$ for some objects $d \in \mcD$ and $d_0, d_0' \in \mcDo$,
    \begin{equation}\label{eq:unique_cofinal_morphism}
        c_{d_0} \circ D(f_d) = c_{d_0'} \circ D(f_{d}')
    \end{equation}
\end{lemma}

\begin{figure}[tbp]
    \centering
    \begin{tikzcd}
	d & d && {D d} && {D\ (H d_0)} \\
	\\
	{H d_0} & {H d_0'} && {D (H d_0')} && {D(H \bar{d}_0)} \\
	\\
	&&&&&&& C
	\arrow["{f_d}", from=1-1, to=3-1]
	\arrow["{f_d'}", from=1-2, to=3-2]
	\arrow["{D(f_d)}", from=1-4, to=1-6]
	\arrow["{D(f_d')}"', from=1-4, to=3-4]
	\arrow["{D(H(g))}", from=1-6, to=3-6]
	\arrow["{c_{d_0}}"', bend left=20, from=1-6, to=5-8]
	\arrow["{D(H(g'))}"', from=3-4, to=3-6]
	\arrow["{c_{d_0'}}"', bend right=14, from=3-4, to=5-8]
	\arrow["{c_{\bar{d}_0}}"', from=3-6, to=5-8]
    \end{tikzcd}
    \caption{A diagram showing that each choice of $f_d : d \to H d_0$ results in the same morphism $c_{d_0} \circ D(f_d)$.}
    \label{fig:cocone_DHtoD}
\end{figure}

\begin{proof}
    By Definition \ref{def:cofinal} condition (b), there exists an object $\bar{d}_0 \in \mcDo$ and morphisms $g : d_0 \to \bar{d}_0$ and $g' : d_0' \to \bar{d}_0$ such that
    \begin{equation}\label{eq:condb}
        H(g) \circ f_d = H(g') \circ f_d'
    \end{equation}
    
    
    Given such a $\bar{d}_0$, $g$, and $g'$, \eqref{eq:unique_cofinal_morphism} follows by 
    commutativity of the diagram displayed in Figure \ref{fig:cocone_DHtoD}:
\begin{align*}
    c_{d_0} \circ D(f_d) &= c_{\bar{d}_0} \circ D(H g) \circ D(f_d) && \textsf{(Because $(C, c_x)_{x \in D_0}$ forms a cocone for $D \circ H$)} \\
    &= c_{\bar{d}_0} \circ D(H g \circ f_d) && \textsf{(Because $D$ is a functor)} \\
    &= c_{\bar{d}_0} \circ D(H g' \circ f_d') && \textsf{(By equation \eqref{eq:condb})} \\
    &= c_{\bar{d}_0} \circ D(H g') \circ D(f_d') && \textsf{(Because $D$ is a functor)} \\
    &= c_{d_0'} \circ D(f_d') && \textsf{(Because $(C, c_x)_{x \in D_0}$ forms a cocone for $D \circ H$)}
\end{align*}

Thus, any choice of $f_d$ results in the same morphism $c_{d_0} \circ D(f_d)$.
\end{proof}


\begin{lemma}\label{lemma:coconeDHtoD}
    Every cocone on $D \circ H$ induces a unique cocone on $D$.
\end{lemma}
\begin{proof}
    Let $(C, c_{d_0} : D(H\ d_0) \to C)_{d_0 \in \mcDo}$ form a cocone for $D \circ H$.  
    
    To define a cocone on $D$ from the given cocone, we must first define morphisms from $D d$ to $C$ for each object $d \in \mcD$.  By Definition \ref{def:cofinal} condition (a), for each object $d \in \mcD$, there exists at least one morphism from $d$ to $H d_0$ for some object $d_0 \in \mcDo$.  Let $f_d : d \to H d_0$ be an arbitrary choice of such a morphism.  
    
    We want to show that $(C, c_{d_0} \circ D(f_d) : D\ d \to C)_{d \in \mcD}$ is a cocone for $D$.  Note that by Lemma \ref{lemma:unique_cofinal_morphism}, for any object $d \in \mcD$, any choice of $f_d$ will result in the same morphism $c_{d_0} \circ D(f_d)$.

    We now must show that $(C, c_{d_0} \circ D(f_d) : D\ d \to C)_{d \in \mcD}$ is, in fact, a cocone.  That is, given any $h : d \to d'$ in $\mcD$, we must show
    \[ c_{d_0} \circ D(f_d) = c_{d_0'} \circ D(f_{d'}) \circ D(h)\]
    or, equivalently,
    \[  c_{d_0} \circ D(f_d) = c_{d_0'} \circ D(f_{d'} \circ h) \]

    The above equation follows from Lemma \ref{lemma:unique_cofinal_morphism} with $f_d = f_d$ and $f_d' = f_{d'} \circ h$.  
    
    Thus, $(C, c_{d_0} \circ D(f_d) : D\ d \to C)_{d \in \mcD}$ is a cocone for $D$.
\end{proof}

Lemma \ref{lemma:coconeDHtoD} describes a map from cocones of $D \circ H$ to cocones of $D$. We henceforth call this function $\Psi_C$, where $\Psi_C : Cocone(D \circ H, C) \to Cocone(D, C)$.

\begin{lemma}\label{lemma:cofinal_cocone_iso}
    Given any object $C \in \mcC$, there is an isomorphism between cocones, natural in $C$:
    \[ Cocone(D, C) \simeq Cocone(D \circ H, C) \]
\end{lemma}

\begin{proof}
    We have a map $\Phi_C : Cocone(D, C) \to Cocone(D \circ H, C)$ and a map $\Psi_C : Cocone(D \circ H, C) \to Cocone(D, C)$.  We will show that composing these two functions both ways results in the identity function on cocones of $D$ and $D \circ H$, respectively.
    \begin{description}
        \item[(1) $\Psi_C \circ \Phi_C = id_{Cocone(D, C)}$: ]
        Let $(C, c_{d} : D d \to C)_{d \in \mcD}$ be a cocone for $D$.  Applying $\Phi_C$ gives the cocone $(C, c_{H d_0} : D(H d_0) \to C)_{d_0 \in \mcDo}$ for $D \circ H$, i.e., restricts the morphism components of the given cocone to objects in the image of $H$.  When we apply $\Psi_C$ to $(C, c_{H d_0})_{d_0 \in \mcDo}$, we get the cocone $(C, c_{Hd_0} \circ D(f_d) : D\ d \to C)_{d \in \mcD}$ for $D$, where each $f_d$ is given by condition (a) of Definition \ref{def:cofinal}.  Recall that the cocone defined by $\Psi$ is independent of our choices of $f_d$.

        By the cocone property of $(C, c_{d})_{d \in \mcD}$, we have that
        \[ c_{H d_0} \circ D(f_d) = c_d \]
        Therefore, $\Psi_C \circ \Phi_C = id_{Cocone(D, C)}$.

        \item[(2) $\Phi_C \circ \Psi_C = id_{Cocone(D \circ H, C)}$: ]
        Let $(C, c_{d_0} : D(H d_0) \to C)_{d_0 \in \mcDo}$ be a cocone for $D \circ H$.  Applying $\Psi$ to $(C, c_{d_0} : D(H d_0) \to C)_{d_0 \in \mcDo}$ gives the cocone $(C, c_{d_0} \circ D(f_d) : D\ d \to C)_{d \in \mcD}$, where $f_d : d \to H d_0$ is given by condition (a) of Definition \ref{def:cofinal}. 

        % Applying Lemma \ref{lemma:unique_cofinal_morphism} with $f_d = f_d$ and $f_d' = H(id_{d_0})$ shows that for all $f_d$,
        % \[ c_{d_0} \circ D(f_d) = c_{d_0} \circ D(H(id_{d_0})) = c_{d_0} \circ D(id_{H d_0}) = c_{d_0} \circ id_{D(H(d_0))} = c_{d_0} \]
        
        Lemma \ref{lemma:unique_cofinal_morphism} proves that any choice of morphism $f_d : d \to H d_0$ results in the same morphism component of $(C, c_{d_0} \circ D(f_d) : D d \to C)_{d \in \mcD}$ at object $d$.  In particular, for $d = H d_0$, the component of $\Psi((C, c_{d_0})_{d_0 \in \mcDo})$ at $d$ is equal to 
        \[ c_{d_0} \circ D(f_d) = c_{d_0} \circ D(H(id_{d_0})) = c_{d_0} \circ D(id_{H d_0}) = c_{d_0} \circ id_{D(H(d_0))} = c_{d_0} \]
        % Consider the set of morphisms $\{c_{d_0} \circ D(f_d)\ |\ d = H d_0\ \textsf{for some $d_0 \in \mcDo$}\}$. Certainly $H(id_{d_0})$ is a morphism from $H d_0$ to $H d_0$, so let $f_d = H(id_{d_0})$ for each component of this set.  Thus, we can consider this set to be any of the equal sets given by
        % \begin{align*}
        %     &\{c_{d_0} \circ D(id_{d_0})\ |\ d = H d_0\ \textsf{for some $d_0 \in \mcDo$}\} \\
        %     &= \{c_{d_0} \circ id_{Hd_0}\ |\ d = H d_0\ \textsf{for some $d_0 \in \mcDo$}\} \\
        %     &= \{c_{d_0}\ |\ d = H d_0\ \textsf{for some $d_0 \in \mcDo$}\}
        % \end{align*}
        Applying $\Phi_C$ to the cocone $(C, c_{d_0} \circ D(f_d) : D\ d \to C)_{d \in \mcD}$ restricts the morphism components of the cocone to only the morphisms with a domain in the image of $D \circ H$.  We just showed that each of these components is the original component $c_{d_0}$ of $(C, c_{d_0} : D(H\ d_0) \to C)_{d_0 \in \mcDo}$, thus $\Phi_C \circ \Psi_C = id_{Cocone(D \circ H, C)}$.

        \item[(3) Naturality: ]
        Given $\gamma : C \to C'$, the map $Cocone(F,C) \to Cocone(F,C')$ is defined by post-composition.
       
        For $\Phi_C$, the naturality is immediate by construction.

        For $\Psi_C$, let $g : C \to C'$ and, for the remainder of this proof, let $(g \circ -)$ denote the pointwise post-composition of every morphism component of a cocone with the morphism $g$.  We need to show that the diagram shown in Figure \ref{fig:natpsi} commutes.

        \begin{figure}[htbp]
        \centering{
            \begin{tikzcd}
            	{Cocone(D \circ H, C)} && {Cocone(D, C)} \\
            	\\
            	{Cocone(D \circ H, C')} && {Cocone(D, C')}
            	\arrow["{\Psi_C}", from=1-1, to=1-3]
            	\arrow["{(g \circ -)}"', from=1-1, to=3-1]
            	\arrow["{(g \circ -)}", from=1-3, to=3-3]
            	\arrow["{\Psi_{C'}}"', from=3-1, to=3-3]
            \end{tikzcd}
            }
            \caption{The naturality square for the natural transformation $\Psi$}\label{fig:natpsi}
            \end{figure}    

        \newcommand{\ccone}{\mathscr{C}}

        Let $\ccone = (C, c_{d_0} : D(H d_0) \to C)_{d_0 \in \mcDo}$ be a cocone for $D \circ H$.
        Then 
        \begin{align*}
        (g \circ {-}) (\Psi_C(\ccone))
        &= (C',g \circ c_{d_0} \circ D(f_d))_{d\in \mcD}\\
        &= \Psi_{C'}(C',g \circ c_{d_0})\\
        &= \Psi'_{C'}((g \circ {-}) \ccone)
        \end{align*}
        The post-composition of $g$ does not affect the choice of $d_0$ for $d$, per Lemma \ref{lemma:unique_cofinal_morphism}.  So, $\Psi_C$ is natural in $C$.
    \end{description}

    The composition $\Phi_C \circ \Psi_C$ is the identity function on the set $Cocone(D \circ H, C)$ and the composition $\Psi_C \circ \Phi_C$ is the identity function on the set $Cocone(D,  C)$.  Therefore, $\Psi_C$ and $\Phi_C$ form a natural isomorphism between these two sets.
\end{proof}

With the above lemmas, we prove our main theorem about cofinal diagrams.

\begin{theorem}\label{theorem:cofinal}
    Let $\mcDo$, $\mcD$, and $\mcC$ be categories.  Let $D : \mcD$ be a functor and $H : \mcDo \to \mcD$ be a cofinal functor.  Then
    \[ \varinjlim D \simeq \varinjlim D \circ H \]
\end{theorem}
\begin{proof}
    Every morphism $g$ out of a cocone $E$ induces a cocone structure on the codomain of $g$ by post-composition of each of the morphism components of $E$ with $g$.  That is, the set $\Hom_\mcC(L, C)$ naturally injects into $Cocone(D, C)$.
    \[ \Hom_\mcC(E, C) \hookrightarrow Cocone(D, C) \]
    A colimit for a functor $D : \mcD \to \mcC$ is a cocone $L$ for which such an inclusion is an isomorphism for all objects $C$.  That is, it is a cocone $L$ which, for all objects $C \in \mcC$, morphisms from $L$ to $C$ are in bijective correspondence with cocones with object component $C$.
    
    By Lemma \ref{lemma:cofinal_cocone_iso}, we have for all objects $C \in \mcC$,
    \[ \Hom_\mcC(\varinjlim D, C) \simeq Cocone(D, C) \simeq Cocone(D \circ H, C) \simeq \Hom_\mcC(\varinjlim (D \circ H), C)\]
    We have thus proven that $\varinjlim D$ has the universal property of the colimit of $D \circ H$, and that $\varinjlim(D \circ H)$ has the universal property of the colimit of $D$.  Because colimits are unique up to isomorphism,
    \[ \varinjlim D \simeq \varinjlim(D \circ H) \qedhere \]
\end{proof}

\begin{corollary}\label{cor:thincofinal}
    For thin (posetal) categories,
    if $P_0 \subseteq P$ is cofinal in the sense of order theory, that is, $\forall x \in P.\ \exists y \in P_0$ such that $x \le y$, then 
    \[ \colim_{p \in P} Fp \simeq \colim_{p \in P_0} Fp \]
\end{corollary}
\begin{proof}
    Condition (a) of Definition \ref{def:cofinal} follows by the cofinality hypothesis on $P_0$.  Condition (b) follows trivially from the fact that $P$ is a thin category.  By application of Theorem \ref{theorem:cofinal}, we have the desired result.
\end{proof}

\begin{corollary}\label{cor:omegacofinal}
    For an $\omega$-indexed diagram $F : \omega \to \mcC$, and any strictly increasing function $f : \omega \to \omega$, we have
    \[ \colim F \simeq \colim (F \circ f)\]
\end{corollary}
\begin{proof}
    Observe that the image of a strictly increasing function is cofinal in the sense of order theory.  The desired result follows immediately from Corollary \ref{cor:thincofinal}.
\end{proof}

We have thus proven a sufficient condition for a subdiagram to share a colimit with its parent diagram.  Corollaries \ref{cor:thincofinal} and \ref{cor:omegacofinal} will be especially useful later in this document, where we interpret certain recursive data types as colimits of $\omega$-indexed diagrams.

\subsection{Iterated colimits}\label{sec:iteratedcolimits}

In this section, we relate iterated colimits to colimits indexed by product categories. Specifically, for a bifunctor $F : \mcC \times \mcC \to \mcD$, we show that taking colimits in each variable successively agrees, up to isomorphism, with taking a single colimit over the product category:
\[ \colim_{(X, Y) \in \mcC \times \mcC}(F(X, Y)) \simeq \colim_{X \in \mcC}(\colim_{Y \in \mcC}(F(X,Y))) \]

This result will be used later in this chapter, where we show that consecutive applications of the least fixed point operator (which is later used to interpret recursive data types) can always be reduced to a single application.

First, we recall the well-known fact that colimits in product categories are computed coordinatewise.

\begin{theorem}\label{thm:productcolmits}
    Let $G : \mcC \to \mcD \times \mcD$ be a functor. Write $GX = (G_1 X, G_2 X)$.  Then
    \[\varinjlim G \simeq (\varinjlim G_1, \varinjlim G_2)\]
\end{theorem}

\begin{proof}
Let $\{ \gamma_X^1 : G_1(X) \to \colim G_1\ |\ X \in \mcC\}$ be the morphism components of the colimit $\varinjlim G_1$ and $\{ \gamma_X^2: G_2(X) \to \colim G_2\ |\ X \in \mcC\}$ be the morphism components of $\varinjlim G_2$.

We verify that $(\varinjlim G_1, \varinjlim G_2)$ has the universal property of the initial cocone for $G$.  First, we establish that $(\varinjlim G_1, \varinjlim G_2)$ forms a cocone for $G$.  Next, we prove that there exists a cocone morphism $h$ from $(\varinjlim G_1, \varinjlim G_2)$ to any other cocone $((C_1,C_2), (\lambda^1_X, \lambda_X^2) : GX \to (C_1, C_2))_{X \in \mcC}$.  Finally, we show that any other cocone morphism $h' : (\varinjlim G_1, \varinjlim G_2) \to (C_1, C_2)$ coincides with $h$.

First, note that $(\varinjlim G_1, \varinjlim G_2)$ forms a cocone for $G$ by the following reasoning.  Let $f : X \to Y$ be a morphism in $\mcC$. Then
\[ (\gamma^1_Y, \gamma^2_Y) \circ G (f) = (\gamma^1_Y, \gamma^2_Y) \circ (G_1(f), G_2(f)) = (\gamma^1_Y \circ G_1(f), \gamma^2_Y \circ G_2(f))  = (\gamma^1_X, \gamma^1_X)\]
where the last equality holds because $\varinjlim G_1$ and $\varinjlim G_2$ are cocones. 

Thus, $(\varinjlim G_1, \varinjlim G_2)$ forms a cocone for $G$.

Now we fix an arbitrary cocone $((C_1, C_2), (\lambda^1_X, \lambda_X^2) : G X \to (C_1, C_2))_{X \in \mcC}$.
We will show that there exists a cocone morphism from $(\varinjlim G_1, \varinjlim G_2)$ to $((C_1, C_2), (\lambda^1_X, \lambda_X^2))_{X \in \mcC}$.

By the universal property of $\varinjlim G_1$, there exist unique morphisms $h_1 : \varinjlim G_1 \to C_1$ and $h_2 : \varinjlim G_2 \to C_2$ such that
\[ h_1 \circ \gamma^1_X = \lambda^1_X  \qquad h_2 \circ \gamma^2_X = \lambda^2_X \]
It follows that $(h_1, h_2)$ is a cocone morphism from $(\varinjlim G_1, \varinjlim G_2)$ to $((C_1, C_2), (\lambda^1_X, \lambda^2_X))_{X \in \mcC}$, since
\[ (h_1, h_2) \circ (\gamma^1_X, \gamma^2_X) = (h_1 \circ \gamma^1_X, h_2 \circ \gamma^1_X) = (\lambda^1_X, \lambda^2_X) \]

Lastly, we show uniqueness.  Let $(h_1', h_2')$ be a cocone morphism from $(\varinjlim G_1, \varinjlim G_2)$ to $((C_1, C_2), (\lambda^1_X, \lambda^2_X))_{X \in \mcC}$.  We will show $(h_1, h_2) = (h_1', h_2')$.  It is sufficient to show that $h_1 = h_1'$ and that $h_2 = h_2'$.  We know that
\[ (h_1', h_2') \circ (\gamma^1_X, \gamma^2_X) = (\lambda^1_X, \lambda^2_X) \]
because $(h_1', h_2')$ is a cocone morphism.  From this, we can conclude that
\[ h_1' \circ \gamma^1_X = \lambda^1_X \qquad h_2' \circ \gamma^2_X = \lambda^2_X \]
But $h_1$ and $h_2$ are the unique morphisms with these respective properties, so $h_1 = h_1'$ and $h_2 = h_2'$.
\end{proof}

Next, we define the following functors on arbitrary categories $\mcC$ and $\mcD$:
{\small
\begin{multicols}{2}
\[
\begin{aligned}
    \$ : \mcD^{\mcC \times \mcC} &\to (\mcD^{\mcC})^\mcC \\
    \$ F(X)(Y) &= F(X, Y) \\
    \$(\eta : F \naturalto G) &= (\eta' : \$F \naturalto \$G) \\
    \text{where } \eta'_X &: \$FX \naturalto \$GX\\
    (\eta'_X)_Y &= \eta_{(X, Y)}
\end{aligned}
\]

\columnbreak

\[
\begin{aligned}
    \lambda : (\mcD^\mcC)^\mcC &\to \mcD^{\mcC \times \mcC} \\
    \lambda F(X, Y) &= F(X)(Y) \\
    \lambda (\eta : F \naturalto G) &= (\eta' : \lambda F \naturalto \lambda G) \\
    \text{where }\eta'_{(X,Y)} &= (\eta_X)_Y
\end{aligned}
\]
\end{multicols}
}

\begin{lemma}
    The functors $\$$ and $\lambda$ defined above form an isomorphism of categories between the categories $\mcD^{\mcC \times \mcC}$ and $(\mcD^{\mcC})^\mcC$.
\end{lemma}

\begin{proof}
    By inspecting the definitions of $\$$ and $\lambda$, it is immediate that both their object parts and functorial actions are (strict) inverses to one another.
\end{proof}

\begin{corollary}\label{cor:curry_cat_iso}
    Given two functors $F, G : \mcC \times \mcC \to \mcD$, 
    \begin{equation}\label{eq:dollarisff}
        \Hom_{\mcD^{\mcC \times \mcC}}(F, G) \simeq \Hom_{(\mcD^{\mcC})^\mcC}(\$F, \$G)
    \end{equation}
\end{corollary}

\begin{proof}
    Since $\$$ is an isomorphism of categories, it is full and faithful.  Therefore, for any functors $F, G : \mcC \times \mcC \to \mcD$, there is a bijection of the form given in equation \eqref{eq:dollarisff}.  This bijection is natural in both $F$ and $G$ because $\$$ and $\lambda$ form a category isomorphism.
\end{proof}

\begin{lemma}\label{lemma:cocone_as_nat_trans}
    Let $\mcA$ and $\mcB$ be categories.  Let $K_X : \mcA \to \mcB$ be the constantly $X$-valued functor.  For all functors $F : \mcA \to \mcB$ and all objects $X \in \mcB$, there exists the following isomorphism
    \[ \Hom_{\mcB}(\colim F, X) \simeq \Hom_{\mcB^\mcA}(F, K_X)\]
    which is natural in $X$.
\end{lemma}

\begin{proof}
    Every cocone of $F$ is an object $X$ together with a family of morphisms $(a_Y : F Y \to X)_{Y \in \mcA}$ such that for all morphisms $g : W \to Z$ in $\mcA$, the diagram shown in Figure \ref{fig:cocone} commutes.
    \begin{figure}[htbp]
    \centering
        \begin{tikzcd}
	FW && FZ \\
	\\
	X && X
	\arrow["{F(g)}", from=1-1, to=1-3]
	\arrow["{a_W}", from=1-1, to=3-1]
	\arrow["{a_Z}", from=1-3, to=3-3]
	\arrow[no head, from=3-1, to=3-3]
	\arrow[shift left, no head, from=3-1, to=3-3]
        \end{tikzcd}
        \caption{The definition of a cocone.}
        \label{fig:cocone}
    \end{figure}

    Now observe that a natural transformation $\eta : F \to K_X$ consists of 
    \begin{itemize}[nosep]
        \item a morphism $\eta_Y : F Y \to X$ for every object $Y \in \mcA$
        \item such that for every morphism $g : W \to Z$ in $A$,
        \[ K_X(g) \circ \eta_W = \eta_Z \circ F(g) \]
    \end{itemize}
    Since $K_X(g) = id_X$ for all morphisms $g$, this condition becomes
    \[ \eta_W = \eta_Z \circ F(g) \]
    which is precisely the cocone property.

    Therefore, giving a cocone for $F$ is the same as giving a natural transformation from $F$ to $K_X$, and
    \[ Cocone(F, X) \simeq Hom_{\mcB^\mcA}(F, K_X) \]

    Now recall that the colimit of a functor $F : \mcA \to \mcB$ is defined as a cocone $(\colim F, c_Y : F Y \to \colim F)_{Y \in \mcA}$ such that giving a morphism out of the object component $\colim F$ of the colimit of $F$ is equivalent to giving a cocone of $F$. That is, the following isomorphism exists and is natural in $X$:
    \[ \Hom_B(\colim F, X) \simeq Cocone(F, X) \]
    Note that the above isomorphism is trivially natural in $X$, since both sides transfer under arbitrary $f : X \to Y$ by post-composition.

    Composing these two natural isomorphisms yields the desired result.
\end{proof}

For the main theorem of this section, we will also need the following well-known fact which is a consequence of the Yoneda lemma.

\begin{prop}\label{prop:yoneda_iso}
    Given a category $\mcD$, and objects $A$ and $B$ in $\mcD$, if there exists a family of isomorphisms
    \[ \Hom_\mcD(A, X) \simeq \Hom_\mcD(B, X) \]
    natural in $X \in \mcD$, then
    \[ A \simeq B\]
    in $\mcD$.
\end{prop}

\begin{proof}
    By the Yoneda lemma, the Yoneda embedding functor
    \[ y : \mcD \to \Set^\mcD,\quad X \mapsto Hom(X, -)\]
    is full and faithful, and therefore reflects isomorphisms.
    The hypothesis gives an isomorphism between $y_A$ and $y_B$, hence $A \simeq B$.
\end{proof}

We now give the main theorem of Section \ref{sec:iteratedcolimits}.

\begin{theorem}\label{thm:iterated_colimits}
    Let $F : \mcC \times \mcC \to \mcD$ be a functor. If $F$ has a colimit, then
    \[ \colim_{(X, Y) \in \mcC \times \mcC}(F(X, Y)) \simeq \colim_{X \in \mcC}(\colim_{Y \in \mcC}(F(X,Y)) \]
\end{theorem}
\begin{proof}
    Let $K_Z$ be the constantly $Z$-valued functor.  We apply Corollary \ref{cor:curry_cat_iso} together with Lemma \ref{lemma:cocone_as_nat_trans} to form the following chain of isomorphisms which are natural in $Z$.

    \begin{align*}
        &\phantom{\simeq}\;\; \Hom_{\mcD}\left(\colim_{(X,Y) \in \mcC \times \mcC}\Big(F(X,Y)
        \Big), Z\right) \\
        &\simeq \Hom_{\mcD^{\mcC \times \mcC}}(F, K_Z)  && \textsf{(by Lemma \ref{lemma:cocone_as_nat_trans})}\\
        &\simeq \Hom_{(\mcD^{\mcC})^\mcC}(\$F, \$K_Z) = 
        \Hom_{(\mcD^{\mcC})^\mcC}(\$F, K_{K_Z})   && \textsf{(by Corollary \ref{cor:curry_cat_iso})}\\
        &\simeq \Hom_{\mcD^\mcC}(\colim_{X \in \mcC}(\lambda Y. F(X, Y)), K_Z)   && \textsf{(by Lemma \ref{lemma:cocone_as_nat_trans})}\\
        &\simeq \Hom_\mcD(\colim_{Y \in \mcC}(\colim_{X \in \mcC}(F(X, Y))), Z)   && \textsf{(by Lemma \ref{lemma:cocone_as_nat_trans})}
    \end{align*}

    By Proposition \ref{prop:yoneda_iso}, we have
    \[ \colim_{(X, Y) \in \mcC \times \mcC}(F(X, Y)) \simeq \colim_{X \in \mcC}(\colim_{Y \in \mcC}(F(X,Y)) \]
    which is exactly what we wanted to show.
\end{proof}

This theorem shows that colimits computed iteratively can instead be computed in a product category. This observation will be used in the next section to simplify certain constructions involving nested recursion, allowing them to be expressed in terms of a single recursive step.

\section{Construction of initial algebras}\label{sec:constructinginitalg}

In this section, we lay the groundwork for the interpretation of a class of recursive data types called algebraic data types.  We begin by introducing $F$-algebras and initial $F$-algebras, highlighting how they generalize monotone functions on partially ordered sets and how the least fixed point of a monotone function corresponds directly to the initial $F$-algebra.  Later, we introduce the \emph{least fixed point} functor, the central construction used to interpret recursively defined types, and establish some of its key properties: it is a functor and it gives the carrier of the initial algebra of a functor to which it is applied.  Finally, we show how the initial algebra of a functor can be constructed using the least fixed point functor.

\subsection{Algebras and initial algebras}\label{sec:init_alg}
We start by defining $F$-algebras and initial $F$-algebras, which form the categorical foundation for interpreting recursive types.

Throughout this section, let $\mcC$ be an $\omega$-cocomplete category.

\begin{definition}
    Given an endofunctor $F : \mcC \to \mcC$, an \emph{$F$-algebra} $(X, f)$ is an object $X$ of $\mcC$ together with a morphism $f : F(X) \to X$ in $\mcC$.   
\end{definition}

\begin{figure}[htbp]
        \centering
        \begin{tikzcd}[scale=2.0,sep=large]
            F(X) \arrow[r, "f"]
            & X
        \end{tikzcd}
        \caption{A diagram showing an $F$-algebra}
        \label{fig:falg}
\end{figure}

$F$-algebras form their own category with $F$-algebras as objects and certain morphisms in $\mcC$ as morphisms.  Given two $F$-algebras $(Y : \mcC, g : F(Y) \to Y)$ and $(X : \mcC, f : F(X) \to X)$, a morphism in the category of $F$-algebras with domain $(Y, g)$ and codomain $(X, f)$ is a morphism $h$ in $\mcC$ such that $h \circ g = f \circ F (h) $. That is, it is a morphism $h$ such that the diagram shown in Figure \ref{fig:falgmorphism} commutes.

\begin{figure}[htbp]
    \centering
    \begin{tikzcd}[scale=2.0,sep=large]
        F(Y) \arrow[r, "g"] \arrow[d, "F(h)"] & Y \arrow[d, "h"] \\
        F(X) \arrow[r, "f"] & X
    \end{tikzcd}
    \caption{A diagram showing an $F$-algebra morphism}
    \label{fig:falgmorphism}
\end{figure}


In such a category, an initial object is an $F$-algebra $(\mu F, in_F)$ such that for any other $F$-algebra $(X, f)$, there exists a unique morphism $h : \mu F \to X$ such that the diagram shown in Figure \ref{fig:initialfalg} commutes.

\begin{figure}[htbp]
    \centering
    \begin{tikzcd}[scale=2.0,sep=large]
        F(\mu F) \arrow[r, "in_F"] \arrow[d, "F(h)"] & \mu F \arrow[d, "h"] \\
        F(X) \arrow[r, "f"] & X
    \end{tikzcd}
    \caption{A diagram showing an initial $F$-algebra}
    \label{fig:initialfalg}
\end{figure}

An initial algebra is a direct generalization of the concept of a least fixed point of a monotone function on a partially ordered set.  Given a partially ordered set $(P, \sqsubseteq)$ and a monotone function $f : P \to P$, a least fixed point of $f$ is an element $\mu f$ of $P$ that is a fixed point of $f$, i.e., $f(\mu f) = \mu f$, such that for any other fixed point $x$ of $f$, $\mu f \sqsubseteq x$.  This notion can be generalized to categories.  

Let $F : \mcC \to \mcC$ be an endofunctor on the category $\mcC$.  A fixed point of $F$ is an object $X$ of $\mcC$ such that $X$ is isomorphic to $F(X)$.  Elements of $P$ are generalized by objects of $\mcC$.  A relation between two elements of $P$, for example $x \sqsubseteq y$, is generalized as a morphism in $\mcC$ between objects $X$ and $Y$.  For $x \in P$ to be \textit{less than} $y \in P$ means that $x \sqsubseteq y$.  

An $F$-algebra generalizes the notion of a prefix point of a monotone function: given any $F$-algebra $(X, h)$, we have a morphism $h : F(X) \to X$.  For an $F$-algebra $(X, h)$ to be \textit{less than} another $F$-algebra $(Y, g)$ means that a morphism exists in the category of $F$-algebras with domain $(X, h)$ and codomain $(Y , g)$.  This is analogous to having $x$ and $y$ as prefix points of a monotone function $f$, and $f\ x \sqsubseteq f\ y$ as well as $x \sqsubseteq y$.  

Every least prefix point of a monotone function is also a fixed point (and is the least fixed point of such a function).  A \textit{least} fixed point of a functor $F$ is therefore an $F$-algebra $(\mu F : \mcC, in_F : F(\mu F) \to \mu F)$ such that given any other $F$-algebra $(X, f)$, there exists a (unique) $F$-algebra morphism with domain $(\mu F, in_F)$ and codomain $(X, f)$.\footnote{Lambek's Lemma guarantees that $in_F$ is an isomorphism, so $\mu F$ is indeed a fixed point of $F$; however, we will not need this fact. We will construct initial algebras directly, and the invertibility of $in_F$ will follow by construction.} Having a unique morphism from $(\mu F, in_F)$ to every other $F$-algebra characterizes $(\mu F, in_F)$ as the initial object in the category of $F$-algebras.  The initial algebra of a functor $F$ is therefore the least fixed point of $F$.

\subsection{The least fixed point functor}\label{sec:lfp}

We can make this correspondence between initial algebras and least fixed points systematic by defining a functor
\[ LFP : [\mcC,\mcC] \to \mcC \]
which maps an $\omega$-cocontinuous endofunctor $F : \mcC \to \mcC$ to the least fixed point of $F$—that is, to the object component of $F$’s initial algebra.  We now proceed to define the action of $LFP$ on objects the functor category $[\mcC,\mcC]$.

Given a functor $F : \mcC \to \mcC$, we define the functor $F_{seq} : \omega \to \mcC$ as follows: the action of $F_{seq}$ on objects of $\omega$ is given as
\[ F_{seq}(n) = F^n(\emptyset) \]
where $\emptyset$ is initial in $\mcC$ and $F^n$ is $n$-fold repeated composition of $F$.  The action of $F_{seq}$ on morphisms is given as
\begin{align*}
    F_{seq}(0 \le 1) &= \emptyset_{F(\emptyset)} \\
    F_{seq}(n+1 \le n+2) &= F(F(n \le n+1))
\end{align*}
where $\emptyset_{X}$ is the unique morphism from the initial object $\emptyset$ to $X$.
This is extended by transitivity to the entire relation $\le$.

The action of $LFP$ on objects can now be stated as follows: 
$LFP(F)$ is the object component of the cocone given by $\colim F_{seq}$.
That is, $LFP(F)$ is the object component of the colimit of the diagram shown in Figure \ref{fig:lfpseq}.
\begin{figure}[htbp]
    \centering
{    \begin{tikzcd}[scale cd=1,sep=large]
        \emptyset \arrow[r, "\emptyset_{F(\emptyset)}"] & F(\emptyset) \arrow[r, "F(\emptyset_{F(\emptyset)})"] & F^2(\emptyset) \arrow[r]  & \cdots \arrow[r] 
        & F^n(\emptyset) \arrow[r, "F^n(\emptyset_{F(\emptyset)})"] & F^{n+1}(\emptyset) \arrow[r] & \cdots
    \end{tikzcd}
    \caption{The diagram whose colimit defines the action of $LFP$ on objects}
    \label{fig:lfpseq}
}
\end{figure}

$LFP$'s action on morphisms in the category of $\omega$-cocontinuous endofunctors on $\mcC$ (i.e., natural transformations) is given as follows: given a natural transformation $\eta : F \naturalto G$, we can define an infinite set of natural transformations $\eta^i : F^i \naturalto G^i$ indexed by objects of $\omega$, that is, by natural numbers.  We define $\eta^i$ by induction on $i$.
\begin{align*}
    \eta^0_X &= id_X \\
    \eta^{i+1}_X &= G(\eta^i_X) \circ \eta_{F^i(X)}
\end{align*}
Note that by considering the naturality square for $\eta$ shown in Figure \ref{fig:eta_square}, $\eta^{i+1}_X$ can equally be defined as $\eta_{G^i(X)} \circ F(\eta^i_X)$.

\begin{figure}[htbp]
    \centering
    \begin{tikzcd}[scale=2.0,sep=large]
        F^{i+1}(X) \arrow[r, "\eta_{F^i(X)}"] \arrow[d, "F(\eta^i_X)"] \arrow[rd, "\eta^{i+1}_X", dashed]
        & G (F^i(X)) \arrow[d, "G(\eta^i_X)"] \\
        F(G^i(X)) \arrow[r, "\eta_{G^i(X)}"]
        & G^{i+1}(X)
    \end{tikzcd}
    \caption{The naturality square for $\eta$ used in defining $\eta^{i+1}$}
    \label{fig:eta_square}
\end{figure}

In the remainder of this subsection, we fix two $\omega$-cocontinuous functors $F : \mcC \to \mcC$ and $G : \mcC \to \mcC$.  Now that we have defined $\eta^i$ for an arbitrary $\eta$, consider Figure \ref{fig:lfpmorphism} showing a diagram in $\mcC$ containing the diagram $F_{seq}$, the diagram $G_{seq}$, and each of their related colimits.  In this diagram, and in the remainder of this chapter, we denote the object part of the colimits of $F_{seq}$ and of $G_{seq}$ as $LFP(F)$ and $LFP(G)$, respectively.  We denote the morphism components of the colimit of $F_{seq}$ as $\phi_i$, where $\phi_i : F^i(\emptyset) \to LFP(F)$.  We denote the morphism components of the colimit of $G_{seq}$ as $\psi^i$, where $\psi^i : G^i(\emptyset) \to LFP(G)$.  Additionally, to avoid cumbersome notation, we denote the unique morphism with domain $F^i(\emptyset)$ and codomain $F^j(\emptyset)$ (with $i \le j$) in the image of $F_{seq}$ as $f_{i \to j}$.  We denote the unique morphism with domain $G^i(\emptyset)$ and codomain $G^j(\emptyset)$ in the image of $G_{seq}$ as $g_{i \to j}$.

\begin{figure}[htbp]
    \centering
{    \begin{tikzcd}[scale cd=0.67,sep=large]
        \emptyset \arrow[r, "f_{0 \to 1}"] \arrow[drrrrrrr, "\phi_0", bend right=9] \arrow[dd, "\eta^0_\emptyset"] &
            F(\emptyset) \arrow[r, "f_{1 \to 2}"] \arrow[drrrrrr, "\phi_1", bend right=5] \arrow[dd, "\eta^1_\emptyset"] & 
            F^2(\emptyset) \arrow[r] \arrow[drrrrr, "\phi_2"] \arrow[dd, "\eta^2_\emptyset"]  & 
            \cdots \arrow[r] &
            F^n(\emptyset) \arrow[r, "f_{n \to n+1}"] \arrow[drrr, "\phi_n", pos=0.6] \arrow[dd, "\eta^n_\emptyset", pos=0.75] & 
            F^{n+1}(\emptyset) \arrow[r] \arrow[drr, "\phi_{n+1}"] \arrow[dd, "\eta^{n+1}_\emptyset", pos=0.75] & \cdots \\
        & & & & & & & LFP(F) \arrow[dd, "LFP(\eta)", dashed] \\
        \emptyset \arrow[r, "g_{0 \to 1}"] \arrow[drrrrrrr, "\psi_0", bend right=9] &
            G(\emptyset) \arrow[r, "g_{1 \to 2}"] \arrow[drrrrrr, "\psi_1", bend right=5] & 
            G^2(\emptyset) \arrow[r] \arrow[drrrrr, "\psi_2"] & 
            \cdots \arrow[r] &
            G^n(\emptyset) \arrow[r, "g_{n \to n+1}"] \arrow[drrr, "\psi_n"] & 
            G^{n+1}(\emptyset) \arrow[r] \arrow[drr, "\psi_{n+1}"] & \cdots \\
        & & & & & & & LFP(G) \\
    \end{tikzcd}
    \caption{A morphism of cocones of $F_{seq}$ and $G_{seq}$}
    \label{fig:lfpmorphism}
}
\end{figure}

To define $LFP(\eta : F \naturalto G) : LFP(F) \to LFP(G)$, we first must generate a morphism from $LFP(F)$ to $LFP(G)$.  We do this by showing that $LFP(G)$ together with each $\eta^i_\emptyset$ and $\psi^i$, form a cocone for $F_{seq}$, and therefore there exists a unique cocone morphism from the initial cocone $LFP(F)$ to $LFP(G)$.

To prove that $LFP(G)$ is a cocone for $F_{seq}$, we make use of the following lemma:

\begin{lemma}\label{lemma:lfp_morphism_square}
    Given a natural transformation $\eta : F \naturalto G$, and any $n \in N$, we have that $\eta^{n+1}_\emptyset \circ f_{n\to n+1} = g_{n\to n+1} \circ \eta^n_\emptyset$.  That is, each square in the diagram shown in Figure \ref{fig:lfpmorphism} commutes.
\end{lemma}

\begin{proof} We proceed by induction on $n$.

    If $n = 0$, we have that $\eta^1_\emptyset \circ f_{0\to 1} = g_{0 \to 1} \circ \eta^0_\emptyset$.  Reducing these morphisms to their definitions, we get
    \[ \eta_\emptyset \circ \emptyset_{F(\emptyset)} = \emptyset_{G(\emptyset)} \circ id_\emptyset \]
    Each of these morphisms has an initial object $\emptyset$ as its domain.  Because there is a unique morphism from any initial object to any object, these morphisms must be equal.

    Given
    \[ \eta^{n+1}_\emptyset \circ f_{n\to n+1} = g_{n \to n+1} \circ \eta^n_\emptyset \]
    we must now show
    \[ \eta^{n+2}_\emptyset \circ f_{n+1\to n+2} = g_{n+1 \to n+2} \circ \eta^{n+1}_\emptyset \]
    We prove this as follows:
    \begin{align*}
        \eta^{n+2}_\emptyset \circ f_{n+1\to n+2} &= \eta^{n+2} \circ F(f_{n\to n+1})  && (\textsf{By definition of $f_{n+1 \to n+2}$}) \\
        &= \eta_{G^{n+1}(\emptyset)} \circ F(\eta^{n+1}_\emptyset) \circ F(f_{n \to n + 1}) && (\textsf{By definition of $\eta^{n+2}$}) \\
        &= \eta_{G^{n+1}(\emptyset)} \circ F(\eta^{n+1}_\emptyset \circ f_{n \to n + 1}) && (\textsf{By functoriality of $F$}) \\
        &= \eta_{G^{n+1}(\emptyset)} \circ F(g_{n \to n+1} \circ \eta^n_\emptyset) && (\textsf{By naturality of $\eta$}) \\
        &= G(g_{n \to n+1} \circ \eta^n_\emptyset) \circ \eta_{F^{n}(\emptyset)} && (\textsf{By naturality of $\eta$}) \\
        &= G(g_{n \to n+1}) \circ G(\eta^n_\emptyset) \circ \eta_{F^n(\emptyset)} && (\textsf{By functoriality of $G$}) \\
        &= G(g_{n \to n+1}) \circ \eta^{n+1}_\emptyset && (\textsf{By definition of $\eta^{n+1}$)} \\
        &= g_{n+1 \to n+2} \circ \eta^{n+1}_\emptyset && (\textsf{By definition of $g_{n+1 \to n+2}$})
    \end{align*}

    Thus, each square in the diagram shown in Figure \ref{fig:lfpmorphism} commutes.
\end{proof}

Now we can show that $LFP(G)$ forms a cocone for $F_{seq}$.

\begin{lemma}
    Given a natural transformation $\eta : F \naturalto G$, $LFP(G)$ together with the set of morphisms $\{\psi_i \circ \eta^i_\emptyset\ |\ i \in \omega\}$ is a cocone of  $F_{seq}$.
\end{lemma}

\begin{proof}
    We must show the following:
    \[ \forall n \in \omega.\ \psi_{n+1} \circ \eta^{n+1}_\emptyset \circ f_{n \to n+1} = \psi_n \circ \eta^n_\emptyset \]
    We do this by the following sequence of equations:
    \begin{align*}
        \psi_n \circ \eta^n_\emptyset &= \psi_{n+1} \circ g_{n \to n+1} \circ \eta^n_\emptyset && (\textsf{Because $(LFP(G), \psi_i)_{i \in \omega}$ forms a cocone for $G_{seq}$})\\
        &= \psi_{n+1} \circ \eta^{n+1}_\emptyset \circ f_{n \to n+1} && (\textsf{By Lemma \ref{lemma:lfp_morphism_square}})
    \end{align*}

    All morphisms in the set $\{\psi_i \circ \eta^i_\emptyset\ |\ i \in \omega\}$ commute with the morphisms in the image of $F_{seq}$, so $(LFP(G), \psi_i \circ \eta^i_\emptyset)_{i \in \omega}$ forms a cocone for $F_{seq}$.
\end{proof}

$(LFP(F), \phi_i)_{i \in \omega}$ is the initial cocone for $F_{seq}$.  Because $(LFP(G), \psi_i \circ \eta^i_\emptyset)_{i \in \omega}$ forms a cocone for $F_{seq}$, there is a unique cocone morphism $m : (LFP(F), \phi_i)_{i \in \omega} \to (LFP(G), \psi_i~\circ~\eta^i_\emptyset)_{i \in \omega}$.  That is, there is a unique morphism $m$ in $\mcC$ satisfying
\[ \forall n \in \omega.\ m \circ \phi_n = \psi_n \circ \eta^n_\emptyset \]
We define the action of $LFP$ on a morphism $\eta$ in $C$ as the unique morphism $m$ satisfying the above requirement.

We must now prove that the action of $LFP$ on morphisms is functorial.  For this proof, we refer to the diagram shown in Figure \ref{fig:lfpcomposition}.

\begin{figure}[htbp]
    \centering
{    \begin{tikzcd}[scale cd=0.8]
        \emptyset \arrow[r, "f_{0\to 1}"] \arrow[drrrrrrr, "\phi_0", bend right=9] \arrow[dd, "\eta^0_\emptyset"] &
            F(\emptyset) \arrow[r, "f_{1\to 2}"] \arrow[drrrrrr, "\phi_1", bend right=5, pos=0.55] \arrow[dd, "\eta^1_\emptyset"] & 
            F^2(\emptyset) \arrow[r] \arrow[drrrrr, "\phi_2"] \arrow[dd, "\eta^2_\emptyset"]  & 
            \cdots \arrow[r] &
            F^n(\emptyset) \arrow[r, "f_{n\to n+1}"] \arrow[drrr, "\phi_n", pos=0.5] \arrow[dd, "\eta^n_\emptyset", pos=0.75] & 
            F^{n+1}(\emptyset) \arrow[r] \arrow[drr, "\phi_{n+1}"] \arrow[dd, "\eta^{n+1}_\emptyset", pos=0.75] & \cdots \\
        & & & & & & & LFP(F) \arrow[dd, "LFP(\eta)", dashed, pos=0.75] \arrow[dddd, "LFP(\zeta \circ \eta)", dashed, bend left=40] \\
        \emptyset \arrow[r, "g_{0 \to 1}"] \arrow[drrrrrrr, "\psi_0", bend right=9] \arrow[dd, "\zeta^0_\emptyset"] &
            G(\emptyset) \arrow[r, "g_{1 \to 2}"] \arrow[drrrrrr, "\psi_1", bend right=5, pos=0.55] \arrow[dd, "\zeta^1_\emptyset"] & 
            G^2(\emptyset) \arrow[r] \arrow[drrrrr, "\psi_2"] \arrow[dd, "\zeta^2_\emptyset"] & 
            \cdots \arrow[r] &
            G^n(\emptyset) \arrow[r, "g_{n \to n+1}"] \arrow[drrr, "\psi_n", pos=0.5] \arrow[dd, "\zeta^n_\emptyset", pos=0.75] & 
            G^{n+1}(\emptyset) \arrow[r] \arrow[drr, "\psi_{n+1}"] \arrow[dd, "\zeta^{n+1}_\emptyset", pos=0.75] & \cdots \\
        & & & & & & & LFP(G) \arrow[dd, "LFP(\zeta)", dashed, pos=0.25] \\
        \emptyset \arrow[r, "h_{0 \to 1}"] \arrow[drrrrrrr, "\chi_0", bend right=9] &
            H(\emptyset) \arrow[r, "h_{1 \to 2}"] \arrow[drrrrrr, "\chi_1", bend right=5] & 
            H^2(\emptyset) \arrow[r] \arrow[drrrrr, "\chi_2"] & 
            \cdots \arrow[r] &
            H^n(\emptyset) \arrow[r, "h_{n \to n+1}"] \arrow[drrr, "\chi_n"] & 
            H^{n+1}(\emptyset) \arrow[r] \arrow[drr, "\chi_{n+1}"] & \cdots \\
        & & & & & & & LFP(H) \\
    \end{tikzcd}
    \caption{The action of $LFP$ on a composition of natural transformations}
    \label{fig:lfpcomposition}
}
\end{figure}

First, we show that $LFP$'s action on morphisms preserves composition.

\Needspace{4\baselineskip}
\begin{lemma}\label{lemma:lfp_composition}
    Let $F, G, H : \mcC \to \mcC$ be $\omega$-cocontinuous functors on $\mcC$.
    
    Let $\eta : F \naturalto G$ and $\zeta : G \naturalto H$ be natural transformations. Then
    \[ LFP(\zeta \circ \eta) = LFP(\zeta) \circ LFP(\eta) \]
\end{lemma}

\begin{proof}
    By the universal property of the initial cocone $LFP(F)$, we know that $LFP(\zeta \circ \eta)$ is the unique morphism in $\mcC$ satisfying for all $k$, 
    \[ \chi^k \circ \zeta^k_\emptyset \circ \eta^k_\emptyset = LFP(\zeta \circ \eta) \circ \phi_k \]
    We can show that $LFP(\zeta) \circ LFP(\eta)$ satisfies this requirement for all $k$ as follows:
    \begin{align*}
        LFP(\zeta) \circ LFP(\eta) \circ \phi_k
        &=  LFP(\zeta) \circ \psi_k \circ \eta_\emptyset^k && \textsf{(By the universal property of $LFP(F)$)} \\
        &= \chi_k \circ \zeta^k_\emptyset \circ \eta^k_\emptyset && \textsf{(By the universal property of $LFP(G)$)}
    \end{align*}

    $LFP(\zeta) \circ LFP(\eta)$ satisfies the unique property of the cocone morphism $LFP(\zeta \circ \eta)$ with domain $LFP(F)$ and codomain $LFP(H)$, therefore $LFP(\zeta \circ \eta) = LFP(\zeta) \circ LFP(\eta)$ and $LFP$'s action on morphisms preserves composition.
\end{proof}

It remains to show that $LFP$'s action on morphisms preserves identity natural transformations.  To prove this, we make use of the following lemma.

\begin{lemma}
    Given any functor $F : \mcC \to \mcC$, its identity natural transformation $id_F$ has the following property:

    \[ \forall k.\ (id_F)^k = id_{F^k} \]
\end{lemma}

\begin{proof}
    We proceed by induction on $k$.

    Let $k = 0$. Then, by definition, $(id_F)^0 = id_{F^0}$.

    Let $k = n + 1$.  Assume that $(id_F)^n = id_{F^n}$. Then we have
    \begin{align*}
        (id_F)_X^k &= F((id_F)^n_X) \circ (id_F)_{F^n(X)} \\
        &= F((id_{F^n})_X) \circ (id_F)_{F^n(X)} \\
        &= F(id_{F^{n}(X)}) \circ (id_F)_{F^n(X)} \\
        &= id_{F^{n+1}(X)} \circ id_{F^{n+1}(X)} \\
        &= id_{F^{n+1}(X)} \\
        &= id_{F^{k}(X)}
    \end{align*}

    $(id_F)^k_X = (id_{F^k})_X$ for each component morphism of the natural transformation $(id_F)^k$, therefore $(id_F)^k = (id_{F^k})$.
\end{proof}

We may now prove that $LFP$ preserves identity morphisms.

\begin{lemma}\label{lemma:lfp_identities}
    For all identity natural transformations $id_F : F \naturalto F$, $LFP(id_F) = id_{LFP(F)}$.
\end{lemma}

\begin{proof}
    We prove this by showing that $id_{LFP(F)}$ satisfies the unique property of the cocone morphism $LFP(id_F)$ with domain $LFP(F)$ and codomain $LFP(F)$.

    For all $k$, we have that

    \[ id_{LFP(F)} \circ \phi_{k} = \phi_k = (id_{F^k})_\emptyset \circ \phi_k = (id_F)^k_\emptyset \circ \phi_k\]

    Thus, the action on $LFP$ preserves identity morphisms in $[\mcC, \mcC]$.
\end{proof}

We now have everything we need to prove that $LFP$ is a functor.

\Needspace{3\baselineskip}
\begin{theorem}\label{theorem:lfpfunctor}
    $LFP$ is a functor.
\end{theorem}

\begin{proof}
    Assemble Lemmas \ref{lemma:lfp_composition} and \ref{lemma:lfp_identities}.
\end{proof}

The preceding theorem establishes that $LFP$ preserves identities and composition on morphisms and is therefore a functor. In the next subsection, we use this fact to construct initial algebras.  While the least fixed point functor provides the object component of an initial algebra, we still need to define its morphism component and verify that the resulting pair indeed forms an initial algebra.

\subsection{Initial algebras from least fixed points}\label{sec:initial_algebras_morphisms}

In this subsection, we fix an $\omega$-cocontinuous endofunctor $F : \mcC \to \mcC$.  We are interested in defining the morphism component $in_F : F(LFP(F)) \to LFP(F)$ of the initial algebra of $F$.  We define such a morphism as the inverse of the unique cocone morphism from $\colim(F_{Seq})$ to another cocone with vertex $F(LFP(F))$ which we now define.

Let $\phi_i : F^i(\emptyset) \to LFP(F)$ denote the morphism components of the colimit $\colim(F_{Seq})$.  Consider the set of morphisms given by $\{\emptyset_{F(LFP(F))}\} \cup \{ F(\phi^{i}) : F^{i+1}(\emptyset) \to F(LFP(F))\}$.  For these morphisms to form a cocone for $F_{Seq}$ would mean both that
\begin{equation}\label{eq:flfpf0cocone}
    F(\phi_m) \circ f_{0 \to m} = \emptyset_{F(LFP(F))}
\end{equation}
and that
\begin{equation}\label{eq:flfpfcocone}
    F(\phi_m) \circ f_{n+1\to m+1} = F(\phi_n)
\end{equation}
Equation \eqref{eq:flfpf0cocone} holds trivially because each morphism has an initial object as its domain.  Equation \eqref{eq:flfpfcocone} holds by the following reasoning:
\begin{align*}
    F(\phi_m) \circ f_{n+1\to m+1} &= F(\phi_m) \circ F(f_{n \to m})  && \textsf{(By Definition of $f_{n+1 \to m+1}$)} \\
    &= F(\phi_m \circ f_{n \to m})&& \textsf{(By functoriality of $F$)} \\
    &= F(\phi_n) &&\textsf{(The set of $\phi_i$ form the initial cocone for $F_{Seq}$)}
\end{align*}
We therefore obtain a unique cocone morphism from $LFP(F)$ into $F(LFP(F))$ shown in Figure \ref{fig:initalg_in}.  We call this morphism $in^-_F$.  By hypothesis, $F$ is $\omega$-cocontinuous, so $in^-_F$ is an isomorphism, i.e., it is invertible.  We define $in_F$ as the inverse of the morphism $in^-_F$.  In particular, we see that as an invertible cocone morphism, $in_F$ is a cocone isomorphism.  The cocone structure we have imposed on $F(LFP(F))$ therefore makes it initial in the category of cocones of $F_{Seq}$.  We will use this fact in proving that $in_F$ forms the initial $F$-algebra, which we now proceed to do.

\begin{figure}[htbp]
    \centering
{    \begin{tikzcd}[scale cd=0.8,sep=large]
        \emptyset \arrow[r, "f_{0 \to 1}"] \arrow[drrrrrr, "\phi_0", bend right=9] \arrow[ddrrrrrr, "\emptyset_{F(LFP(F))}", bend right=28] &
            F(\emptyset) \arrow[r, "f_{1 \to 2}"] \arrow[drrrrr, "\phi_1", bend right=5] \arrow[ddrrrrr, "F(\phi_0)", bend right=25] & 
            F^2(\emptyset) \arrow[r] \arrow[drrrr, "\phi_2", pos=0.6] \arrow[ddrrrr, "F(\phi_1)", bend right=23, pos=0.6]  & 
            \cdots \arrow[r] &
            F^n(\emptyset) \arrow[r, "f_{n \to n+1}"] \arrow[drr, "\phi_n", pos=0.65]  \arrow[ddrr, "F(\phi_{n-1})", bend right=30, pos=0.6] & 
            F^{n+1}(\emptyset) \arrow[r] \arrow[dr, "\phi_{n+1}"] \arrow[ddr, "F(\phi_{n})", bend right=30, pos=0.75] & \cdots \\
        & & & & & & LFP(F) \arrow[d, "in_F^-", dashed, shift left] \\
        & & & & & & F(LFP(F)) \arrow[u, "in_F", dashed, shift left] \\
    \end{tikzcd}
    \caption{The definition of $in_F$}\label{fig:infcoconemorphism}
    \label{fig:initalg_in}
}
\end{figure}

For the following definition and proof, we hold constant  an arbitrary $F$-algebra $(Y, g : F(Y) \to Y)$.  Given $(Y, g)$, we can define a cocone for $F_{Seq}$ (see Figure \ref{fig:falgcoconemorphism}).  This cocone has object component $Y$ and morphism components $\{\gamma_i : F^i(\emptyset) \to Y\ |\ i \in \omega\}$ where $\gamma_i$ is defined inductively as
    \begin{align*}
        \gamma_0 =& \emptyset_Y \\
        \gamma_{n+1} &= g \circ F(\gamma_n)
    \end{align*}
We consider this cocone (and the morphisms it induces) throughout the following proof.

\begin{lemma}
    $(Y, \gamma_i)_{i \in \omega}$ forms a cocone for $F_{Seq}$.
\end{lemma}

\begin{proof}
    We can show that $(Y, \gamma_i)_{i \in \omega}$ forms a cocone for $F_{Seq}$ by proving for all $m, n \in \nat$ that
    \begin{equation}\label{eq:falg_cocone}
        \gamma_m \circ f_{n \to m} = \gamma_n
    \end{equation}
    By transitivity, it is sufficient to prove that $\gamma_{n+1} \circ f_{n \to n+1} = \gamma_n$ for all $n$.  We prove this by induction as follows:
    \begin{description}
        \item[Base Case:] 
        \[ \gamma_1 \circ f_{0 \to 1} = \gamma_0 \]
        because the domains of each of these morphisms is an initial object.
        \item[Inductive Case:]
        Assume 
        \[ \gamma_n+1 \circ f_{n \to n+1} = \gamma_n \]
        Then
        \begin{align*}
            \gamma_{n+2} \circ f_{n+1 \to n+2} &= g \circ F(\gamma_{n+1}) \circ F(f_{n \to n+1}) && \textsf{(Definitions of $f_{n+1 \to n+2}$ and of $\gamma_{n+2}$)}\\
            &= g \circ F(\gamma_{n+1} \circ f_{n \to n+1}) && \textsf{(By functoriality of $F$)}\\
            &= g \circ F(\gamma_n) && \textsf{(By the inductive hypothesis)}\\
            &= \gamma_{n+1} && \textsf{(By definition of $\gamma_{n+1})$}
        \end{align*}
    \end{description}
    Thus, $(Y, \gamma_i)$ is a cocone for $F_{Seq}$.
\end{proof}

\begin{prop}\label{prop:ininitial}
    $(LFP(F), in_F)$ is the initial $F$-algebra.
\end{prop}

\begin{proof}
    In this proof, we refer to Figure \ref{fig:falgcoconemorphism}.  We will show there is a unique $F$-algebra morphism from $(LFP(F), in_F)$ to the given $F$-algebra $(Y, g)$.

    \begin{figure}[htbp]
    \centering
{    
\begin{tikzcd}
	\cdots & {F^n(\emptyset)} & {F^{n+1}(\emptyset)} & \cdots \\
	& {F(LFP(F))} & {LFP(F)} \\
	& {F(Y)} & Y
	\arrow[from=1-1, to=1-2]
	\arrow["{f_{n \to n+1}}", from=1-2, to=1-3]
	\arrow["F(\phi_{n-1})"', from=1-2, to=2-2]
	\arrow["\phi_{n}", from=1-2, to=2-3, pos=0.2]
	\arrow["\gamma_{n-1}"', bend right=90, from=1-2, to=3-2]
	\arrow[from=1-3, to=1-4]
	\arrow["F(\phi_n)"', from=1-3, to=2-2, pos=0.8]
	\arrow["\phi_{n+1}", from=1-3, to=2-3]
	\arrow["\gamma_n", controls={+(2.5,-0.3) and +(3,-2.5)}, from=1-3, to=3-2]
	\arrow["{in_F}", from=2-2, to=2-3]
	\arrow["{F(m)}", from=2-2, to=3-2]
	\arrow["m", dashed, from=2-3, to=4-4, to=3-3]
	\arrow["g", from=3-2, to=3-3]
\end{tikzcd}
%     \begin{tikzcd}
% 	\emptyset & {F(\emptyset)} & \cdots & {F^n(\emptyset)} & {F^{n+1}(\emptyset)} & \cdots \\
% 	\\
% 	&&& {F(LFP(F)} & {LFP(F)} \\
% 	&&& {F(Y)} & Y
% 	\arrow["{f_{0 \to 1}}", from=1-1, to=1-2]
% 	\arrow["\phi_0", from=1-1, to=3-4]
% 	\arrow[from=1-1, to=3-5]
% 	\arrow[bend right=40, from=1-1, to=4-5]
% 	\arrow[from=1-2, to=1-3]
% 	\arrow[from=1-2, to=3-4]
% 	\arrow[from=1-2, to=3-5]
% 	\arrow[bend right=30, yshift=-3, from=1-2, to=4-4]
% 	\arrow[from=1-3, to=1-4]
% 	\arrow["{f_{n \to n+1}}", from=1-4, to=1-5]
% 	\arrow[from=1-4, to=3-4]
% 	\arrow[from=1-4, to=3-5]
% 	\arrow[bend right=90, from=1-4, to=4-4]
% 	\arrow[from=1-5, to=1-6]
% 	\arrow[from=1-5, to=3-4]
% 	\arrow[from=1-5, to=3-5]
% 	%\arrow[bend left=230, from=1-5, to=4-4]
% 	\arrow["{in_F}", from=3-4, to=3-5]
% 	\arrow["{F(m)}"{description}, from=3-4, to=4-4]
% 	\arrow["m"{description}, dashed, from=3-5, to=4-5]
% 	\arrow["g"{description}, from=4-4, to=4-5]
% \end{tikzcd}
    \caption{An $F$-algebra morphism is a cocone morphism}
    \label{fig:falgcoconemorphism}
}
\end{figure}

    It suffices to show two things:
    \begin{enumerate}[label=(\arabic*),nosep]
        \item The unique cocone morphism from $\colim(F_{Seq})$ to $(Y, \gamma_i)_{i \in \omega}$ is an $F$-algebra morphism.
        \item Every $F$-algebra morphism from $(LFP(F), in_F)$ to $(Y, g)$ is an $F_{Seq}$-cocone morphism from $\colim(F_{Seq})$ to $(Y, \gamma_i)_{i \in \omega}$.
    \end{enumerate}

    Proving (1) will show that there exists an $F$-algebra morphism from $(LFP(F), in_F)$ to $(Y, g)$.  Proving (2) will show that this morphism is unique because there is a unique cocone morphism from the colimit $\colim(F_{Seq})$ to $(Y, \gamma_i)_{i \in \omega}$.

    We begin by proving (1).  Let $m : LFP(F) \to Y$ be the unique cocone morphism from $\colim(F_{Seq})$ to $(Y, \gamma_i)_{i \in \omega}$.  We want to know that $m$, which is the unique morphism satisfying
    \begin{equation}\label{eq:falgcoconemorphism}
        m \circ \phi_i = \gamma_i
    \end{equation}
    (see Figure \ref{fig:falgcoconemorphism}) is an $F$-algebra morphism, i.e., the square shown in Figure \ref{fig:initialfalg} commutes with $X = Y$ and $f = g$, or
    \[ m \circ in_F = g \circ F(m) \]
    
    For this proof, recall that $in_F$ is a cocone isomorphism between the two cocones
    \[(F(LFP(F)), \{\emptyset_{F(LFP(F))}\} \cup \{ F(\phi^{i}) : F^{i+1}(\emptyset) \to F(LFP(F))\})\] and 
    \[(LFP(F), \phi_i : F^i(\emptyset) \to LFP(F))_{i \in \omega} \]
    That is to say, $(F(LFP(F)), \{\emptyset_{F(LFP(F))}\} \cup \{ F(\phi^{i})\})$ is an initial cocone of $F_{Seq}$. The composition $m~{\circ}~in_F$ is also a cocone morphism, in this case from $(F(LFP(F)), \{\emptyset_{F(LFP(F))}\} \cup \{ F(\phi^{i})\})$ to $(Y, \gamma_i)_{i \in \omega}$, because it is a composition of cocone morphisms.  By initiality of the cocone $(F(LFP(F)), \{\emptyset_{F(LFP(F))}\} \cup \{ F(\phi^{i})\})$, $m \circ in_F$ is the unique cocone morphism from $(F(LFP(F)), \{\emptyset_{F(LFP(F))}\} \cup \{ F(\phi^{i})\})$ to $(Y, \gamma_i)_{i \in \omega}$.  We will prove that $g \circ F(m)$
    is also such a cocone morphism, and therefore
    \[ g \circ F(m) = m \circ in_F \]

    For $g \circ F(m)$ to be a cocone morphism, we need that
    \[ g \circ F(m) \circ \emptyset_{F(LFP(F))} = \gamma_0 \]
    and that
    \[ g \circ F(m) \circ F(\phi_i) = \gamma_{i+1} \]
    for all $i$.  The first equation holds trivially.  The second we prove by the following chain of equalities.
    \begin{align*}
        g \circ F(m) \circ F(\phi_i) &= g \circ F(m \circ \phi_i) && \textsf{(By functoriality of $F$)} \\
        &= g \circ F(\gamma_i) && \textsf{(By equation \eqref{eq:falgcoconemorphism})} \\
        &= \gamma_{i+1} && \textsf{(By definition of $\gamma_{i+1}$)}
    \end{align*}

    Thus, $g \circ F(m)$ is also the unique cocone morphism from $(F(LFP(F)), \{\emptyset_{F(LFP(F))}\} \cup \{ F(\phi^{i})\})$ to $(Y, \gamma_i)_{i \in \omega}$. We therefore have that
    \[ g \circ F(m) = m \circ in_F \]
    and $m$ is an $F$-algebra morphism.

    Next, we prove (2): every $F$-algebra morphism from $(LFP(F), in_F)$ to $(Y, g)$ is an $F_{Seq}$-cocone morphism from $\colim(F_{Seq})$ to $(Y, \gamma_i)_{i \in \omega}$.

    Let $m$ be an $F$-algebra morphism from $(LFP(F), in_F)$ to $(Y, g)$.  That is, $m$ is a morphism such that
    \begin{equation}\label{eq:mfalgmorphism}
        m \circ in_F = g \circ F(m)
    \end{equation}
    We will prove directly that $m$ is a cocone morphism from $(LFP(F), \phi_i)_{i \in \omega}$ to $(Y, \gamma_i)_{i \in \omega}$ by showing for all $i$ that
    \[ m \circ \phi_i = \gamma_i \]
    We proceed by induction on $i$.
    \begin{description}
        \item[Base Case:]
            \[ m \circ \phi_i = \gamma_0 \]
            Each of the above morphisms are the unique morphism from an initial object to $Y$, thus the equality holds.
        \item[Inductive Case:] Assume $m \circ \phi_n = \gamma_n$.  Then
        \begin{align*}
            m \circ \phi_{n+1} &= m \circ in_F \circ F(\phi_n) && \textsf{(Because $in_F$ is a cocone morphism)} \\
            &= g \circ F(m) \circ F(\phi_n)   && \textsf{(Because $m$ is an $F$-algebra morphism)} \\
            &= g \circ F(m \circ \phi_n) && \textsf{(By Functoriality of $F$)} \\
            &= g \circ F(\gamma_n) && \textsf{(By the inductive hypothesis)} \\
            &= \gamma_{n+1} && \textsf{(By definition of $\gamma_{n+1}$)}
        \end{align*}
    \end{description}
    And we have that every $F$-algebra morphism is a cocone morphism from $(LFP(F), \phi_i)_{i \in \omega}$ to $(Y, \gamma_i)_{i \in \omega}$. 
    
    We have just shown that every $F$-algebra induces a cocone morphism that is an $F$-algebra morphism, and that every $F$-algebra morphism is a cocone morphism with the same domain and codomain as the cocone morphism induced by the given $F$-algebra.  Because this domain is initial in the category of cocones, the $F$-algebra morphism is unique.  Therefore $(LFP(F), in_F)$ is initial in the category of $F$-algebras.
\end{proof}

\begin{corollary}\label{cor:inexistsandisiso}
    If $F$ is $\omega$-cocontinuous, then the initial $F$-algebra exists and it is an isomorphism.
\end{corollary}

\begin{proof}
    If $F$ is $\omega$-cocontinuous, then the initial $F$-algebra exists by Proposition \ref{prop:ininitial}.  It is an isomorphism by construction.
\end{proof}

Next, we show that the functorial action of $LFP$ always results in an algebra morphism.

\begin{prop}\label{prop:lfpalgmorphism}
    Let $F, G : \mcC \to \mcC$ be $\omega$-cocontinuous functors.  Let $\eta : F \naturalto G$.  Then $LFP(\eta)$ is an $F$-algebra morphism from the initial $F$-algebra $(LFP(F), in_F)$ to $(LFP(G), in_G \circ \eta_{LFP(G)} : F(LFP(G)) \to LFP(G))$.
\end{prop}

\begin{proof}
    We refer again to the diagram shown in Figure \ref{fig:lfpmorphism} for this proof.  We know that $LFP(\eta)$ is the unique morphism satisfying for all $n$
    \begin{equation}
    \label{eq:LFPeta}      
 LFP(\eta) \circ \phi_n = \psi_n \circ \eta^n_\emptyset
     \end{equation}
    We will show that the morphism $in_G \circ \eta_{LFP(G)} \circ F(LFP(\eta)) \circ in^-_F$ satisfies equation \eqref{eq:LFPeta} for all $n$, and therefore
    \[ LFP(\eta) = in_G \circ \eta_{LFP(G)} \circ F(LFP(\eta)) \circ in^-_F \]
    That is, the diagram shown in Figure \ref{fig:lfpalgmorphism} commutes.
    We do this using the following chain of equalities:
   {\scriptsize \begin{align*}
        &in_G \circ \eta_{LFP(G)} \circ F(LFP(\eta)) \circ in^-_F \circ \phi_n  \\
        &= in_G \circ \eta_{LFP(G)} \circ F(LFP(\eta)) \circ F(\phi_n) \circ f_{n \to n+1} && \textsf{(By commutativity of the diagram shown in Figure \ref{fig:initalg_in})} \\
        &= in_G \circ \eta_{LFP(G)} \circ F(LFP(\eta) \circ \phi_n) \circ f_{n \to n+1} && \textsf{(By functoriality of $F$)} \\
        &= in_G \circ \eta_{LFP(G)} \circ F(\psi_n \circ \eta^n_\emptyset) \circ f_{n \to n+1} && \textsf{(Because $LFP(\eta)$ is a cocone morphism, see Figure \ref{fig:lfpmorphism})} \\
        &= in_G \circ G(\psi_n \circ \eta^n_\emptyset) \circ \eta_{F^n(\emptyset)} \circ f_{n \to n+1} && \textsf{(By naturality of $\eta$)} \\
        &= in_G \circ G(\psi_n) \circ G(\eta^n_\emptyset) \circ \eta_{F^n(\emptyset)} \circ f_{n \to n+1} && \textsf{(By functoriality of $G$)} \\
        &= in_G \circ G(\psi_n) \circ \eta^{n+1}_\emptyset \circ f_{n \to n+1}
        && \textsf{(By definition of $\eta^{n+1}$)} \\
        &= \psi_{n+1} \circ \eta^{n+1}_\emptyset \circ f_{n \to n+1} && \textsf{(By commutativity of the diagram shown in Figure \ref{fig:initalg_in} with $F=G$)} \\
        &= \psi_n \circ \eta^n_\emptyset && \textsf{(By commutativity of the diagram shown in Figure \ref{fig:lfpmorphism})}
    \end{align*} }

    The above sequence of equalities shows that $in_G \circ \eta_{LFP(G)} \circ F(LFP(\eta)) \circ in^-_F$ is a cocone morphism from the colimit $\colim(F_{Seq})$ to the cocone $(LFP(G), \psi_i \circ \eta^i_\emptyset)_{i \in \omega}$.  We know that the unique cocone morphism between these two cocones is $LFP(\eta)$.  Therefore, 
    \[ LFP(\eta) = in_G \circ \eta_{LFP(G)} \circ F(LFP(\eta)) \circ in^-_F \]
    and the diagram shown in Figure \ref{fig:lfpalgmorphism} commutes.
    Indeed, $LFP(\eta)$ is an $F$-algebra morphism.
    
\end{proof}

\begin{figure}[htbp]
        \centering{
        \begin{tikzcd}
	{F(LFP(F))} && {LFP(F)} \\
	\\
	{F(LFP(G))} & {G(LFP(G))} & {LFP(G)}
	\arrow["{in_F}", from=1-1, to=1-3]
	\arrow["{F(LFP(\eta))}"', from=1-1, to=3-1]
	\arrow["{LFP(\eta)}", from=1-3, to=3-3]
	\arrow["{\eta_{LFP(G)}}"', from=3-1, to=3-2]
	\arrow["{in_G}"', from=3-2, to=3-3]
        \end{tikzcd}
        }
        \caption{$LFP$'s functorial action results in an algebra morphism}
        \label{fig:lfpalgmorphism}
\end{figure}

In this subsection, we showed that the least fixed point functor $LFP$ gives rise to the initial algebra of a functor to which it is applied. We now turn to the nontrivial isomorphisms that arise between objects in $\mcC$ under applications of $LFP$.

\section{Properties of the least fixed point functor}

In this section, we explore several isomorphisms induced by the $LFP$ functor.  In particular, we describe two interesting phenomena: first, that applying $LFP$ to a composition of functors allows one functor to be ``pulled out'' of the application of $LFP$ and second, that iterated applications of $LFP$ can be collapsed to a single application.  In Chapter \ref{sec:ADTs}, we will leverage these results to establish corresponding isomorphisms between data types.

\subsection{Least fixed points of a composition of functors}

We begin by formalizing the first of these phenomena: a nontrivial isomorphism that arises when $LFP$ is applied to a composition of two functors. The corresponding theorem is stated below.

Let $\mcC$ be an $\omega$-cocomplete category.

\Needspace{5\baselineskip}
\begin{theorem}\label{theorem:FGcolimit}
    For all $\omega$-cocontinuous functors $F, G : \mcC \to \mcC$
    \[ F(LFP(G\circ F)) \simeq LFP(F \circ G) \]
    and
    \[ G(LFP(F \circ G)) = LFP(G \circ F)\]
\end{theorem}

\begin{proof}
Let $F,G :\mcC \to \mcC$ be given.

Consider the product category $\mcC \times \mcC$ whose objects are pairs of objects in $\mcC$ and whose morphisms are pairs of morphisms in $\mcC$. By Theorem \ref{thm:productcolmits}, this category has colimits and, therefore, least fixed points.

Let $Swap(X, Y) = (Y, X)$ be the endofunctor on $\mcC \times \mcC$ that swaps the components of its argument.
Note that 
 \begin{align*}
    (Swap \circ (F \times G) \circ Swap)(X, Y)
    &= Swap((F \times G)(Swap(X, Y))) \\
    &= Swap((F \times G)(Y, X))  \\
    &= Swap(F\ Y, G\ X) \\
    &= (G\ X, F\ Y) \\
    &= (G \times F) (X,Y)
\end{align*}
That is, 
\begin{equation} \label{eq:swaplemma}
Swap \circ (F \times G) \circ Swap = G \times F
\end{equation}

Let $H : \mcC \times \mcC \to \mcC$ be the functor defined by
    \[ H(X,Y) = (F \times G) \circ Swap \]
Observe that 
\begin{align*}
    H \circ H
    &= (F \times G) \circ Swap \circ (F \times G) \circ Swap \\
    &= (F \times G) \circ (Swap \circ (F \times G) \circ Swap) \\
    &= (F \times G) \circ (G \times F) && \textsf{(by equation \eqref{eq:swaplemma})}\\
    &= (F \circ G) \times (G \circ F)
\end{align*}
    That is, $H \circ H$ is the product of two functors $F \circ G$ and $G \circ F$ from $\mcC$ to $\mcC$.  
    
    By Theorem \ref{thm:productcolmits}, colimits in product categories are computed coordinatewise, so 
\begin{equation} \label{eq:HcircH}
     LFP(H \circ H) \simeq (LFP(F \circ G), LFP(G \circ F)) 
\end{equation}

    To finish the proof, let $in_H : H(LFP(H)) \to LFP(H)$ be the morphism component of $H$'s initial algebra.  Recall that $in_H$ is an isomorphism by construction.  It therefore follows that
        \begin{align*}
        &(F(LFP(G \circ F)), G(LFP(F \circ G)))\\
        &\qquad= ((F \times G) \circ Swap)(LFP(F \circ G), LFP(G \circ F)) && \textsf{(by the definition of $Swap$)} \\
        &\qquad= H(LFP(F \circ G), LFP(G \circ F)) && \textsf{(by the definition of $H$)} \\
          &\qquad\simeq H(LFP(H \circ H))         && \textsf{(by equation \eqref{eq:HcircH})} \\
          &\qquad\simeq H(LFP(H))            && \textsf{(By Corollary \ref{cor:omegacofinal} with $f(x) = 2x$)} \\
          &\qquad\simeq LFP(H) && \textsf{(By the definition of $in_H$)} \\
          &\qquad\simeq LFP(H \circ H) && \textsf{(By Corollary \ref{cor:omegacofinal} with $f(x) = 2x$)} \\
          &\qquad\simeq (LFP(F \circ G), LFP(G \circ F))  && \textsf{(by equation \eqref{eq:HcircH})}
        \end{align*}

       The above isomorphism is an isomorphism in the product category $\mcC \times \mcC$, so it is composed of a pair of isomorphisms in $\mcC$.  The first of these isomorphisms has type $F(LFP(G \circ F)) \to LFP(F \circ G)$ and the second has type $G(LFP(F \circ G)) \to LFP(G \circ F))$.  That is, they are exactly the isomorphisms we wanted to show exist.
\end{proof}

With the proof complete, we have shown that $LFP$ applied to a composition of functors indeed satisfies the stated isomorphism.

\subsection{Iterated fixed points}

In this subsection, we show that iterated applications of the $LFP$ functor can be collapsed to a single application.  We make use of some additional notation and two supporting lemmas for the proof of this result.

\begin{notation}
    From this point forward, we denote
    \[ LFP(\lambda X.F(X))\]
    as 
    \[ LFP_X(F(X)) \]
\end{notation}

For example, we write $LFP_X(LFP_Y(F(X,Y)))$ for $LFP(\lambda X. LFP (\lambda Y. F(X,Y)))$. \\

In this subsection, we will prove that
\[ LFP_X(LFP_Y(F(X,Y))) \simeq LFP_X(F(X,X)) \]
Note that for the above equation to be well-defined, $LFP_Y(F(-,Y))$ must be an omega cocontinuous functor. We therefore prove the following lemma with reference to Johann and Polonsky \cite{jp2021} and Mac Lane \cite{maclane1998categories}.

\Needspace{4\baselineskip}
\begin{lemma}
    Given any $\omega$-cocontinuous functor $F : \mcD \times \mcD \to \mcC$, The functor $H$ defined as
    \[ H(X) = LFP_Y(F(X, Y)) \]
    is $\omega$-cocontinuous.
\end{lemma}

\begin{proof}
    We want to show that 
    \[ \colim_{i \in \omega}(HX_i) = H(\colim_{i \in \omega} X_i) \]
    We prove this as follows:
    \begin{align*}
        \colim_{i \in \omega}(HX_i) 
        &= \colim_{i \in \omega}(LFP(\lambda Y.F(X_i,Y)))
        &&\textsf{(By definition of $H$)} \\
        &= LFP(\colim_{i \in \omega}(\lambda Y.F(X_i, Y))
        &&\textsf{($LFP$ is $\omega$-cocontinuous by Theorem 5 in \cite{jp2021})} \\
        &= LFP(\lambda Y. F(\colim_{i \in \omega}(X_i), Y)) 
        &&\textsf{(by Theorem V.3.1 in \cite{maclane1998categories})}\\
        &= H(\colim_{i \in \omega} X_i) &
        &\textsf{(By definition of $H$)} 
    \end{align*}
The penultimate equality in the above sequence is sometimes stated as ``colimits in functor categories are computed pointwise''.
\end{proof}

\begin{lemma}\label{lemma:inisnatural}
    Let $F : \mcC \times \mcC \to \mcC$ be a functor on $C$ that is $\omega$-cocontinuous in both of its arguments.

    Let
    \[ K_X(Y) = F(X, Y) \]

    Then $in_{K_X} : K_X(LFP(K_X)) \to LFP(K_X)$ is a natural transformation in its variable $X$.  That is, it is a natural transformation from the functor $\lambda X. K_X(LFP(K_X))$ to the functor $\lambda X. LFP(K_X)$
\end{lemma}

\begin{proof}
    Let $f : X \to Z$ be an arbitrary morphism in $\mcC$.  Then $K_f = F(f, -)$ is a natural transformation from $\lambda Y.F(X, Y)$ to $\lambda X.F(Z, Y)$.  We want to show that $in_{K_X}$ is natural in its variable $X$, i.e., that the diagram shown in Figure \ref{fig:inisnatural} commutes.

    \begin{figure}[htbp]
        \centering{
        \begin{tikzcd}
	{K_X(LFP(K_X))} && {LFP(K_X)} \\
	{K_Z(LFP(K_Z))} && {LFP(K_Z)}
	\arrow["{in_{K_X}}", from=1-1, to=1-3]
        \arrow["{F(f, LFP(K_f))}", from=1-1, to=2-1]
        \arrow["{LFP(K_f)}", from=1-3, to=2-3]
        \arrow["{in_{K_Z}}", from=2-1, to=2-3]
        \end{tikzcd}
        }
        \caption{$in_{K_X}$ is a component of a natural transformation}
        \label{fig:inisnatural}
    \end{figure}

    By application of Proposition \ref{prop:lfpalgmorphism}, we obtain that
    \[ LFP(K_f) \circ in_{K_X} = in_{K_Z} \circ F(f, LFP(K_Y)) \circ K_X(LFP(K_f)) \]
    By definition of $K_X(LFP(K_f))$ and then by functoriality of $F$, we have
    \[ LFP(K_f) \circ in_{K_X} = in_{K_Z} \circ F(f, LFP(K_Y)) \circ F(X, LFP(K_f)) = in_{K_Z} \circ F(f, LFP(K_f)) \]
    which is exactly what we wanted to show.
\end{proof}

We can now prove the main result of this subsection.  It is stated below.

\begin{prop}\label{prop:iteratedlfp}
    Let $F : \mcC \times \mcC \to \mcC$ be $\omega$-cocontinuous in each of its arguments.  Then
    \begin{equation}
    \label{eq:LFPXY}
    LFP_X(LFP_Y(F(X,Y))) \simeq LFP_X(F(X,X))
    \end{equation}
\end{prop}

\begin{proof}
    Let
    \begin{align*}
        H(X) &= LFP_Y(F(X, Y)) \\
        G(X) &= F(X,X) \\
        K_X(Y) &= F(X,Y) \\
        A &= LFP(H)
    \end{align*}
    We will show that 
    \[ LFP(G) = A \]

    First, we will give $A$ the structure of a $G$-algebra.  Then, we will show that there exists a unique $G$-algebra morphism from this $G$-algebra structure on $A$ to any other $G$-algebra.  Since initial algebras, like initial objects in any category, are unique up to isomorphism, $A$ must be isomorphic to $LFP(G)$, proving \eqref{eq:LFPXY}.

    First, note that $H(X)$ is a least fixed point, so for all $X \in \mcC$
    \[ in_{K_X} : F(X, H(X)) \to H(X) \]
    and $in_{K_X}$ is an isomorphism.
    That is, $H(X)$ is $LFP(K_X)$.

    Next, note that $A$ is the carrier of the initial algebra of $H$. That is,
    \[ in_H : H(A) \to A\]

    We can now impose the following $G$-algebra structure on $A$ (see Figure \ref{fig:fxxalgmorphism}):
    \[ in_H \circ in_{K_{A}} \circ F(A, in^-_H) : F(A, A) \to A \]
    Note that because each component of the above morphism is an isomorphism, the composition remains an isomorphism.

    \begin{figure}
        \centering
        {
        \begin{tikzcd}
	{F(A, A)} && {F(A, H(A))} && {H(A)} & A \\
	\\
	{F(C,C)} &&&&& C
	\arrow["{F(A, in_H^-)}", from=1-1, to=1-3]
	\arrow["{F(g,g)}", from=1-1, to=3-1]
	\arrow["{in_{K_A}}", from=1-3, to=1-5]
	\arrow["{in_H}", from=1-5, to=1-6]
	\arrow["g", from=1-6, to=3-6]
	\arrow["\gamma", from=3-1, to=3-6]
        \end{tikzcd}
        }
        \caption{A $G$-algebra morphism from a $G$-algebra structure on $LFP(H)$}
        \label{fig:fxxalgmorphism}
    \end{figure}

    Next, we will show that given any $G$-algebra $(C, \gamma : F(C, C) \to C)$, there exists a unique $G$-algebra morphism $g : (A, in_H \circ in_{K_{A}} \circ F(A, in^-_H)) \to (C, \gamma)$.

    We want a morphism $g$ from $A$ to $C$ such that the diagram shown in Figure \ref{fig:fxxalgmorphism} commutes.
    First, we can obtain a morphism $g : A \to C$ through the following two-step process:
    \begin{enumerate}[nosep]
        \item Notice that $\gamma : F(C,C) \to C$ is a $K_C$-algebra, so there exists a unique morphism $f : H(C) \to C$ such that the diagram shown in Figure \ref{fig:kcalgmorphism} commutes.

        \begin{figure}
        \centering
        {
        \begin{tikzcd}
	{F(C, H(C))} && {H(C)} \\
	{F(C, C)} && C
	\arrow["{in_{K_C}}", from=1-1, to=1-3]
	\arrow["{F(C, f)}", from=1-1, to=2-1]
	\arrow["f", from=1-3, to=2-3, dashed]
	\arrow["\gamma", from=2-1, to=2-3]
        \end{tikzcd}
        }
        \caption{The unique $K_C$-algebra morphism  $f$ induced by $\gamma$}
        \label{fig:kcalgmorphism}
        \end{figure}

        \item Notice that $f$ provides an $H$-algebra structure on $C$, so there exists a unique morphism $g : A \to C$ such that the diagram shown in Figure \ref{fig:halgmorphism} commutes.

        \begin{figure}
        \centering
        {
        \begin{tikzcd}
	{H(A)} && A \\
	{H(C)} && C
	\arrow["{in_H}", from=1-1, to=1-3]
	\arrow["{H(g)}", from=1-1, to=2-1]
	\arrow["g", from=1-3, to=2-3, dashed]
	\arrow["f", from=2-1, to=2-3]
        \end{tikzcd}
        }
        \caption{The unique $H$-algebra morphism $g$ induced by $f$}
        \label{fig:halgmorphism}
        \end{figure}
    \end{enumerate}

    Now we will show that the morphism $g$ that we just defined is a $G$-algebra morphism.  For this to be true, the diagram shown in Figure \ref{fig:fxxalgmorphism} must commute. That is, we must have that
    \[ g \circ in_H \circ in_{K_A} \circ F(A, in^-_H) = \gamma \circ F(g, g) \]
    We can show this to be the case by the following chain of equations:
    {\small \begin{align*}
        &g \circ in_H \circ in_{K_A} \circ F(A, in^-_H) \\
        &= f \circ H(g) \circ in_{K_A} \circ F(A, in^-_H) &&  \textsf{(By commutativity of the diagram shown in Figure 
        \ref{fig:halgmorphism})} \\
        &= f \circ LFP(K_g) \circ in_{K_A} \circ F(A, in^-_H) &&\textsf{(By the definitions of $H$ and of $K_A$)} \\
        &= f \circ in_{K_C} \circ F(g, LFP(K_g)) \circ F(A, in^-_H) &&\textsf{(By naturality of $in_{K_-}$, see Lemma \ref{lemma:inisnatural})} \\
        &= \gamma \circ F(C, f) \circ F(g, LFP(K_g)) \circ F(A, in^-_H) &&\textsf{(By commutativity of the diagram shown in Figure \ref{fig:kcalgmorphism})} \\
        &= \gamma \circ F(g, f \circ LFP(K_g) \circ in^-_H) &&\textsf{(By functoriality of $F$)}\\
        &= \gamma \circ F(g, f \circ H(g) \circ in^-_H) &&\textsf{(By the definitions of $K$ and of $H$)} \\
        &= \gamma \circ F(g, g) &&\textsf{(By commutativity of the diagram shown in Figure \ref{fig:halgmorphism})}
    \end{align*}}
    Therefore, $g$ is a $G$-algebra morphism.  It remains to show that $g$ is the unique $G$-algebra morphism.
    
    The uniqueness of $g$ will follow from the fact that \emph{any} $G$-algebra morphism $m : A \to C$ from the $G$-algebra structure we have imposed on $A$ to the given $(C, \gamma)$ induces a $K_A$-algebra morphism $m \circ in_H$ from the initial $K_A$-algebra $(A, in_{K_A})$ to the $K_A$-algebra $(C, F(m,C) \circ \gamma)$. 
    
    Recall that $m \circ in_H$ is a $K_A$-algebra morphism if and only if
    \[ m \circ in_H \circ in_{K_A} = \gamma \circ F(m, C) \circ F(A, m \circ in_H)\]
    By functoriality of $F$, the above equality is logically equivalent to
    \[ m \circ in_H \circ in_{K_A} = \gamma \circ F(m, m) \circ F(A, in_H) \]
    which is true because $m$ is a $G$-algebra morphism and makes the diagram shown in Figure \ref{fig:fxxalgmorphism} commute with $g = m$.

    Because $m \circ in_H$ is a morphism out of the initial $K_A$-algebra and is therefore unique, we have that, given any two $G$-algebra morphisms $g$ and $g'$ from $(A, in_H \circ in_{K_{A}} \circ F(A, in^-_H))$ to $(C, \gamma)$,
    \[ g \circ in_H = m \circ in_H = g' \circ in_H \]
    $in_H$ is an isomorphism, so it is an epimorphism.  Therefore,
    \[ g = g' \]
    and there is a \emph{unique} $G$-algebra morphism from $(A, in_H \circ in_{K_{A}} \circ F(A, in^-_H))$ to any $G$-algebra.  Because initial algebras are unique up to isomorphism,
    \[ LFP_X(LFP_Y(F(X, Y))) = A \simeq LFP(G) = LFP_X(F(X,X)) \qedhere \]
\end{proof}

In this chapter, we established the necessary background to interpret recursive types as least fixed points (and also initial algebras) of functors.  We also proved theorems that we will later use to prove the existence of isomorphisms between categorical interpretations of recursive types.  Notice that in the last section of this Chapter, the only assumption we consistently made of the category in which we are working is that it is $\omega$-cocomplete.  This will be useful later, when we want to prove that certain isomorphisms hold in general, regardless of the setting in which certain types are interpreted.

\chapter{First-order Equational Theories and Their Models}\label{sec:eqtheory}

In this chapter we introduce several basic notions from universal algebra including signatures, equational theories, and models.  Later, we apply these concepts to the concrete algebraic structures of rings and semirings.  We assume the reader has foundational knowledge in abstract algebra including the definitions of ring and semiring, and the concepts of quotient algebras and ideals.

A signature specifies the operations that are available in a particular algebraic structure. It consists of a collection of operation symbols together with their arities. At this stage the symbols are purely formal---they simply indicate that operations of certain shapes exist, without yet specifying what those operations do.

An equational theory builds on a signature by specifying laws that the operations should satisfy. These laws are expressed as equations between terms built from the operation symbols and variables. In this way, an equational theory describes the algebraic behavior expected of the operations.

Finally, a model of an equational theory is a concrete mathematical structure that interprets the operation symbols as actual operations on a set and in which all of the specified equations hold. Models therefore give concrete structures in which the abstract operations and laws of the theory hold.

The goal of this chapter is to introduce the reader to these ideas for later use in describing, manipulating, and reasoning about the algebraic structure present in certain data types.

\section{Signatures and algebras}
\begin{definition}
A \emph{first-order signature} is a set $\Sigma$ of function symbols labeled by non-negative integers called their \emph{arity}.
\[ \Sigma = \{f_1^{a_1}, f_2^{a_2}, ..., f_k^{a_k}\}\]
\end{definition}

For the rest of this chapter, let $\Sigma$ be a first-order signature as above.

\begin{definition}
    A \emph{$\Sigma$-algebra} is a pair $(A, \{F_i : D^{a_i} \to D\ |\ f_i \in \Sigma\})$ where $D$ is a set and $\{F_i : A^{a_i} \to A\ |\ f_i \in \Sigma\}$ is a set of $k$ functions, one for each function symbol in $\Sigma$.  
\end{definition}

\begin{definition}\label{def:sig_hom}
    Given two $\Sigma$-algebras, $(A, \{\alpha_1, \alpha_2, ..., \alpha_k\})$ and $(B, \{\beta_1, \beta_2, ..., \beta_k\})$, a \emph{$\Sigma$-algebra homomorphism} is a function $h : A \to B$ such that for all $\alpha_i$,
    \[ h(\alpha_i(m_1, m_2, ..., m_{a_i})) = \beta_i(h(m_1), h(m_2), ..., h(m_{a_i})) \]
\end{definition}



\begin{definition}
The \emph{initial $\Sigma$-algebra} $\Sigma^*$ consists of the set generated by the following rule:

\[ \inferrule{
        t_1, t_2, ... , t_{a_i} \in \Sigma^*\ \ f_i \in \Sigma
    }{
        f_i(t_1, t_2, ... t_{a_i}) \in \Sigma^*
    } \]
together with the set of functions $\{f_1, ..., f_k\}$ where each $f_i$ maps $(t_1, ..., t_{a_i})$ to $f_i(t_1, ..., t_{a_i})$, with each $t_j \in \Sigma^*$. 
Note that this is an inductive definition, with base case obtained at those $f_i \in \Sigma$ with $a_i = 0$.
\end{definition}

\begin{prop}
    (Universal property of the initial $\Sigma$-algebra) 
    For every $\Sigma$-algebra $(A, \{F_i : A^{a_i} \to A\ |\ f_i \in \Sigma\})$, there exists a unique $\Sigma$-algebra homomorphism $\phi$ from $\Sigma^*$ to $B$; it is defined as follows:
\begin{equation} \label{eq:init-hom}
\phi(f_i(t_1, ..., t_{a_i})) = F_i(\phi(t_1), ... , \phi(t_{a_i}))
\end{equation}
\end{prop}

\begin{proof}
    By the definition of $\Sigma$-algebra homomorphism, \eqref{eq:init-hom} must be true of any homomorphism $\phi$ from $\Sigma^*$ into $(A, \{F_i\ |\ f_i \in \Sigma\})$.
    
    By induction on $t \in \Sigma^*$, this uniquely defines $\phi(t)$.
\end{proof}

Let $V$ be a set.

\begin{definition}
    The \emph{free $\Sigma$-algebra over $V$}, which is here denoted as $\Sigma^*_V$, is the set generated by the following rules:

\[ \inferrule{
        t_1, t_2, ... , t_{a_i} \in \Sigma^*_V\ \ f_i \in \Sigma
    }{
        f_i(t_1, t_2, ... t_{a_i}) \in \Sigma^*_V
    }\qquad
    \inferrule{
        v \in V
    }{
        v \in \Sigma^*_V
    }
\]
\noindent
with $\{f_1,\dots,f_k\}$ again giving it the structure of a $\Sigma$-algebra. 
\end{definition}

Note that $\Sigma^*_\emptyset = \Sigma^*$.

\begin{definition}
    Given a $\Sigma$-algebra $\Amod = (A, \{F_i\ |\ f_i \in \Sigma\})$, a $V$-\emph{environment} in $\Amod$ (or just ``environment'', when $V$ is easily inferred) is a function $\theta : V \to A$.  
\end{definition}
\newcommand{\msL}{\mathscr{L}}
\begin{prop}
    (Universal property of the free $\Sigma$-algebra over $V$) For any $\Sigma$-algebra $\Amod=(A, \{F_i\ |\ f_i \in \Sigma\})$ together with a $V$-environment $\theta$ in $\mathscr{A}$, there exists a unique $\Sigma$-algebra homomorphism $\semof{-}_\theta^\Amod : \Sigma^*_V \to \Amod$ satisfying $\semof{v}_\theta^\Amod = \theta(v)$ for each $v \in V$. This homomorphism is uniquely defined below.
\begin{align*}
    \semof{-}_\theta^\Amod &: \Sigma^*_V \to \Amod \\
    \semof{v}_\theta^\Amod &= \theta(v) && (v \in V)\\
    \semof{f_i(t_1,\dots,t_{a_i})}_\theta^\Amod 
    &= F_i(\semof{t_1}_\theta^\Amod,\dots,\semof{t_{a_i}}_\theta^\Amod) && (f_i \in \Sigma)
\end{align*}
\end{prop}

\begin{proof}
    By the definition of $\Sigma$-algebra homomorphism, these two equations must be true of any homomorphism $h$ from $\Sigma^*_V$ into $(A, \{F_i\ |\ f_i \in \Sigma\})$ satisfying  $x \in V \implies h(x) = \theta(x)$.
    
    By induction on $t \in \Sigma^*_V$, this uniquely defines $\semof{t}_\theta^\Amod$.
\end{proof}

When $\Amod$ is clear from the context, we just write $\semof{-}_\theta$ for $\semof{-}_\theta^\Amod$.

\begin{corollary}
    For all $s \in \Sigma^*_V$, $\semof{s}_{id_V} = s$.
\end{corollary}


%Let $\Lmod=(L, \{K_i\ |\ f_i \in \Sigma\})$ be a $\Sigma$-algebra, and let $V$ be a set. Every environment in $\Lmod$ extends, by the universal property of $\Sigma^*_V$, to a $\Sigma$-algebra homomorphism:

\section{Congruences and quotients}
\newcommand{\thof}[1]{\ensuremath{E^*_{#1}}}

Let $\Amod = (A, \{F_i : A^{a_i} \to A\ |\ f_i \in \Sigma\})$ be a $\Sigma$-algebra. Let $R \subseteq A \times A$ be a relation on $A$.

\begin{definition}
    $R$ is an \emph{equivalence relation} if $R$ is reflexive, symmetric, and transitive, i.e., $R$ validates the following rules.

    \[ 
    \inferrule{t \in A}{(t, t) \in R}\ \ \inferrule{(s, t) \in R}{(t, s) \in R}\ \ 
    \inferrule{(s, t) \in R\ \ (t, u) \in R}{(s, u) \in R}
    \] 
\end{definition}

\begin{definition}
    $R$ is a \emph{congruence} with respect to $\Amod$ if $R$ is an equivalence relation and $R$ validates the following rule.

    \[ \inferrule{(s_1, t_1) \in R\ \ ...\ \ (s_{a_i}, t_{a_i}) \in R\ \ f_i \in \Sigma}{(F_i(s_1, ..., s_{a_i}), F_i(t_1, ..., t_{a_i})) \in R} 
    \]
\end{definition}

\begin{definition} \label{def:rel-cong}
The \emph{congruence generated by $R$}, denoted $\sim_R$, is the least congruence on $\Amod$ containing $R$. 
When $(x,y) \in\ \sim_R$, we write $x \sim_R y$. It is defined by 
\[
    \inferrule{(s, t) \in R}{s \sim_R t}\ \ 
    \inferrule{t \in A}{t \sim_R t}\ \ \inferrule{s \sim_R t}{t \sim_R s}\ \ 
    \inferrule{s \sim_R t\ \ t \sim_R u}{s \sim_R u} 
\]
\[
    \inferrule{s_1 \sim_R t_1\ \ ...\ \ s_{a_i} \sim_R t_{a_i}\ \ f_i \in \Sigma}{F_i(s_1, ..., s_{a_i}) \sim_R F_i(t_1, ..., t_{a_i})} 
\] 
\end{definition}
Note that when $R$ is a congruence, $x \sim_R y$ if and only if $(x,y) \in R$.

\begin{notation}
    In the remainder of this document, we denote the equivalence class of $x$ with respect to the equivalence relation $R$ as $[x]_R$.
\end{notation}


For any $\Sigma$-algebra $\Amod = (A, \{F_i\ |\ f_i \in \Sigma\})$ and any relation $T$ that is a congruence with respect to $\Amod$, there exists a quotient $\Amod/T$.

\begin{prop}
    $\Amod/T$ has the structure of a $\Sigma$-algebra with operations defined as follows for every $f_i \in \Sigma$:
\begin{align*}
    F^T_i &: (A/T)^{a_i} \to A/T \\
    F^T_i([x_1]_T, \dots, [x_{a_i}]_T) &= [F_i(x_1, \dots, x_{a_i})]_T
\end{align*}
\end{prop}
\begin{proof}
    Because $T$ is a congruence, we know that 
    if $[x_i]_T = [y_i]_T$ for all $i$ from 1 to $a_i$, then
    \[ [F_i(x_1, \dots, x_{a_i})]_T = [F_i(y_1, \dots, y_{a_i})]_T\]
    The values of $F^T_i$ are thus independent of the representatives $x_i$ chosen from the given equivalence classes.  That is, $F^T_i$ is well-defined.
\end{proof}

\newcommand{\Bmod}{\mathscr{B}}
\begin{prop} \label{prop:univ-quot}
    (Universal property of the quotient $\Amod/T$) For any  $\Sigma$-algebra $\mathscr{B} = (B, \{G_i : B^{a_i} \to B\ |\ f_i \in \Sigma\})$ and $\Sigma$-algebra homomorphism $h : A \to B$ with $h(x)=h(y)$ for all $(x,y) \in T$, there exists a unique $\Sigma$-algebra homomorphism $\tilde{h} : A/T \to B$ such that
    \[\forall x \in A.\ \tilde{h}([x]_T) = h(x)\]
\end{prop}

\begin{proof}
    Let $(x, y) \in T$. Then 
    \[ \tilde{h}([x]_T) = h(x) = h(y) = \tilde{h}([y]_T) \]
    by hypothesis.  Thus, $\tilde{h}$ is well-defined.

    $\tilde{h}$ is the only $\Sigma$-algebra homomorphism with the property that
    \[\forall x \in A.\ \tilde{h}([x]_F) = h(x)\]
    because this equation completely determines $\tilde{h}$ in terms of $h$.
\end{proof}

\begin{definition}\label{def:quotient_congruence}
    If $T$ is a congruence on $A$ and $R \subseteq A \times A$ is a relation, then the \emph{quotient congruence of $R$ with respect to $T$}, denoted $[R]_T \subseteq \Amod/T$, is the congruence on the quotient algebra $A/T$ defined by
    \[ [R]_T = \{ ([x]_T,[y]_T) \mid x \sim_{R \cup T} y\} \]  
\end{definition}

If $T$ is clear from context, we write $[R]$ for $[R]_T$.

\begin{notation}
    If $\sigma : V \to A$ is an environment and $T \subseteq A \times A$ is a congruence, we write $\tilde{\sigma} : V \to A/T$ for the environment defined by
    \[ \tilde{\sigma}(v) = [\sigma(v)]_T\]
\end{notation}

\begin{lemma}\label{lemma:quotient_environments}
    Let $\Amod = (A, \{F_i\ |\ f_i \in \Sigma\})$ be a $\Sigma$-algebra. Let $T \subseteq A \times A$ be a congruence.  Let $\sigma : V \to A$ be an environment.  Then
    \[ \forall t \in \Sigma^*_V.\ [\semof{t}_\sigma]_T =\semof{t}_{\tilde{\sigma}} \]
\end{lemma}

\Needspace{3\baselineskip}
\begin{proof}
    We proceed by induction on $t \in \Sigma^*_V$.

    Let $t \in V$. Then

    \[ [\semof{t}_\sigma]_T = [\sigma(t)]_T = \semof{t}_{\tilde{\sigma}}\]

    Let $t = f_i(t_1, ..., t_{a_i})$. Suppose that for all $k$, $[\semof{t_k}]_T = \semof{t_k}_{\tilde{\sigma}}$.  Then

    \begin{align*}
        [\semof{t}_\sigma]_T &= [F_i(\semof{t_1}_\sigma, ..., \semof{t_{a_i}}_\sigma)]_T \\
        &= F^T_i([\semof{t_1}_\sigma]_T, ..., [\semof{t_{a_i}}_\sigma]_T) \\
        &= F^T_i(\semof{t_1}_{\tilde{\sigma}}, ..., \semof{t_{a_i}}_{\tilde{\sigma}}) \\
        &= \semof{f_i(t_i, ..., t_{a_i})}_{\tilde{\sigma}} \\
        &= \semof{t}_{\tilde{\sigma}} 
    \end{align*}
This completes the proof.
\end{proof}

Let $\Amod = (A, \{F_i\ |\ f_i \in \Sigma\})$ and $\Bmod = (B, \{ G_i\ |\ f_i \in \Sigma\})$ be $\Sigma$-algebras.  Let $h : \Amod \to \Bmod$ be a $\Sigma$-algebra homomorphism. 

\begin{definition}
    Given a $\Sigma$-algebra homomorphism $h$ from $\Amod$ to $\Bmod$, the \emph{kernel} of $h$, denoted $\ker(h)$, is a relation on $\Amod$ defined as
    \[ \ker(h) = \{(x, y)\ |\ h(x) = h(y)\} \]
\end{definition}

\begin{prop}
    $\ker(h)$ is a congruence on $\Amod$.
\end{prop}

\begin{proof}
    $\ker(h)$ is an equivalence relation because it relates two elements of $\Amod$ when they are equal in the image of $h$.

    Given $x_1, ..., x_{a_i}, y_1, ..., y_{a_i} \in \Amod$ and that $\forall i.\ (x_i, y_i) \in \ker(h)$, we have that $h(x_i) = h(y_i)$ for all $i$.  Then
    \[ G_i(h(x_1), ..., h(x_{a_i})) = G_i(h(y_1), ..., h(y_{a_i})) \]

    Because $h$ is a $\Sigma$-algebra homomorphism, the above implies
    \[ h(F_i(x_1, ..., x_{a_i})) = h(F_i(y_1, ..., y_{a_i})) \]
    which means that $(F_i(x_1, ..., x_{a_i}), F_i(y_1, ..., y_{a_i})) \in \ker(h)$, as required.
\end{proof}

\begin{prop}
The quotient $\Amod/\ker(h)$ is isomorphic to 
the image of $h$ in $\Bmod$: \[ \Amod/\ker(h)~\simeq~\mathsf{Im}(h)\]
\end{prop}
\begin{proof}
    By the universal property of quotients of $\Sigma$-algebras (Proposition \ref{prop:univ-quot}), there exists a unique $\Sigma$-algebra homomorphism $\tilde{h} : \Amod/\ker(h) \to \Bmod$ such that
    \[ \forall x \in A.\ \tilde{h}([x]_{\ker(h)}) = h(x) \]
    
    Trivially, the image of $\tilde{h}$ is the image of $h$.  Because every homomorphism is surjective on its image, $\tilde{h}$, when viewed as a homomorphism from $\Amod/\ker(h)$ to $\mathsf{Im}(h)$, is surjective.  We now must show that $\tilde{h}$ is injective.  
    
    Let $x, y \in A$ be such that $\tilde{h}([x]_{\ker(h)}) = \tilde{h}([y]_{\ker(h)})$.  By the definition of $\tilde{h}$, $h(x) = h(y)$.
    So $(x, y) \in \ker(h)$ and $[x]_{\ker(h)} = [y]_{\ker(h)}$.  Indeed, $\tilde{h}$ is injective.

    Because $\tilde{h}$, when viewed as a homomorphism from $\Amod/\ker(h)$ to $\mathsf{Im}(h)$ is both injective and surjective, we have that
    \[ \Amod/\ker(h) \simeq \mathsf{Im}(h) \qedhere \]
\end{proof}


\section{Theories and models}

\begin{definition}
    If $E \subseteq \Sigma^*_V \times \Sigma^*_V$ is a relation on $\Sigma^*_V$, then we call 
    $E$ a \emph{$\Sigma$-theory} over $V$.
\end{definition}

For the remainder of this section, let $E$ be an arbitrary $\Sigma$-theory over a set $V$.
    
\begin{definition}
    For an arbitrary set $W$, the \emph{saturation of $E$ with respect to $W$}, is a least congruence on $\Sigma^*_W$ containing every substitution instance of each $(s,t)\in E$ in $\Sigma^*_W$. 
    
    That is, it is the inductive relation on $A$ generated by the following rules:
    
\[
    \inferrule{(s, t) \in E \qquad \sigma : V \to \Sigma^*_W}{(\semof{s}_\sigma, \semof{t}_\sigma) \in \thof{W}}\ \ \inferrule{t \in \Sigma^*_W}{(t, t) \in \thof{W}}\ \ \inferrule{(s, t) \in \thof{W}}{(t, s) \in \thof{W}}
\]
\[
    \inferrule{(s, t) \in \thof{W}\ \ (t, u) \in \thof{W}}{(s, u) \in \thof{W}}\ \  \inferrule{(s_1, t_1) \in \thof{W}\ \ ...\ \ (s_{a_i}, t_{a_i}) \in \thof{W}\ \ f_i \in \Sigma}{(f_i(s_1, ..., s_{a_i}), G_i(t_1, ..., t_{a_i})) \in \thof{W}}
\]
\end{definition}

\begin{definition}\label{def:models}
A $\Sigma$-algebra $\Amod = (A, \{F_i\ |\ f_i \in \Sigma\})$ is an \textit{$E$-model}, written $\Amod \models E$, if the following sentence is true:
\[ \forall (s, t) \in E.\ \forall \theta : V \to A.\ \semof{s}_\theta = \semof{t}_\theta \]
\end{definition}

If $E = \{s = t\}$ is a singleton set, we write
\[ A \models s = t \]

\begin{lemma}\label{lemma:EtoE*}
    Let $(s, t) \in E^*_W$ for an arbitrary set $W$.  Let $\sigma$ be a $W$-environment in $\Lmod$.  %Then     
    \[ \Lmod \models E \implies \semof{s}_\sigma = \semof{t}_\sigma \]
\end{lemma}

\begin{proof}
    Assume $\Lmod \models E$.
    We show $\semof{s}_\sigma = \semof{t}_\sigma$ by induction on the proof of $(s,t) \in E^*_W$.

    \begin{description}
        \item[Base Cases: ]\hfill
        \begin{itemize}[nosep]
            \item Let $(s,t) \in E$. Then because $\Lmod \models E$, $\semof{s}_\sigma = \semof{t}_\sigma$.
            \item Suppose $(s,t) \in E^*_W$ follows by the reflexivity axiom with $s$ and $t$ 
            being the same element of $\Sigma^*_W$.  Then $\semof{s}_\sigma = \semof{t}_\sigma$.        
        \end{itemize}
        \item[Inductive Cases:]\hfill
        \begin{itemize}[nosep]
            \item Suppose $(s,t) \in E^*_W$ follows by the symmetry axiom with $(t,s) \in E^*_W$. By induction hypothesis, $\semof{t}_\sigma = \semof{s}_\sigma$.
            By symmetry of equality,  $\semof{s}_\sigma = \semof{t}_\sigma$. 
            \item Suppose $(s, v) \in E^*_W$ and $(v, t) \in E^*_W$.  By the induction hypothesis, we have that $\semof{s}_\sigma = \semof{v}_\sigma$
            and
            $\semof{v}_\sigma = \semof{t}_\sigma$.
            By transitivity of equality,
            $\semof{s}_\sigma = \semof{t}_\sigma$.
            \item Suppose $(s, t) \in E^*_W$ is proved by congruence with $s = f_i(s_1, ..., s_{a_i})$ and $t = f_i(t_1, ..., t_{a_i})$. By the induction hypothesis, we have $\semof{s_i}_\sigma = \semof{t_i}_\sigma$ for all $i$.  By definition of $\semof{-}_\sigma$, we have
            \[ \semof{s}_\sigma = \semof{f_i(s_1, ..., s_{a_i})}_\sigma = \semof{f_i(t_1, ..., t_{a_i})}_\sigma = \semof{t}_\sigma \qedhere\]
        \end{itemize}
    \end{description}
\end{proof}

\begin{definition}
    An \emph{$E$-model homomorphism} (e.g., a ring homomorphism) is a $\Sigma$-algebra homomorphism between $E$-models.
\end{definition}

Because both the domain and codomain of an $E$-model homomorphism are $E$-models, any such homomorphism preserves relations defined by $E$. Specifically, for any $(l=r) \in E$, and any $\theta : X \to A$, we have 
\[ \semof{l}_{h \circ \theta}^B = h(\semof{l}_\theta^A) = h(\semof{r}_\theta^A) = \semof{r}_{h \circ \theta}^B\]
For example, a semigroup is a model of the first-order signature $\Sigma_{SG} = \{*^2\}$ and the set of equations $E_{SG}=\{x * (y * z) = (x * y) * z\}$.  A $\Sigma$-algebra homomorphism $h$ between two $\Sigma_{SG}$-algebras that are also semigroups, i.e., have an \textit{associative} binary operation, also preserves associativity. 
\[h(x) * (h(y) * h(z)) = h(x * (y * z)) = h((x * y) * z) = (h(x) * h(y)) * h(z) \]

In the remainder of this chapter, let $\sigma_\emptyset$ be the empty environment, i.e., the environment with domain $\emptyset$.

\begin{definition}
    Given a $\Sigma$-theory $E$, the \emph{initial $E$-model} is $\Sigma^*/E^*_\emptyset$, the quotient of $\Sigma^*$ over the congruence $E^*_\emptyset$.
\end{definition}

Note that because $E^*_\emptyset$ is a congruence with respect to $\Sigma$, the quotient $\Sigma^*/E^*_\emptyset$ maintains a $\Sigma$-algebra structure.  The universal property of the initial $E$-model is inherited from this quotient $\Sigma^*/E^*_\emptyset$. 

\begin{prop}(Universal property of the initial $E$-model)
     For any other $\Sigma$-algebra $\Amod=(A, \{F_i\ |\ f_i \in \Sigma\})$, if $\Amod$ models $E$, then $\semof{-}_{\sigma_\emptyset}$ extends to an $E$-model homomorphism from the initial $E$-model $\Sigma^*/E^*_\emptyset$ to $\Amod$.  This homomorphism is uniquely defined by
\begin{equation}\label{eq:init_E_hom}
\semof{[t]_{\thof{\emptyset}}}_{\sigma_\emptyset} = \semof{t}_{\sigma_\emptyset}
\end{equation}
\end{prop}
\begin{proof}
    $\Amod$ models $E$ so, by Lemma \ref{lemma:EtoE*}, the action of $\semof{-}_{\tilde{\sigma_\emptyset}}$ on $\Sigma^*/E^*_\emptyset$ is well-defined.
\end{proof}

\begin{definition}
    For any fixed set $X$, there exists an \emph{initial $E$-model over $X$}, also called the \emph{free $E$-model over $X$}.  This model is $\Sigma^*_X/\thof{X}$.
\end{definition}

Again, note that $\Sigma^*_X/\thof{X}$ is a $\Sigma$-algebra because $\thof{X}$ is a congruence.

\Needspace{5\baselineskip}
\begin{prop}
    \label{prop:free-mod}
    (Universal property of the free $E$-model over $X$) For any other $\Sigma$-algebra $\Amod=(A, \{F_i\ |\ f_i \in \Sigma\})$, if $\Amod$ models $E$, then given any $X$-environment $\sigma$ in $\Amod$, $\semof{-}_\sigma$ extends to an $E$-model homomorphism from $\Sigma^*_X/\thof{X}$ to $\Amod$.  This homomorphism is uniquely defined by
    \begin{equation}\label{eq:free_e_model_hom}
    \semof{[t]_{\thof{X}}}_{\sigma} = \semof{t}_\sigma
    \end{equation}
\end{prop}

\begin{proof}
    $\Amod$ models $E$ so, by Lemma \ref{lemma:EtoE*}, the action of $\semof{-}_\sigma$ on $\Sigma^*/E^*_X$ is well-defined for an arbitrary environment $\sigma$.
\end{proof}

\section{Quotients of free models}

\begin{notation}
    If $E$ is a $\Sigma$-theory, we write $\fv{E} = \bigcup \{\fv{s} \cup \fv{t} \mid (s=t) \in E\}$, where $\fv{u}$ denotes the set of free variables occurring in $u$.
\end{notation}

\begin{definition}\label{def:entails}
Let $E_1, E_2$ be $\Sigma$-theories.
We say that $E_1$ \emph{entails} $E_2$, written $E_1 \models E_2$, if for every $\Sigma$-algebra  $\Amod$, 
\[ \Amod \models E_1 \implies \Amod \models E_2 \]
\end{definition}

When $E_2 = \{s = t\}$ is a singleton, we simply write $E_1 \models s = t$.

Note that $E \models s = t$ if and only if, for every $\Sigma$-algebra $\Amod$, every $(u, w) \in E$, and every environment $\rho : \fv{E} \to \Amod$, $\semof{u}_\rho = \semof{w}_\rho$, then for every $\Sigma$-algebra $\Amod$ and every environment $\rho : \fv{E} \to \Amod$, $\semof{s}_\rho = \semof{t}_\rho.$

\begin{prop}\label{prop:union_models}
Let $T$ and $E$ be $\Sigma$-theories.  Let $s,t \in \Sigma^*_V$. 

$T^*_V \cup E \models s = t$ if and only if $[s]_{T^*_V} \sim_{[E]_{T^*_V}} [t]_{T^*_V}$ in
 $\Sigma^*_V/T^*_V$, the free $T$-model over $V$.
    \end{prop}

\newcommand{\RA}{\Rightarrow}
\newcommand{\LA}{\Leftarrow}
\begin{proof}
\begin{description}
\item[$(\RA)$]
    Suppose $T^*_V \cup E \models s=t$.  We need to show that $[s]_{T^*_V} \sim_{[E]_{T^*_V}} [t]_{T^*_V}$.
    
    Recall that $\Sigma^*_V/T^*_V \models T^*_V$.

%     Let $\sigma : V \to \Sigma^*_V/T^*_V$ be an environment 
% satisfying $\semof{l}_\sigma = \semof{r}_\sigma$ for all 
% $(l=r) \in T^*_V \cup E$.

    Let $\sigma : V \to \Sigma^*_V$ be the environment defined by $\sigma(v) = v$.
    
    Let $\tilde{\sigma} : V \to \Sigma^*_V/T^*_V$ be the
    environment defined by $\tilde{\sigma}(v) = [v]_{T^*_V}$.

    Then by Lemma \ref{lemma:quotient_environments}, $[s]_{T^*_V} = \semof{s}_\sigma$, 
    and $[t]_{T^*_V} = \semof{t}_\sigma$.

    Let $\hat{\sigma} : V \to \left(\Sigma^*_V/T^*_V\right)/[E]_{T^*_V}$ 
    be the composition of $\tilde{\sigma}$ with the quotient map:
    \[ \hat \sigma(v) = [\tilde \sigma(v)]_{[E]_{T^*_V}}\]

    Let $l = r \in T^*_V \cup E$.  By Lemma \ref{lemma:quotient_environments}, we have
    \begin{align*}
        \semof{l}_{\hat \sigma} &=
        [\semof{l}_{\tilde \sigma}]_{[E]_{T^*_V}} \\
        &= [[\semof{l}_\sigma]_{T^*_V}]_{[E]_{T^*_V}} \\
        &= [[l]_{T^*_V}]_{[E]_{T^*_V}} \\
        &= [l]_{T^*_V \cup E} && \textsf{(By Definition \ref{def:quotient_congruence})}\\
        &= [r]_{T^*_V \cup E} && \textsf{(By hypothesis)}\\
        &= [[r]_{T^*_V}]_{[E]_{T^*_V}} && \textsf{(By Definition \ref{def:quotient_congruence})}\\
        &= [[\semof{r}_\sigma]_{T^*_V}]_{[E]_{T^*_V}} \\
        &= [\semof{r}_{\tilde \sigma}]_{[E]_{T^*_V}} \\
        &= \semof{r}_{\hat \sigma}
    \end{align*}

    Since $T^*_V \cup E \models s=t$, $\semof{s}_{\tilde \sigma}
    = \semof{t}_{\tilde \sigma}$ (See Definitions \ref{def:models} and \ref{def:entails}).  Thus 
    \[ [[s]_{T^*_V}]_{[E]_{T^*_V}} =
    \semof{s}_{\tilde \sigma}
    = \semof{t}_{\tilde \sigma}
    = [[t]_{T^*_V}]_{[E]_{T^*_V}}
    \]
    and $[s]_{T^*_V} \sim_{[E]_{T^*_V}} [t]_{T^*_V}$.

    \item[$(\LA)$]
    Suppose $[s]_{T^*_V} \sim_{[E]_{T^*_V}} [t]_{T^*_V}$.

    Let $\Amod$ be a $\Sigma$-algebra, and suppose that 
    $\Amod \models T^*_V \cup E$.
    
    Let $\sigma : V \to \Amod$ be an environment.  By hypothesis, $\sigma$ validates every equation in $T^*_V \cup E$.  We need to show that 
    $\semof{s}_\sigma = \semof{t}_\sigma$.

    From the universal property of the free $T$-model over $V$, 
    let $h : \Sigma^*_V/T^*_V \to \Amod$ be the map induced 
    by $\sigma : V \to \Amod$ and the fact that $\Amod \models T$.
    Per Proposition \ref{prop:free-mod}, we now write 
    \[ \semof{s}_\sigma = 
    h ([s]_{T^*_V}), \qquad \semof{t}_\sigma = h([t]_{T^*_V})\] 
    
    Next, we observe that $\Amod$ satisfies the hypotheses on $\Bmod$ in
    Proposition \ref{prop:univ-quot} with $\Amod = \Sigma^*_V/T^*_V$, $F = [E]_{T^*_V}$,
    and $h$ as above.  The fact that $h(x)=h(y)$ for all 
    $x \sim_{[E]_{T^*_V}} y$ follows from the fact that 
    \[h([l]_{T^*_V}) = \semof{l}_\sigma = \semof{r}_\sigma 
    = h([r]_{T^*_V})\] 
    for all $(l=r) \in E$, with the middle equality obtained from $\Amod \models E$.

    We thus have $h([x]_{T^*_V}) = \tilde h([[x]_{T^*_V}]_{[E]_{T^*_V}})$ for all $x \in \Sigma^*_V/T^*_V$. 
    
    The claim now follows from the principal hypothesis on $s$ and $t$, as 
    \begin{align*}
    \semof{s}_\sigma &= h([s]_{T^*_V})  \\
    &= \tilde h([[s]_{T^*_V}]_{[E]_{T^*_V}}) \\
    &= \tilde h([[t]_{T^*_V}]_{[E]_{T^*_V}}) \\
    &= h([t]_{T^*_V})\\
    &= \semof{t}_\sigma
    \end{align*}
        
\end{description}
\end{proof}

\begin{notation}
    When either of the equivalent conditions of this proposition holds, we write $E \models_T s=t$.
\end{notation}

\begin{remark}
The second equivalent condition of Proposition \ref{prop:union_models} projects $E$ to equivalence classes of $T^*_V$.  Therefore, replacing any terms in $E$ with $T$-equivalents does not change the set of consequences: when $p \sim_{T^*_V} p'$ and $q \sim_{T^*_V} q'$, then 
\[ E, p=q \models_T s=t \iff E,p'=q' \models_T s=t \]
In light of this, we may sometimes write, e.g., $p_1 = p_2 \models_T q_1 = q_2$ when $p_1, p_2 \in \Sigma^*/{T^*_V}$
are $T$-equivalence classes (elements of a free $T$-model).
\end{remark}


\section{Rings and rigs}

We are particularly interested in commutative rings and commutative semirings (rigs) in this work.  Throughout this document, any time we refer to a ring, the reader can assume said ring is commutative.  The same can be assumed for references to rigs.

Rings form a $\Sigma$-algebra with 
\[ \Sigma_{Ring}=\{+^2, \times^2, 0^0, 1^0, -^1\}\]

Rigs form a $\Sigma$-algebra with
\[ \Sigma_{Rig}=\{+^2, \times^2, 0^0, 1^0\} = \Sigma_{Ring}\backslash\{-^1\}\]

We may write $0_R$, $+_R$, etc., to specify that a particular function symbol, or its equivalence class under the ring/rig axioms, comes from the ring/rig $R$.

Rings and rigs each model an equational theory as well.  The set of equations in the equational theory of a ri(n)g correspond exactly to the familiar ri(n)g axioms.

\begin{definition}\label{def:ring}
A \emph{ring} is a $\Sigma_{Ring}$-algebra that models the theory
\begin{align*}
    E_{Ring} =&\ \{x+(y+z) =(x+y)+z, \\
                   &\ \ \ x + y = y + x,\\
                   &\ \ \ x + (-x) = 0, \\
                   &\ \ \ 0 + x = x, \\
                   &\ \ \ x\times(y\times z) = (x \times y) \times z, \\
                   &\ \ \ x \times y = y \times x, \\
                   &\ \ \ 1 \times x = x, \\
                   &\ \ \ x \times (y + z) = (x \times y) + (x \times z)\}
\end{align*}
\end{definition}
That is, for all rings $R$, the ring axioms hold for all elements of $R$. 

\begin{definition}\label{def:rig}
A \emph{rig} is a $\Sigma_{Rig}$-algebra that models the equational theory
\begin{align*}
    E_{Rig} =&\ \{x+(y+z) =(x+y)+z, \\
             &\ \ \ x + y = y + x,\\
             &\ \ \ 0 + x = x, \\
             &\ \ \ x\times(y\times z) = (x \times y) \times z, \\
             &\ \ \ x \times y = y \times x, \\
             &\ \ \ 1 \times x = x, \\
             &\ \ \ x \times (y + z) = (x \times y) + (x \times z)\}
\end{align*}
\end{definition}

That is, the rig axioms hold for all elements of all rigs.

\Needspace{5\baselineskip}
\begin{prop}\label{prop:freering_is_zx}
The following statements hold true:
\begin{enumerate}[nosep] 
\item 
The initial ring is the ring $\ints$ with the usual operations 
of addition and multiplication of integers.
\item 
The initial rig is the rig $\nat$ with the usual operations 
of addition and multiplication of non-negative integers.
\item \label{prop:poly-ring}
The free ring over a finite set $X = \{x_1,\dots,x_n\}$ is the 
 ring of polynomials \[\ints[X] = \ints[x_1,\dots,x_n]\] with integer coefficients and the usual operations.
\item 
The free rig over a finite set $X = \{x_1,\dots,x_n\}$ is the 
 rig of polynomials \[\nat[X] = \nat[x_1,\dots,x_n]\] with non-negative integer coefficients and the usual operations.
\end{enumerate}
\end{prop}

\begin{proof}
We prove only part \ref{prop:poly-ring}; the rest are similar.

Throughout this proof, we refer to the vector of arguments $x_1, ..., x_n$ passed to a polynomial in $\ints[X]$ as $\bar{x}$.

First, note that $\ints[X]$ is indeed a ring: all the ring axioms are satisfied.

Suppose $R$ is a ring, and let $\theta : X \to R$.

For $p(\bar{x}) \in \ints[X]$,  $\semof{p(\bar{x})}_\theta \in R$ is defined by induction on the structure of $p$. The variable base case is given by $\theta$.  Because $R$ is a ring, whenever $p_1(\bar{x}) = p_2(\bar{x})$, $\semof{p_1}_\theta = \semof{p_2}_\theta$. That is, $\semof{-}_\theta$ is well-defined.  We proceed by induction on $p(\bar{x})$ to show that any ring homomorphism $h : \ints[X] \to R$ with $h(x) = \theta(x)$ for all $x \in X$ coincides with $\semof{-}_\theta$.  

\medskip
\Needspace{3\baselineskip}
\noindent
\textbf{Base cases.}

\medskip

\noindent
Let $p(\bar{x}) = x_i$ for some $x_i \in X$. Then by hypothesis
\[
  h(x) = \theta(x) = \semof{x}_\theta.
\]

\medskip

\noindent
Let $p(\bar{x}) = 0$. We know that $h$ must satisfy $h(0) = 0_R$ and 
$\semof{-}_\theta$ must also satisfy $\semof{0}_\theta = 0_R$. Thus,
\[
  h(0) = \semof{0}_\theta.
\]

\medskip

\noindent
Let $p(\bar{x}) = 1$. We know that $h$ must satisfy $h(0) = 1_R$ and 
$\semof{-}_\theta$ must also satisfy $\semof{1}_\theta = 1_R$. Thus,
\[
  h(1) = \semof{1}_\theta.
\]

\medskip

\noindent
\textbf{Inductive cases.}

\medskip

\noindent
Let $p(\bar{x}) = p_1(\bar{x}) + p_2(\bar{x})$. Assume
\[
  h(p_i(\bar{x})) = \semof{p_i(\bar{x})}_\theta.
\]
Then
\begin{align*}
  h(p_1(\bar{x}) + p_2(\bar{x}))
  &= h(p_1(\bar{x})) +_R h(p_2(\bar{x})) \\
  &= \semof{p_1(\bar{x})}_\theta +_R \semof{p_2(\bar{x})}_\theta \\
  &= \semof{p_1(\bar{x}) + p_2(\bar{x})}_\theta.
\end{align*}

\medskip

\noindent
Let $p(\bar{x}) = p_1(\bar{x}) \times p_2(\bar{x})$. Assume
\[
  h(p_i(\bar{x})) = \semof{p_i(\bar{x})}_\theta.
\]
Then
\begin{align*}
  h(p_1(\bar{x}) \times p_2(\bar{x}))
  &= h(p_1(\bar{x})) \times_R h(p_2(\bar{x})) \\
  &= \semof{p_1(\bar{x})}_\theta \times \semof{p_2(\bar{x})}_\theta \\
  &= \semof{p_1(\bar{x}) \times p_2(\bar{x})}_\theta.
\end{align*}

\medskip

\noindent
Let $p(\bar{x}) = -p_1(\bar{x})$. Assume
\[
  h(p_1(\bar{x})) = \semof{p_1(\bar{x})}_\theta.
\]
Then
\begin{align*}
  h(-p_1(\bar{x}))
  &= -_R h(p_1(\bar{x})) \\
  &= -_R \semof{p_1}_\theta \\
  &= \semof{-p_1(\bar{x})}_\theta.
\end{align*}

\medskip

\noindent
Given any environment $\theta$, there exists a unique ring homomorphism 
$\semof{-}_\theta$ from the ring $\ints[X]$ to any other ring $R$. 
$\ints[X]$ has the universal property of the free ring over $X$, therefore, 
it is the free ring over $X$.
\end{proof}

\begin{notation}
Throughout the rest of this document, we write $\models_{Ring}$ to mean $\models_{E_{Ring}}$ and $\models_{Rig}$ to mean $\models_{E_{Rig}}$.  
\end{notation}

\begin{notation}
    For the remainder of this document, if $p \in \ints[X]$ for some set $X = \{x_1, ..., x_n\}$, and $r_1, ..., r_n \in R$ for some ring $R$, then we write $p(r_1, ..., r_n)$ to mean $\semof{p}_\theta$, where $\theta(x_i) = r_i$. Similarly, if $p \in \nat[X]$, and $s_1, ..., s_n \in S$ for some semiring $S$, then we write $p(s_1, ..., s_n)$ to mean $\semof{p}_\theta$, where $\theta(x_i) = s_i$ for the equivalent interpretation function $\semof{-}_\theta$ for semirings.
\end{notation}

\begin{notation}
    If $p \in (\Sigma_{Ring})^*_X$ for some set $X = \{x_1, ..., x_n\}$ and $r_1, ..., r_n \in R$ for some $\Sigma_{Ring}$-algebra $R$, then we write $p(r_1, ..., r_n)$ to mean $\semof{p}_\sigma$ where $\sigma(x_i) = r_i$.  We do similarly for $p \in (\Sigma_{Rig})^*_X$.
\end{notation}

\begin{corollary}[Corollary of Proposition \ref{prop:union_models}]
\label{cor:rigmodels}
The following statements hold true:
    \begin{enumerate}[nosep]
        \item Let $E$ be a $\Sigma_{Ring}$ theory.  Then given $s, t \in (\Sigma_{Ring})^*_X$
            \[ E \models_{Ring} s = t \iff \semof{s}_{\tilde{id_X}}^{\ints[X]} \sim_{[E]} \semof{t}_{\tilde{id_X}}^{\ints[X]} \]
            % \[ E \models_{Ring} s = t \iff [s]_{(E_{Ring})^*_X} \sim_{[E]} [t]_{(E_{Ring})^*_X} \]
        \item Let $E$ be a $\Sigma_{Rig}$ theory.  Then given $s, t \in (\Sigma_{Rig})^*_X$
            \[ E \models_{Rig} s = t \iff \semof{s}_{\tilde{id_X}}^{\nat[X]} \sim_{[E]} \semof{t}_{\tilde{id_X}}^{\nat[X]}\]
    \end{enumerate}
\end{corollary}

\begin{proof}
We prove only part 1. The proof of part 2 is similar.

First, note that by Proposition \ref{prop:freering_is_zx}, $\ints[X]$ is the free $E_{Ring}$ model over the set $X$.

Now, by application of Proposition \ref{prop:union_models} with $\Sigma = \Sigma_{Ring}$, $T = E_{Ring}$, $E = E$, $V = X$, and $s, t = s, t$, the desired result is immediate.
\end{proof}

\begin{notation}
In light of the corollary above, we will often abuse the notation, writing, e.g., $p = q \models_{Ring} s = t$ also when $s, t \in \ints[X]$.
This should always be taken to mean that $p = q \models_{Ring} s' = t'$ for some $s', t' \in (\Sigma_{Ring})^*_X$, where $\semof{s'}^{\ints[X]}_{\tilde id}~=~s$ and $\semof{t'}^{\ints[X]}_{\tilde id}~=~t$.  
The corollary shows that this statement is independent of the choices of $s'$ and $t'$.

We do similarly for $\models_{Rig}$.
\end{notation}


Many of the morphisms induced by the universal properties of the structures introduced in this chapter gain more structure when their domain and codomain share structure.  The following proposition is one such example which will be useful to us.


\begin{prop}\label{prop:rigringhom}
    A rig homomorphism between rings is a ring homomorphism.
\end{prop}
\begin{proof}
    The only additional structure that a rig homomorphism must preserve is additive inverses, i.e., a rig homomorphism $h : A \to B$ must preserve that for all $x \in A$, $h(-x) = -h(x)$.  We can show this as follows: 
    
    We know that a rig homomorphism must preserve addition, so
    \[ h(x) + h(-x) = h(x + (-x)) = h(0_A) = 0_B \]
    Because $B$'s addition forms a group, we can subtract $h(x)$ from both sides of this equation to obtain
    \[h(-x) = -h(x) \qedhere \]
\end{proof}

To conclude this chapter, we make note of a relationship between ideals and congruences that will be useful to us in the coming chapter.

\begin{remark}[Ideals and congruences in a ring]
    The concept of an ideal is fully equivalent to the concept of a congruence on a ring.  Given an ideal $I$ on a ring, we can obtain the corresponding congruence by relating two elements $q_1$ and $q_2$ when $q_1 - q_2 \in I$.  Given a congruence $C$ on a ring, the ideal corresponding to $C$ is the set of elements $x \in R$ such that $[x]_C = [0]_C$.

    In light of this, we will freely transition between congruences and ideals. For example, given a set of equations $E$, we equally treat $\sim_E$ as a congruence, and as the ideal that is induced by this congruence as just stated.  This is consistent with standard notations, for example, the quotient $R/\sim_E$ is the quotient $R/I$ if $I$ is the ideal corresponding to $\sim_E$.
\end{remark}

Now that we have the background to recognize and manipulate algebraic structure, we proceed in the next chapter by defining ADTs, and imposing algebraic structure on them.  Moreover, we use this algebraic structure to prove our main theorem about ADT isomorphisms.


\chapter{Isomorphisms of Algebraic Data Types}\label{sec:ADTs}

\section{Algebraic data types}
Algebraic data types (ADTs) are a class of sum of products data types that may be defined recursively, i.e., may refer to themselves in their definitions.  Functional programming languages like Haskell, OCaml, Agda, and Lean provide support for ADTs by allowing users to define their own.  

In Haskell, a sum of types is encoded as a type with multiple data constructors.  In the example below, every element of the \verb"Value" type is \emph{either} constructed using the \verb"Num" data constructor and contains an integer or is constructed using the \verb"Error" data constructor and contains a string.  The type \verb"Value" is a sum of the types \verb"Int" and \verb"String".
\begin{verbatim}
    data Value = Num Int | Error String
\end{verbatim}

In Haskell, a product of types is encoded as a type with a data constructor that takes multiple arguments.  Shown below is the definition of a type \verb"Point".  Its singular data constructor \verb"Pair" takes two elements of \verb"Float" as arguments and produces an element of \verb"Point".  Every element of \verb"Point", being constructed with the \verb"Pair" data constructor, contains two elements of \verb"Float".  That is, this type is a product of two copies of the \verb"Float" type.
\begin{verbatim}
    data Point = Pair Float Float
\end{verbatim}


In Haskell, data types may refer to themselves in their definitions.  The data type given in the following Haskell type definition--the type of binary trees with unlabeled nodes and unlabeled leaves--has two data constructors \verb"BLeaf" and \verb"BNode".  The \verb"BNode" data constructor takes two elements of \verb"BTree" and produces a new element of \verb"BTree"; it is recursively defined.
\begin{verbatim}
    data BTree = BLeaf | BNode BTree BTree
\end{verbatim}

While Haskell provides convenient examples of algebraic data types, we now introduce a formal grammar for ADTs in order to reason about them abstractly. This separates the mathematical structure of ADTs from any particular programming language and allows us to define their interpretation precisely.
\newcommand{\mb}{\mathbf}
\[ ADT_{\vars} = \vars\ |\ \mathbf{1}\ |\ \mathbf{0}\ |\ ADT_\vars + ADT_\vars\ |\ ADT_\vars \times ADT_\vars\ |\ \mu X.ADT_{\vars, X} \]
The above grammar defines an ADT (in some finite set of variables $\vars$) as either a variable in $\vars$, $\mathbf{1}$, $\mathbf{0}$, a sum of two elements of $ADT_\vars$, a product of two elements of $ADT_\vars$, or a $\mu$ with a variable $X \notin \vars$ and an element of $ADT_{\vars,X}$.  We can view $\mu$ as binding $X$ in an element of $ADT_{\vars,X}$.  Given an element $A$ of $ADT_{\vars,X}$, we denote the syntactic substitution of all occurences of $X$ in $A$ with an element $B$ of $ADT_\vars$ as $A[B/X]$.

\begin{notation}
Throughout the rest of this document, to avoid notational overhead, we write $\mathbf{n}$ to refer to the ADT
\[\underbrace{(\mathbf{1} + \mathbf{1} + \cdots + \mathbf{1})}_{n\text{-times}}\]
We write $A^n$ to refer to the ADT
\[ \underbrace{(A \times A \times \cdots \times A)}_{n\text{-times}} \]
We write $BA$ to refer to the ADT $B \times A$.
\end{notation}

\begin{definition}
    We refer to an ADT in $ADT_{\emptyset}$ constructed using exactly one $\mu$ on the top level as a closed recursive polynomial type.  
\end{definition}
For example, the following ADTs are closed recursive polynomial types:
\begin{itemize}[nosep]
    \item $\mu X.\mathbf{1} + X + X^2$
    \item $\mu Y.\mathbf{3} + \mathbf{7}Y + \mathbf{3}Y^3$
    \item $\mu Z.\mathbf{10}$
\end{itemize}
The following ADTs are not closed recursive polynomial types:
\begin{itemize}[nosep]
    \item $\mu X. \mathbf{5} + X + (\mu Y.\mathbf{3} + Y)$
    \item $\mb{10} + Y + Y^2$
    \item $\mu X.\mb{1} + Y + X^2$
\end{itemize}

We can now define a semantic interpretation of the ADTs generated by the above grammar.  We interpret ADTs in the most general categorical setting that accounts for the features inherent in ADTs: a distributive category $\mcC$ with products, coproducts, an initial object, and a terminal object (that is, a \emph{rig category}) and, importantly, $\mcC$ must be $\omega$-cocomplete, allowing us to interpret recursively defined data types as least fixed points of functors using the $LFP$ functor defined in Chapter \ref{sec:lfp}.  We fix such a category $\mcC$ throughout the remainder of this document.

We interpret ADTs in an environment.  An environment in $\vars$ is a function $\Gamma : \vars \to \mcC$ that assigns an object of $\mcC$ to every variable in $\vars$.  We denote the empty environment, which assigns an object in $\mcC$ to every element in the empty set, as $\Gamma_0$.  We denote updating an environment $\Gamma$ at variable $X$, assigning $X$ the object $Z$ in $\mcC$ by $\Gamma[X := Z]$. Note that if $X \in \vars$, and $\Gamma : \vars \to \mcC$, then $\Gamma[X:=W] : \vars \to \mcC$.  That is, the environment variable assignment function $-[-:=-]$ reassigns a variable $X$ in an environment $\Gamma$ if the domain of $\Gamma$ already includes $X$. Given an environment $\Gamma$, we interpret ADTs according to the following interpretation function $\semof{-}_\Gamma$.

\begin{align*}
    \semof{-}_\Gamma &: ADT_\vars \to \mcC \\
    \semof{v}_\Gamma &= \Gamma\ v && (v \in \vars) \\
    \semof{\mathbf{1}}_\Gamma &= 1 \\
    \semof{\mathbf{0}}_\Gamma &= 0 \\
    \semof{A + B}_\Gamma &= \semof{A}_\Gamma + \semof{B}_\Gamma \\
    \semof{A \times B}_\Gamma &= \semof{A}_\Gamma \times \semof{B}_\Gamma \\
    \semof{\mu X.A}_\Gamma &= LFP(\lambda Y. \semof{A}_{\Gamma[X := Y]})
\end{align*}

A variable $v \in \vars$ is interpreted as the object in $\mcC$ that the environment $\Gamma$ assigns to $v$.  $\mathbf{1}$ is interpreted as $1$, a terminal object in $\mcC$.  $\mathbf{0}$ is interpreted as $0$, an initial object in $\mcC$.  Given two elements $A$ and $B$ of $ADT_\vars$, $A + B$ is interpreted as the categorical coproduct of the intepretations of $A$ and $B$ in $\mcC$.  Given two elements $A$ and $B$ of $ADT_\vars$, $A \times B$ is interpreted as the categorical product of the interpretations of $A$ and $B$ in $\mcC$.  We use the notations for $\times$ and $+$ to mean both to the binary constructors in the syntax of ADTs as well as categorical products and coproducts.  It will always be clear from context which ones we are referring to. Given an element $A$ of $ADT_{\vars,X}$, $\mu X.A$ is interpreted as the least fixed point of the functor that takes an object $Y$ in $\mcC$ and gives back the interpretation of $A$ in $\mcC$ under the environment $\Gamma[X := Y]$.

\subsection*{Functoriality of interpretation}
\begin{remark}
We also note that interpretation of ADTs is functorial in the environment $\Gamma$.

That is, the interpretation function $\semof{A}_{(-)}$ gives rise to a functor 
$\mcC^\vars \to \mcC$, where $\Gamma : \vars \to \mcC$ is viewed as an object in the category $\mcC^\vars$
and a morphism $\gamma$ from $\Gamma \to \Gamma'$ is a function mapping $v \in \vars$
to a morphism in $\mcC$, that is, $\gamma(v) : \Gamma(v) \to \Gamma'(v)$.

The functoriality of interpretation is immediate if one considers that it is defined by functorial operations on a category.
The variable case refers to the morphism $\gamma$. The $\mu$ case follows from Theorem \ref{theorem:lfpfunctor}.
\end{remark}


\section{Rig isomorphisms}\label{sec:rigiso}

We can impose the structure of a semiring on the set $ADT_\vars$ by considering elements as equal when there exists a canonical isomorphism between their interpretations.  We call these isomorphisms, which are described in this section, \emph{rig isomorphisms}. The multiplication operator in this semiring is $\times$.  Its identity is $\mathbf{1}$.  The addition operator in this semiring is $+$. Its identity is $\mathbf{0}$.  As defined in Definition \ref{def:rig}, the semiring axioms are
\begin{align*}
    E_{Rig} =&\ \{x+(y+z) =(x+y)+z, \\
             &\ \ \ x + y = y + x,\\
             &\ \ \ 0 + x = x, \\
             &\ \ \ x\times(y\times z) = (x \times y) \times z, \\
             &\ \ \ x \times y = y \times x, \\
             &\ \ \ 1 \times x = x, \\
             &\ \ \ x \times (y + z) = (x \times y) + (x \times z)\}
\end{align*}

If we view $+$ and $\times$ above as $\mcC$'s categorical coproduct and product, respectively, then we can clearly see that each of these axioms is reflected as an isomorphism in $\mcC$.  For example, we have that coproducts are associative up to natural isomorphism, and thus the first rig axiom is satisfied.  Similarly, coproducts are commutative up to natural isomorphism.  The last listed rig axiom in $E_{Rig}$ states that $\times$ distributes over $+$.  Because $\mcC$ is a distributive category, this axiom is also satisfied by a natural isomorphism in $\mcC$.  Each of these rig axioms is satisfied by natural isomorphisms in $\mcC$.  For more details about these natural isomorphisms, see ``Reason Isomorphically!'' \cite{10.1145/1863495.1863507} by Ralf Hinze and Daniel James.

For the natural isomorphisms above to give the structure of a rig to ADTs, we also must know that isomorphism is a congruence with respect to the functions that define our semiring.  First, isomorphism is an equivalence relation. It is reflexive, since every object has an identity isomorphism; symmetric, since every isomorphism has an inverse; and transitive, since the composition of isomorphisms is again an isomorphism.  Next, since isomorphism is a congruence with respect to products and coproducts (the interpretations of the syntax $\times$ and $+$, respectively, for ADTs), equality in our semiring is, in fact, a congruence.  That is, given isomorphisms $i_1 : A_1 \to A_2$ and $i_2 : B_1 \to B_2$, because functors preserve isomorphisms, the functorial action of the coproduct functor $+$ on $i_1$ and $i_2$, that is, $i_1 + i_2 : A_1 + B_1 \to A_2 + B_2$, is also an isomorphism.  A corresponding rule applies for products.  Also, note that $\mu$ is interpreted as the least fixed point ($LFP$) of a functor, and Chapter \ref{sec:lfp} proved that $LFP$ is a functor. Because functors preserve isomorphisms, equality in this semiring is a congruence with respect to $\mu$.

\begin{figure}
\fbox{
  \begin{minipage}{0.9\linewidth}
    \begin{mathpar}
      \inferrule{A \in ADT_\vars}{A \simeq A} \qquad
    \inferrule{A \simeq B}{B \simeq A} \qquad
    \inferrule{A \simeq B \qquad B \simeq C}{A \simeq C} \\
    \inferrule{A_1 \simeq A_2 \qquad B_1 \simeq B_2}{A_1 + A_2 \simeq B_1 + B_2} \qquad
    \inferrule{A_1 \simeq A_2 \qquad B_1 \simeq B_2}{A_1 \times A_2 \simeq B_1 \times B_2} \qquad
    \inferrule{A \simeq B}{\mu X.A \simeq \mu X.B} \\
    \inferrule{A, B, C \in ADT_\vars}{A + (B + C) \simeq (A + B) + C} \qquad
    \inferrule{A, B \in ADT_\vars}{A + B \simeq B + A} \qquad
    \inferrule{A \in ADT_\vars}{\mathbf{0} + A \simeq A} \\
    \inferrule{A, B, C \in ADT_\vars}{A \times (B \times C) \simeq (A \times B) \times C} \qquad
    \inferrule{A, B \in ADT_\vars}{A \times B \simeq B \times A} \qquad
    \inferrule{A \in ADT_\vars}{\mathbf{1} \times A \simeq A} \\
    \inferrule{A, B, C \in ADT_\vars}{A \times (B \times C) \simeq (A \times B) + (A \times C)} \\
    \inferrule{A \in ADT_{\vars,X}}{\mu X.A \simeq A[\mu X.A/X]}\\
    \end{mathpar}
  \end{minipage}
}
\caption{The inference rules that generate derivations of isomorphisms in the category $\mcC$ in which we interpret ADTs}\label{fig:adt_isos}
\end{figure}
\newcommand{\ADTrig}{\mathcal{A}_\vars}

\begin{definition}
    We define the semiring $\ADTrig$ whose carrier set is $ADT_\vars/\simeq$, where $\simeq$ is the congruence defined in Figure \ref{fig:adt_isos}.  The multiplication and addition of $\ADTrig$ are given by the syntactic operators $\times$ and $+$, respectively.  The multiplicative and additive identities of $\ADTrig$ are $\mb{1}$ and $\mb{0}$, respectively.
\end{definition}

For the remainder of this chapter, we lift the syntax $\simeq$ for isomorphism in $\mcC$ to our syntax of ADTs.  That is, given two elements of $ADT_{\vars}$, we say
\[ A \simeq B \]
to mean that $A$ is related to $B$ under the congruence $\simeq$.  That is to say, $A = B$ in $\ADTrig$.  This notation is logical because every equality in $\ADTrig$ is reflected as a natural isomorphism in $\mcC$.

\begin{notation}
    Consider a closed recursive polynomial type $\mu X.p(X)$, where $p \in (\Sigma_{Rig})^*_{\{X\}}$ is evaluated in the $\Sigma_{Rig}$-algebra formed by the syntax of ADTs. This data type is interpreted as the least fixed point ($LFP$) of the functor
    \[ \lambda Y. \semof{p(X)}_{\Gamma_0[X := Y]} \]
    Because $p$ is in $(\Sigma_{Rig})^*_{\{X\}}$, the above functor maps an object $Y$ precisely to $p$ evaluated at $Y$ in the rig category $\mcC$ considered with its $\Sigma_{Rig}$-algebra structure given by $0$, $1$, $+$, and $\times$.  The above functor can thereby be rewritten as
    \[ \lambda Y.p(Y)\]
    or, simply, $p$.
\end{notation}

The interpretation of $\mu$ introduces some additional equalities into our rig.  As an example, consider the ADT
\[ M = \mu X.\mb{1} + X + X^2\]
Under the empty environment, this ADT is interpreted as the least fixed point of a certain endofunctor.  Specifically, let $m \in (\Sigma_{Rig})^*_{\{X\}}$ be the element defined as
\[ m = 1 + X + X^2\]

As above, $m$ can be viewed as the endofunctor defined by
\begin{align*}
    m &: \mcC \to \mcC \\
    m(Y) &= 1 + Y + Y^2
\end{align*}

When $m$ is viewed in this way, the interpretation $\mathcal{M}$ of $M$ is
\[ \mathcal{M} = LFP(\lambda Y. \semof{\mathbf{1} + X + X^2}_{\Gamma_0[X := Y]}) = LFP(\lambda Y.1 + Y + Y^2) = LFP(m)\]
As discussed in Chapter \ref{sec:lfp}, the least fixed point of a functor is the object part of its initial algebra.  In that chapter, we also proved that the initial algebra of a functor is an isomorphism; see Corollary \ref{cor:inexistsandisiso}.  That is, there exists an isomorphism $in_m$ between $\mathcal{M}$ and $m(\mathcal{M})$.  We therefore obtain a new equality in $\ADTrig$ that does not follow from the rig axioms alone.  The isomorphism
\[ \mathcal{M} \simeq m(\mathcal{M}) \]
gives the equality
\[ M \simeq m(M)\]
in $\ADTrig$.  We call this isomorphism, $M$'s fixed point isomorphism.  This behavior is captured more generally by the following definition.

\begin{definition}
Let $A \in ADT_{\vars,X}$. Then, by the reasoning of Section \ref{sec:initial_algebras_morphisms}, for any environment $\Gamma$ in $\vars$, there exists an isomorphism
\[ \semof{\mu X.A}_\Gamma \simeq \semof{A[\mu X.A/X]}_\Gamma \]
called the \emph{fixed point isomorphism} of $A$.
\end{definition}

If $A$ is a closed recursive polynomial type, that is,
\[ A = \mu X.p(X) \]
for some $p \in (\Sigma_{Rig})^*_{\{X\}}$, then the above specializes to
\[ \semof{\mu X.p(X)}_{\Gamma_0} \simeq p(\semof{\mu X.p(X)}_{\Gamma_0}) \]

The inference rules shown in Figure \ref{fig:adt_isos} generate derivations of isomorphisms in the semiring $\ADTrig$.  The first row expresses that isomorphism is an equivalence relation. The second row states that isomorphism is preserved by the semiring operations and by the formation of recursive types. The third, fourth, and fifth rows state that the semiring axioms hold for ADTs when isomorphism is considered as equality.  The last inference rule states that a recursively defined ADT is isomorphic to its body with the recursive variable substituted with the recursive type itself.

\section{Algebraic data types as complex numbers}

Andreas Blass, in his 1994 paper ``Seven Trees in One'', observes that viewing ADTs modulo the semiring axioms and fixed point isomorphisms for recursive types yields isomorphisms with interesting properties.  For one, these isomorphisms are generic in the sense that they do not depend on the underlying structure of the category $\mcC$ in which we interpret ADTs, so long as it has sufficient structure to interpret ADTs in the first place. That is, these isomorphisms exist not just when ADTs are interpreted in $\Set$ as their set of elements.  Moreover, the isomorphisms described in ``Seven Trees in One'' appear to arise from purely algebraic manipulations of polynomials.  Blass's core discovery was that if we treat recursive type definitions as polynomial equations, and the recursive type itself as a complex root of such an equation, then we can lift resulting algebraic equalities to isomorphisms in $\mcC$.

For example, consider the element $m$ of $(\Sigma_{Rig})^*_{\{X\}}$ defined as
\[ m = 1 + X + X^2 \]
The data type defined by
\[ M = \mu X.m(X) \]
has a fixed point isomorphism given in Figure \ref{fig:adt_isos} as
\[ M \simeq m(M) = \mb{1} + M + M^2\]
If we treat the above isomorphism as a numerical equality in $\mbC$, then the equation $M = 1 + M + M^2$ simplifies to $0 = 1 + M^2$.  Solving for $M$, we find that $M = \pm i$.  While we cannot give any true interpretation to a data type which is ``equal to'' a complex number, this result does predict some isomorphisms in $\mcC$.  First, note that from $M = \pm i$, we can derive that $M^4 = 1$.  A simple interpretation of $M$ in $\Set$ shows that, $M^4 \simeq 1$ cannot be true in a general categorical setting, as $M^4$ is interpreted as a set of 4-tuples where each component is a recursively defined set with infinite cardinality.  However, if we multiply the whole equation by $M$, we get $M=M^5$.  And $M\simeq M^5$ can indeed be formally derived in $\mcC$ by the following chain of isomorphisms, where each step applies $M$'s fixed point isomorphism exactly once, and isomorphisms corresponding to the rig axioms are left implicit.

\begin{align*}
    M &\simeq \mb{1} + M + M^2 \\
    &\simeq \mb{1} + \mb{2}M + M^2 + M^3 \\
    &\simeq \mb{1} + \mb{2}M + \mb{2}M^2 + M^3 + M^4 \\
    &\simeq \mb{1} + \mb{2}M + \mb{3}M^2 + M^3 + \mb{2}M^4 \\
    &\simeq \mb{1} + \mb{2}M + \mb{3}M^2 + \mb{2}M^3 + \mb{2}M^4 + M^5 \\
    &\simeq \mb{1} + \mb{2}M + \mb{3}M^2 + \mb{3}M^3 + \mb{2}M^4 + \mb{2}M^5 \\
    &\simeq \mb{1} + \mb{2}M + \mb{3}M^2 + \mb{3}M^3 + \mb{3}M^4 + \mb{2}M^5 + M^6 \\
    &\simeq \mb{2}M + \mb{2}M^2 + \mb{3}M^3 + \mb{3}M^4 + \mb{2}M^5 + M^6 \\
    &\simeq M + \mb{2}M^2 + \mb{2}M^3 + \mb{3}M^4 + \mb{2}M^5 + M^6 \\
    &\simeq \mb{2}M^2 + M^3 + \mb{3}M^4 + \mb{2}M^5 + M^6 \\
    &\simeq M^2 + M^3 + \mb{2}M^4 + \mb{2}M^5 + M^6 \\
    &\simeq M^3 + M^4 + \mb{2}M^5 + M^6 \\
    &\simeq M^4 + M^5 + M^6 \\
    &\simeq M^5
\end{align*}

Blass recognized the connection between algebraic identities and type isomorphisms.  However, he does not prove any conditions that such algebraic manipulations of types considered as complex numbers give rise to generic categorical isomorphisms.  This was not done until 2002, when Fiore and Leinster published their paper ``Objects of Categories as Complex Numbers'' \cite{Fiore2005}.  The next subsection gives an account of their proof.


% \section{Rig Isomorphisms}

% ADTs form a semiring or \textit{rig}  when isomorphism is considered as equality.  In this rig, categorical product is the multiplication operation and categorical coproduct is the addition operation.  Equality in this rig is determined by the following isomorphisms.

% \begin{description}
% \item[Groupoid operations]
% Equality is always an equivalence relation.
% \begin{itemize}[nosep]
%     \item Reflexivity: $\forall (a : ADT\ \mathbb{V})\ .\ a \simeq a$
%     \item Symmetry: $\forall (a\ b : ADT\ \mathbb{V})\ .\ a \simeq b \implies b \simeq a$
%     \item Transitivity: $\forall (a\ b\ c : ADT\ \mathbb{V})\ .\ a \simeq b\ \land\ b \simeq c \implies a \simeq c$
% \end{itemize}
% \item[Congruence laws]
% Equality is always a congruence with respect to the operations of the structure; in the case of ADTs, these are products, coproducts, and least fixed points. 
% \begin{itemize}[nosep]
%     \item Congruence of Products: $\forall (a1\ a2\ b1\ b2 : ADT\ \mathbb{V})\ .\ a1 \simeq a2\ \land\ b1 \simeq b2 \implies a1 \times b1 \simeq a2 \times b2$
%     \item Congruence of Coproducts: $\forall (a1\ a2\ b1\ b2 : ADT\ \mathbb{V})\ .\ a1 \simeq a2\ \land\ b1 \simeq b2 \implies a1 + b1 \simeq a2 + b2$
%     \item Congruence of Least Fixed Points:
%     $\forall (a\ b : ADT\ \mathbb{V})\ .\ a \simeq b \implies \mu a \simeq \mu b$
% \end{itemize}

% \item[Commutative rig axioms]
% ADTs form a commutative rig.
% \begin{itemize}[nosep]
%     \item Associativity of Products: $\forall (a\ b\ c : ADT\ \mathbb{V})\ .\ a \times (b \times c) \simeq (a \times b) \times c$
%     \item Associativity of Coproducts: $\forall (a\ b\ c : ADT\ \mathbb{V})\ .\ a + (b + c) \simeq (a + b) + c$
%     \item Commutativity of Products: $\forall (a\ b : ADT\ \mathbb{V})\ .\ a \times b \simeq b \times a$
%     \item Commutativity of Coproducts: $\forall (a\ b : ADT\ \mathbb{V})\ .\ a + b \simeq b + a$
%     \item Identity of Products: $\forall (a : ADT\ \mathbb{V})\ .\ a \simeq \mathbf{1} \times a$
%     \item Identity of Coproducts: $\forall (a : ADT\ \mathbb{V})\ .\ a \simeq \mathbf{0} \times a$
%     \item Left Distributivity of Products Over Coproducts: $\forall (a\ b\ c : ADT\ \mathbb{V})\ .\ a \times (b + c) \simeq (a \times b) + (a \times c)$
%     \item Left Distributivity of Products Over Coproducts: $\forall (a\ b\ c : ADT\ \mathbb{V})\ .\ (a + b) \times c \simeq (a \times c) + (b \times c)$
%     \item The $\mathbf{0}$ Annihilator Property: $\forall (a : ADT\ \mathbb{V})\ .\ a \times \mathbf{0} \simeq \mathbf{0}$
% \end{itemize}
% \item[Fixed point rule]
% Every ADT formed by applying the $\mu$ operator to an underlying ADT $a$ is a fixed point of $a$'s interpretation as a functor.

% \begin{itemize}[nosep]
%     \item The Fixed Point Property of LFPs: $\forall (a : ADT\ (\mathbb{V} \cup \{X\}))\ .\ \mu X.a \simeq (a [X := \mu X.a])$ where $a [X :=\mu X.a]$ represents replacing all occurrences of $X$ in $a$ with $\mu X. a$. 
% \end{itemize}
% \end{description}

% The property of the rig that we have just defined that separates it from most rigs is the least fixed point operator.  This operator comes equipped with its own congruence rule as well as an additional isomorphism rule (that gives rise to additional equalities in this rig), the fixed point property of least fixed points.  We may want to ask the question: what new equalities do these rules introduce into our rig?  For an arbitrary ADT $a$, when is $\mu a$ isomorphic to $P(\mu a)$, where $P$ is a polynomial functor?  More generally, when is $Q_1(\mu a)$ isomorphic to $Q_2(\mu a)$? 

% Fiore and Leinster have answered these questions in the paper ``Objects of Categories as Complex Numbers'' \cite{Fiore2005}.  In the following subsections, we will give a brief summary of their results.  Additional details can be found in \cite{Fiore2005}.

\subsection{An outline of Fiore and Leinster's prior work}

We know that $\mu X.P(X) \simeq P(\mu X. P(X))$ for all polynomials $P \in \nat[X]$ from the reasoning provided in Section \ref{sec:rigiso}.  One question that might be natural to ask is what other isomorphisms does this isomorphism induce?  A particular instance of this question asks, ``if $X$ is an ADT, and $P$ is a polynomial with natural number coefficients, and we have an isomorphism $X \simeq P(X)$, when do we get an isomorphism between $Q_1(X)$ and $Q_2(X)$ for some polynomials $Q_1$ and $Q_2$ with natural number coefficients?''  Fiore and Leinster solved this problem in their 2002 paper ``Objects of Categories as Complex Numbers'' \cite{Fiore2005}.  We proceed by describing their solution.

Fiore and Leinster use the following terminology, which is necessary to understand their results.

\begin{definition}\label{def:ringtheoretic}
Let $p_1, p_2, q_1, q_2 \in \ints[x]$.  Then $p_1(x) = p_2(x) \implies q_1(x) = q_2(x)$ holds \textit{ring theoretically} when $p_1 = p_2 \models_{Ring} q_1 = q_2$.
\end{definition}

\begin{prop}
    The following are equivalent:
    \begin{enumerate}[label=(\alph*),nosep]
    \item \label{pl:a} $p_1 = p_2 \models_{Ring} q_1 = q_2$
    \item \label{pl:b} $q_1$ and $q_2$ represent the same element of the quotient ring $\ints[x]/(p_1-p_2)$.
    \item \label{pl:c} $p_1 - p_2$ divides $q_1 - q_2$ in $\ints[x]$.
    \end{enumerate}
\end{prop}
\begin{proof} \hfill
\begin{itemize}[nosep]
\item \ref{pl:a} $\iff$ \ref{pl:b}: by Corollary \ref{cor:rigmodels}.
\item \ref{pl:b} $\implies$ \ref{pl:c}:
     Assume (b).  Then $q_1 = q_2$ in $\ints[x]/(p_1 - p_2)$, so $q_1 - q_2 \in (p_1 - p_2)$.  Because the ideal $(p_1 - p2)$ is the principal ideal generated by $p_1 - p_2$, every element in it has the form $k \times (p_1 - p_2)$. $q_1 - q_2$ therefore has the form $k \times (q_1 - q_2)$ for some polynomial $k$, so (b) implies (c).
\item \ref{pl:c} $\implies$ \ref{pl:b}:
     Assume (c). We have that $p_1 - p_2$ divides $q_1 - q_2$, so $q_1 - q_2 = k \times (p_1 - p_2)$, and $q_1 - q_2$ is in the ideal generated by $p_1 - p_2$. If $q_1 - q_2$ is in the ideal $(p_1 - p_2)$, then $q_1 = q_2$ in $\ints[x]/(p_1 - p_2)$ and (c) implies (b).
\end{itemize}
\end{proof}

\begin{definition}\label{def:rigtheoretic}
Let $p_1, p_2, q_1, q_2 \in \nat[x]$.  Then $p_1(x) = p_2(x) \implies q_1(x) = q_2(x)$ holds \textit{rig theoretically} when $p_1 = p_2 \models_{Rig} q_1 = q_2$:
\end{definition}

\begin{prop}
    The following are equivalent:
    \begin{enumerate}[label=(\alph*),nosep]
    \item \label{pl2:a} $p_1 = p_2 \models_{Rig} q_1 = q_2$
    \item \label{pl2:b} $q_1$ and $q_2$ represent the same element of the quotient rig $\nat[x]/(p_1=p_2)$.
    \item \label{pl2:c} There exists some finite sequence of transformations $r_0, ... r_n$ such that $r_1 = q_1$, $r_n = q_2$, and for all $i \in \nat$ from $1$ to $n$, there exists some $f \in \nat[x]$ and $k \in \nat$ such that \[ r_{i-1} = f(x) + x^k\times p_1(x) \]
    and
    \[ r_i = f(x) + x^k \times p_2(x) \]
\end{enumerate}
\end{prop}

\begin{proof} \hfill
    \begin{itemize}[nosep]
\item \ref{pl2:a} $\iff$ \ref{pl2:b}: by Corollary \ref{cor:rigmodels}.
\item \ref{pl2:b} $\iff$ \ref{pl2:c}: The equivalence of \ref{pl2:b} and \ref{pl2:c} will follow from the fact that these two conditions define the same congruence.

   Specifically, let $(p_1 = p_2)$ denote the least congruence relating $p_1$ and $p_2$. Then \ref{pl2:b} exactly asserts that $(q_1,q_2) \in (p_1=p_2)$.

   Let $\mathcal{R}$ denote the congruence defined in \ref{pl2:c}. We will show that $\mathcal{R} \subseteq (p_1=p_2)$; since $\mathcal{R}$ is a congruence that relates $p_1$ and $p_2$, and $(p_1 = p_2)$ is the least congruence with this property, $\mathcal{R} = (p_1 = p_2)$.
   
    To show that $\mathcal{R}$ is a subcongruence of $(p_1 = p_2)$, it is sufficient to show that every possible step from $r_i$ to $r_{i+1}$ is related by the congruence $(p_1 = p_2)$.  Because $(p_1 = p_2)$ is a congruence, $(f(x), f(x)) \in (p1 = p2)$ for all $f$.  Similarly, $(x^k, x^k)$ for all $k$. Because $p_1 = p_2$ generates the congruence $(p_1 = p_2)$, we know $(p_1, p_2) \in (p_1 = p_2)$. Therefore, the following holds for all $f \in \nat[x]$ and $k \in \nat$:
    \[ (f(x) + x^k\times p_1(x) , f(x) + x^k \times p_2(x)) \in (p_1 = p_2) \]
    Therefore, every relation in $\mathcal{R}$ is also related by $(p_1 = p_2)$, so $\mathcal{R}$ is a subcongruence of $(p_1 = p_2)$.  Because $(p_1 = p_2)$ is the least congruence relating $p_1$ and $p_2$, any subcongruence of it that relates $p_1$ and $p_2$ is equal to it, so $\mathcal{R} = (p_1 = p_2)$. $(q_1, q_2) \in (p_1 = p_2)$ exactly means that $q_1$ is equal to $q_2$ in $\nat[x]/(p_1 = p_2)$.  \ref{pl2:b} and \ref{pl2:c} are therefore logically equivalent. \qedhere    
\end{itemize}
\end{proof}

Let $p$, $q_1$, and $q_2$ be polynomials in $\nat[x]$. Fiore and Leinster proved that if $x = p(x) \implies q_1(x) = q_2(x)$ ring theoretically (i.e., when $p$, $q_1$, and $q_2$ are considered as polynomials in $\mathbb{Z}[x]$), then the same holds true rig theoretically (with the implication holding in $\nat[x]$) under some conditions.
Sufficient conditions are that $p$ has a non-zero constant term and degree at least two and that $q_1$ and $q_2$ have degree at least 1.  Note that this result is proved in the algebraic context of rigs and not in a categorical context, but nevertheless holds true in any rig, including in $\ADTrig$, in which equalities are directly interpreted as isomorphisms in $\mcC$.
Moreover, Fiore and Leinster note that the result does generalize to the categorical setting.

For $x = p(x) \implies q_1(x) = q_2(x)$ to be true ring theoretically, $q_1(x)$ and $q_2(x)$ must represent the same element of the quotient ring $\mathbb{Z}[x]/(p(x) - x)$.  Therefore, $p(x) - x$ must divide $q_1(x) - q_2(x)$ in the ring $\mathbb{Z}[x]$, i.e., there exists some $f(x)$ such that $q_1(x) - q_2(x) = f(x) \times (p(x) - x)$.  If we rearrange this equation to avoid negative coefficients, we get the equation 
\[q_1(x) + x\times f(x) = q_2(x) + p(x)\times f(x)\]
The rig $\nat[x]$ trivially injects into $\ints[x]$, i.e., is a subset of $\ints[x]$. Because any injective morphism reflects equalities, the above equation must hold true rig theoretically in addition to holding ring theoretically.  We know that we are working in a rig where $x = p(x)$, so the equation above yields 
\[ q_1(x) + x\times f(x) = q_2(x) + x\times f(x)\]
In general, in a rig, if $x + z = y + z$, we cannot conclude that $x = y$, but Fiore and Leinster give a condition under which we can. 
Whenever this condition is met, we are able to conclude that $q_1(x) = q_2(x)$ rig theoretically.

We therefore turn our attention to Fiore and Leinster's work on cancellation in commutative semigroups \cite{Fiore2005}.

\subsection{The high elements of a semigroup}\label{subsubsection:highelements}

Consider a commutative semigroup $(A, *)$.  Let the relation $\le_A$ be defined as follows:

\[ b \le_A a \iff \exists\ c \in A\ \mathsf{such\ that}\ b * c = a \]
That is, $b$ is related to $a$ under $\le_A$ if $b$ divides $a$.  Note that the relation $\le_A$ is transitive. Now, consider the set of elements of $A$ to which every other element is related under $\le_A$, i.e., the set of all $x \in A$ such that for all $y \in A$, $y \le_A x$.  We call these elements \textit{high} and we denote the set of high elements of $(A, *)$ as $H(A)$.

\begin{theorem}\label{theorem:cancellation}
    If $a$ and $b$ are high elements of a commutative semigroup $(A, *)$, and $x$ is any element of $A$, then $a * x = b * x \implies a = b$. 
\end{theorem}

We prove the above theorem using the following intermediate lemmas.

\begin{lemma}\label{lemma:ha_closure}
    Let $(A, *)$ be a commutative semigroup. $H(A)$ is closed under multiplication by elements of $A$.
\end{lemma}

\begin{proof}
    Let $a \in H(A)$ and $b \in A$.  We know that for all $x \in A$, $x \le_A a$, i.e., there exists some $y \in A$ such that $x * y = a$.  Then $x \le_A a * b$, because $x * y * b = a * b$.
\end{proof}

In particular, $H(A)$ is a subsemigroup of $A$.
\begin{lemma}\label{lemma:clique}
    Let $(A, *)$ be a commutative semigroup.  Then for all $a, b \in H(A)$, $b \le_{H(A)} a$, i.e., there exists some $x \in H(A)$ such that $b * x = a$.
\end{lemma}

\begin{proof}
    Because $a \in H(A)$, we know that for all $b \in H(A)$, $b * a \le_A a$, i.e., there exists some $d \in A$ such that $b * a * d = a$.  By Lemma \ref{lemma:ha_closure}, we know that $a * d \in H(A)$.  Let $x = a * d$.  Then $b * x = a$ and  $b \le_{H(A)} a$.
\end{proof}

\begin{lemma}\label{lemma:HAabeliangroup}
    For any commutative semigroup $(A, *)$, if $H(A)$ is non-empty, then $(H(A), *)$ is an abelian group.
\end{lemma}

\begin{proof}
    Associativity and commutativity are inherited from the commutative semigroup $(A, *)$.  The unit of the abelian group $(H(A), *)$ is given as follows: for any $x \in H(A)$, $x \le_{H(A)} x$ (by Lemma \ref{lemma:clique}), which means there exists a $z \in H(A)$ such that $x * z = x$.  Given any other element $y \in H(A)$, $x \le_{H(A)} y$ (again, by Lemma \ref{lemma:clique}), so there exists a $w \in H(A)$ such that $x * w = y$.  Now
    \[ y * z = x * w * z = w * x * z = w * x = y\]
    % \[ x * z = x \implies w * x * z = w * x \implies y * z = y \]
    \noindent
    and $z$ is the unit of the abelian group.  Because $z \in H(A)$, by Lemma \ref{lemma:clique}, for all $v \in H(A)$, $v \le_{H(A)} z$, which implies the existence of inverses.  $(H(A), *, z)$ is, therefore, an abelian group.
\end{proof}

\begin{corollary}\label{cor:uniquez}
    Let $(A, *)$ be a commutative semigroup.  Given any $x \in H(A)$, any element $y \in H(A)$ witnessing $x \le_{H(A)} x$ is the identity of the abelian group $(H(A), *)$.
    
    That is, for any $y,y' \in H(A)$, $x * y = x = x * y'$ implies $y = y'$.
\end{corollary}

\begin{proof}
    Lemma \ref{lemma:HAabeliangroup} proves that $(H(A), *)$ is an abelian group and the element $z$ witnessing $x \le_{H(A)} x$ for an arbitrary $x$ is an identity for $*$.  Because the identity of an abelian group is unique, any element $y$ witnessing that $x \le_{H(A)} x$ is the identity of $(H(A), *)$.
\end{proof}

\begin{corollary}\label{cor:uniqueinverses}
    Let $(A, *)$ be a commutative semigroup.  Let $z$ be the identity element of the abelian group $(H(A), *, z)$.  Any element $w \in H(A)$ witnessing $x \le_{H(A)} z$ is the unique inverse of $x$.
\end{corollary}

\begin{proof}
    Let $x * w = z$ and $x * w' = z$.  Then 
    \[ x * w = z = x * w' \]
    Multiplying this equation on the left by $x^{-1}$ (by Lemma \ref{lemma:HAabeliangroup}), we get $w = w'$.
\end{proof}

We can now use the fact that $(H(A), *)$ forms an abelian group to prove our main theorem about cancellation in commutative semigroups.

\begin{proof}[Proof of Theorem \ref{theorem:cancellation}]
    Let $(A, *)$ be a commutative semigroup.  Let $a, b \in H(A)$, and $x \in A$.  We have that $a * x = b * x$.  We first post-multiply this equation by $z$ (as described in Lemma \ref{lemma:HAabeliangroup}) to get $a * x * z = b * x * z$.  We now use the fact that $(H(A), *)$ is a group to divide this equation through by the high element $x * z$ to obtain $a = b$.  
\end{proof}

Next, we summarize Fiore and Leinster's classification of the high polynomials in our semiring of interest.

\subsection{Classification of high polynomials}\label{subsubsection:highpolynomials}

Consider the least congruence generated by the equation $x = p(x)$ over the rig of polynomials with natural number coefficients.  We can quotient the rig $\nat[x]$ by this congruence to get the rig $\nat[x]/(x=p(x))$, where equality is determined by the usual rig congruence laws, but also any additional equalities generated by the equation $(x = p(x))$.  Fiore and Leinster show that if $p$ is a polynomial with degree at least 2 and a non-zero constant term, then all non-constant polynomials are high in the commutative semigroup $(\nat[x]/(x=p(x)), +)$ underlying the quotient rig $\nat[x]/(x=p(x))$.  Therefore, if $p$ is a polynomial with degree at least 2 and a non-zero constant term,
\[q_1(x) + x\times f(x) = q_2(x) + x\times f(x) \implies q_1(x) = q_2(x)\]
if $q_1$ and $q_2$ are high (non-constant) polynomials.

For the rest of this section, we write $f \le g$ to mean $f \le_{\nat[x]/(x=p(x))}g$.  If every polynomial in our rig is related to a polynomial $f$, then $f$ is a high polynomial.  Here, we replicate Fiore and Leinster's result that all non-constant polynomials in the rig $\nat[x]/(x=p(x))$ are high if $p$ has a non-zero constant term and degree at least 2.  First, we must prove several related lemmas about the $\le$ relation.

\begin{lemma}\label{lemma:le_mul}
    If $g_1 \le f_1$ and $g_2 \le f_2$ then $g_1g_2 \le f_1f_2$.
\end{lemma}
\begin{proof}
    By definition of $\le$, there exists $h_1$ such that $g_1 + h_1 = f_1$.  Similarly, there exists $h_2$ such that $g_2+h_2 = f_2$.

    \[ f_1f_2 = (g_1+h_1)\times (g_2+h_2) = g_1g_2 + g_1h_2 + h_1g_2 + h_1h_2 \]

    $g_1g_2 \le f_1f_2$ because adding $g_1h_2 + h_1g_2 + h_1h_2$ to $g_1g_2$ results in $f_1f_2$.
\end{proof}
\Needspace{6\baselineskip}
\begin{lemma}\label{lemma:incdeg}
    If $p$ has a non-zero constant term, then $1 \le x \le x^2 \le ...$.
\end{lemma}

\begin{proof}
    By hypothesis, there exists some $p'(x)$ such that $p(x) = p'(x) + 1$,  which means $1 \le p(x)$, so $1 \le x$.
    
    Also $x^n + 0 = x^n$, so $x^n \le x^n$.
    
    By Lemma \ref{lemma:le_mul}, $x^n \le x^{n+1}$.
\end{proof}

\begin{lemma}\label{lemma:kxx}
    If $p$ has a non-zero constant term and degree at least two, then for all $k$, $kx \le x$.
\end{lemma}

\begin{proof}
    By hypothesis, we have $1 + x^d \le p(x) = x$ for some $d \ge 2$.  We therefore have,

    % \[ x \ge 1 + x^d \ge x^d \ge x^{d-1}x = x^{d-1}(p(x)) \ge x^{d-1}(1+x^d) \ge x^{d-1} + x^{2d-1} \ge x + x \ge 2x\]
    \begin{align*}
        2x &= x+x & \\
        &\le x^{d-1} + x^{2d-1} & \textsf{(By Lemma \ref{lemma:incdeg})} \\
        &= x^{d-1}(1+x^d) & \\
        &\le x^{d-1}(p(x)) & \textsf{(By Lemma \ref{lemma:le_mul})} \\
        &= x^{d-1}x & \\
        &= x^d  & \\
        &\le x^d + 1 & \\
        &\le x
    \end{align*}

    By induction on $k$, for all $k$, $kx \le x$.
\end{proof}
\Needspace{5\baselineskip}
\begin{lemma}\label{lemma:xnx}
    If $p$ has a non-zero constant term and degree at least two, then for all $n \ge 1$, $x^n \le x$.
\end{lemma}

\begin{proof}
    By hypothesis, we have $1 + x^d \le x$ for some $d \ge 2$.  We therefore have, by Lemma \ref{lemma:incdeg},
    \[ x^2 \le x^d \le x^d + 1 \le x\]
    Multiplying through by $x^{k-1}$, we get $x^{k+1} \le x^k$ for all $k \ge 1$.  In particular,
    \[x^n \le x^{n-1} \le \cdots \le x\]
    Because $\le$ is transitive, $x^n \le x$.
\end{proof}

Fiore and Leinster's result on high polynomials is reproduced below.

\begin{theorem}\label{theorem:singlevarhighpolynomials}
    If $p$ has a non-zero constant term and degree at least two, then every non-constant polynomial in $\nat[x]/(x=p(x))$ is high.
\end{theorem}

\begin{proof}
    By Lemmas \ref{lemma:kxx} and \ref{lemma:xnx},

    $x \ge kx \ge x^{n_1} + x^{n_2} + x^{n_3} + \dots + x^{n_k}$

    That is, $x$ is high.  Consider a non-constant $n$th degree polynomial $f$ in $\nat[x]/(x=p(x))$.  Let $f = x^n + f'(x)$. By Lemma \ref{lemma:incdeg},

    \[ f = x^n + f'(x) \ge x^n \ge x\]

    Therefore, $f$ is a high polynomial.
\end{proof}

\subsection{The main theorem for the univariate case}

The previous section proves that given any polynomial $p$ with a non-zero constant term and degree at least 2, every non-constant polynomial in $\nat[x]/(x = p(x))$ is high.  We can now prove Fiore and Leinster's main two theorems.

\begin{theorem}\label{theorem:ringtorig}
    Let $p, q_1, q_2 \in \nat[x]$.  Let $p$ have a non-zero constant term and degree at least 2.  Let $q_1$ and $q_2$ be non-constant polynomials.  Then, if
    \[ x = p(x) \implies q_1(x) = q_2(x)\]
    holds ring theoretically, the implication also holds rig theoretically. 
\end{theorem}
\begin{proof}
    By hypothesis, $q_1$ and $q_2$ represent the same element of the quotient ring $\ints[x]/(p(x)~-~x)$, so $q_1 - q_2 \in (p(x) - x)$ and $p(x) - x$ divides $q_1 - q_2$.  We therefore have
    \[q_1(x) - q_2(x) = f(x) \times (p(x) - x) \]
    which implies
    \[ q_1(x) + x \times f(x) = q_2(x) + p(x) \times f(x)\]
    and because $x = p(x)$
    \[ q_1(x) + x \times f(x) = q_2(x) + x \times f(x) \]
    The above equality holds rig theoretically in addition to ring theoretically because the inclusion of $\nat[x]$ into $\ints[x]$, being injective, reflects equality.

    $q_1$ and $q_2$ are non-constant polynomials, so by \ref{theorem:singlevarhighpolynomials}, they are high.  By Theorem \ref{theorem:cancellation}, we can cancel $x \times f(x)$ from both sides of the above equation to get that $q_1(x) = q_2(x)$ in $\nat[x]/(x = p(x))$.
\end{proof}

Fiore and Leinster also give a condition under which \[ x = p(x) \implies q_1(x) = q_2(x)\] holds true both ring theoretically and rig theoretically.

\begin{theorem}\label{theorem:complextoring}
    Let $p, q_1, q_2 \in \nat[x]$ be as stated in the above theorem.  Additionally, suppose that $p(x) - x$ is primitive and has no repeated roots.  Then, if 
    \[ a = p(a) \implies q_1(a) = q_2(a) \]
    holds for all complex roots $a$ of $p(x) - x$, then
    \[ x = p(x) \implies q_1(x) = q_2(x) \]
    holds true rig theoretically.
\end{theorem}
\begin{proof}
    By the polynomial division algorithm, 
    there exists some $f, g \in \mathbb{Q}[x]$ such that 
    \[ q_1 - q_2 = f \times (p - x) + g\]
    and $deg(g) < deg(p - x)$.  Each root of $p - x$ is a root of $q_1 - q_2$, so is also a root of $g$.  Because $p - x$ has no repeated roots, $g$ must have as many unique roots as $p - x$, but has degree strictly less than $p - x$.  By the fundamental theorem of algebra, $g = 0$.  If $f$ is $0$, then $f \in \ints[x]$.  Otherwise, we can write $f$ as $\tilde{f} / k$ where $\tilde{f}$ is a primitive polynomial in $\ints[x]$, and $k$ is an integer.  We can now derive
    \[ k\times (q_1 - q_2) = \tilde{f} \times (p - x)\]
    Gauss's Lemma asserts that the product of two primitive polynomials is primitive.  Since, $\tilde{f}$ is primitive, and $p - x$ is primitive by hypothesis, $k \times (q_1 - q_2)$ is primitive, meaning $k = \pm 1$ and $f \in \ints[x]$.
    Therefore, $q_1 = q_2$ in $\ints[x]/(p(x) - x)$, i.e., $x = p(x) \implies q_1(x) = q_2(x)$ ring theoretically. By Theorem \ref{theorem:ringtorig}, the same implication holds rig theoretically.
\end{proof}

\subsection{An example derivation of isomorphism from equality}
We can use the above theorems to derive isomorphisms between algebraic data types. Consider the type $M = \mu X. m(X)$ where $m = x^2 + x + 1$ is an element of $(\Sigma_{Rig})^*_{\{x\}}$.  Then we know that $M = m(M)$, because $M$ is a (the least) fixed point of $m$.  

Let $q_1 = x^3 + x^2 + x + 1$ and $q_2  = 3x^2 + 3$ be elements of $(\Sigma_{Rig})^*_{\{x\}}$.  The roots of the polynomial $q_1(x) - q_2(x) = x^3 + x^2 + x + 1 - (3x^2 + 3)$ are $\{i, -i, 2\}$, and the roots of $m(x) - x$ are $\{i, -i\}$.  $m$ has a non-zero constant term, degree at least 2.  $m(x) - x$ has no repeated roots and is primitive.  $q_1$ and $q_2$ are non-constant polynomials.  The roots of $q_1 - q_2$ are a superset of the roots of $m(x) - x$.  Considering all of the above, Theorem \ref{theorem:complextoring} applies, and we can generate an isomorphism between $q_1(M)$ and $q_2(M)$.

First, we divide $q_1(x) - q_2(x)$ by $m(x) - x$ to get a quotient of $x - 2$. That is,
\[ x^3 + x^2 + x + 1 - (3x^2 + 3) = (x - 2)(x^2 + x + 1 - x) \]
Now, recall that
\[ q_1 - q_2 = f(p - x) \implies q_1+fx = q_2 + fp\]
In this case, $f(x) = x - 2$ is a polynomial containing a negative coefficient, so let $f(x) = f_1(x) - f_2(x)$ where $f_1(x) = x$ and $f_2(x) = 2$ each have only positive coefficients to get
\[ q_1 - q_2 = (f_1-f_2)(m - x) \]
and
\[ q_1 + f_1x + f_2m = q_2 + f_1m + f_2x \]
Because this polynomial identity holds in $\nat[x]$, it is derivable using only rig axioms.  

So we can construct an isomorphism 
\begin{align*}
    i &: q_1(M) + f_1(M)\times M + f_2(M)\times m(M)  \\
    &\hspace{1.5cm}\simeq q_2(M) + f_1(M) \times m(M) + f_2(M)\times M
\end{align*}
using only compositions of basic isomorphisms that correspond to the commutative rig axioms.  We now use the isomorphisms corresponding to the rig congruence laws together with the isomorphism corresponding to the fixed point property of the data type $M$ to obtain the following two isomorphisms $j_1$ and $j_2$:
\begin{align*}
    j_1 &: q_1(M) + f_1(M)\times M + f_2(M)\times m(M) \\
    &\hspace{1.5cm}\simeq q_1(M) + f_1(M)\times M + f_2(M)\times M
\end{align*}
and
\begin{align*}
    j_2 &: q_2(M) + f_1(M) \times m(M) + f_2(M)\times M \\
    &\hspace{1.5cm}\simeq q_2(M) + f_1(M) \times M + f_2(M)\times M
\end{align*}
We can use $i$, $j_1$, and $j_2$ together with the transitivity to get the isomorphism
\begin{align*}
    k &: q_1(M) + f_1(M)\times M + f_2(M)\times M \\
    &\hspace{1.5cm}\simeq q_2(M) + f_1(M) \times M + f_2(M)\times M
\end{align*}
Lastly, we use Theorem \ref{theorem:cancellation} to ``cancel'' $r=f_1(M) \times M + f_2(M)\times M = M^2 + 2\times M$ from both sides of our isomorphism.  Proceeding as in the proof of that theorem, we do this by first adding the identity element $z$ of the abelian group of high elements of $\nat[M]$, which is guaranteed to exist by Lemma \ref{lemma:HAabeliangroup}, then by adding the high inverse of $r+z$, which is also guaranteed to exist by Lemma \ref{lemma:HAabeliangroup} together with Theorem \ref{theorem:singlevarhighpolynomials}.  Clearly $M^2 + 1$ witnesses the fact that $M \le M$ so, by Corollary \ref{cor:uniquez}, $M^2 + 1$ is the additive identity among the high elements of $\nat[M]$.  $M^2 + 1$ has an identity isomorphism, $id_{M^2+1}$.  We can use the fact that isomorphism is a congruence with respect to addition to obtain the isomorphism $k'$ as $k + id_{M^2 + 1}$.
\begin{align*}
    k' &: q_1(M) + f_1(M)\times M + f_2(M)\times M + M^2 + 1 \\
    &\hspace{1.5cm}\simeq q_2(M) + f_1(M) \times M + f_2(M)\times M + M^2 + 1
\end{align*}

The additive inverse of $r + z = M \times M + 2\times M + (M^2 + 1)$ among the high elements of $\nat[M]$ is derivable constructively using Lemma \ref{lemma:HAabeliangroup} together with Theorem \ref{theorem:singlevarhighpolynomials}. In fact, $2M^3+1$ is such an element; it witnesses that $r + z \le z$.  This is demonstrated by a chain of isomorphisms generated by repeated application of $M$'s fixed point isomorphism.  The isomorphism
$h : r+z-(r+z) \simeq z$
% \begin{align*}
% h &: r+z-(r+z) \simeq z
% % h &: M^2 + 2M + M^2 + 1 + 2M^3 + 1 \simeq 1 + M^2
% \end{align*}
is defined by the following chain of isomorphisms:
\begin{align*}
 % h : r+z-(r+z) &\simeq z\\
     r + z + 2M^3 + 1
    &= M^2 + 2M + M^2 + 1 + 2M^3 + 1 \\
    &\simeq 2M^3 + 2M^2 + 2M + 2 \\
    &\simeq 2M^3 + M^2 + 2M + 1 \\
    &\simeq M^3 + M^2 + M + 1 \\
    &\simeq M^2 + 1 \\
    &\simeq 1 + M^2
\end{align*}
By Corollary \ref{cor:uniqueinverses}, $2M^3 + 1$ is the additive inverse of $r + z$ in the high elements of $\nat[M]$.  

There exists an identity isomorphism $id_{-(r+z)}$ for the additive inverse of $r + z$, that is, for $2M^3 + 1$. Again, using the fact that isomorphism is a congruence with respect to addition together with the isomorphism $k$, we obtain the following isomorphism $l$ from $k' + id_{-(r+z)}$
\begin{align*}
    l &: q_1(M) + f_1(M)\times M + f_2(M)\times M + (M^2 + 1) + (2M^3 + 1) \\
    &\simeq q_2(M) + f_1(M) \times M + f_2(M)\times M + (M^2 + 1) + (2M^3 + 1)
\end{align*}
Unfolding the definitions of $q_1$, $q_2$, $f_1$, and $f_2$ lets us rewrite the above as:
\begin{align*}
    l &: (M^3 + M^2 + M + 1) + M^2 + 2M + (M^2 + 1) + (2M^3 + 1) \\
    &\hspace{1.5cm}\simeq (3M^3+3) + M^2 + 2M + (M^2 + 1) + (2M^3 + 1)
\end{align*}
Applying $h$ twice, once for each side of $l$ results in the isomorphism $l'$:
\begin{align*}
    l' &: (M^3 + M^2 + M + 1) + M^2 + 1 \\
    &\hspace{1.5cm}\simeq (3M^3+3) + M^2 + 1
\end{align*}

Lastly, we must ``cancel'' $z$ from both sides of $l'$.  Again, the sequences of isomorphisms that do this are derivable constructively from Lemma \ref{lemma:HAabeliangroup} together with Theorem \ref{theorem:singlevarhighpolynomials}. We shall skip this step and merely note that $t_1$ and $t_2$
satisfy the requirement, where
\begin{align*}
t_1 &: (M^3 + M^2 + M + 1) + M^2 + 1 \simeq M^3 + M^2  + M + 1 \\
t_2 &: (3M^3+3) + M^2 + 1 \simeq 3M^3 + 3
\end{align*}
Explicitly, $t_1$ is defined by the following chain of isomorphisms:
\begin{align*}
    (M^3 + M^2 + M + 1) + M^2 + 1
    &\simeq M^3 + 2M^2 + M + 2 \\
    &\simeq M^3 + M^2  + M + 1
\end{align*}
and $t_2$ is defined by the following chain of isomorphisms:
\begin{align*}
    (3M^3+3) + M^2 + 1
    &\simeq 3M^3 + M^2 + 4 \\
    &\simeq M^4 + 3M^3 + 2M^2 + 4 \\
    &\simeq M^4 + 4M^3 + 2M^2 + M + 4 \\
    &\simeq M^4 + 4M^3 + M^2 + M + 3 \\
    &\simeq M^4 + 3M^3 + M^2 + 3 \\
    &\simeq 3M^3 + 3
\end{align*}
Composing $t_1$, $t_2$, and $l'$, we obtain the desired isomorphism
\[ M^3 + M^2 + M + 1 \simeq 3M^3 + 3\]

% We now interpret the quotient $x - 2$ as an isomorphism between data types.  First, recall that isomorphism is a congruence with respect to products and coproducts.  Our base isomorphism is
% \[ i : M \simeq M^2 + M + 1\]
% We interpret $(M + (-2))i$ as an isomorphism by applying the isomorphism congruence laws for products and coproducts, which interpret addition and multiplication of polynomials, together with the isomorphism symmetry law, which interprets negative polynomials.
% \begin{align*}
%     (M-2)i &= Mi + (-2i) \\
%     Mi &: M\times M \simeq M \times (M^2 + M + 1) \\
%     2i &: 2\times M \simeq 2 \times (M^2 + M + 1) \\
%     -2i &: 2 \times (M^2 + M + 1) \simeq 2 \times M \\
%     (M - 2)i &: M \times M + 2 \times (M^2 + M + 1) \simeq M \times (M^2 + M + 1) + 2 \times M
% \end{align*}
% By using the commutative rig axioms for ADT isomorphisms, we can simplify each side of the isomorphism $(M-2)i$ to the following isomorphism $((M-2)i)'$:
% \[ ((M-2)i)' : 3M^2 + 2M + 1 \simeq M^3 + M^2 + 3M \]

This example illustrates how the preceding results can be used to explicitly construct isomorphisms between algebraic data types. We now turn to the general problem of determining when such isomorphisms exist, extending the work of Fiore and Leinster to the multivariate setting.

\section{Generalization to the multivariate case}

\newcommand{\ering}{{\mathcal{E}_{\mathrm{CRing}}}}

Fiore and Leinster solved the question of when two polynomials are equal in a rig of polynomials containing one variable, and with additional equalities specified by a single equation of the form $x = p(x)$, where $p$ is some polynomial over $x$ with a non-zero constant term and degree at least 2.  We now wish to extend their work to rigs of polynomials containing an arbitrary number of variables, with one equation specifying additional equalities for each variable, i.e., we want to decide equality in a rig that has the form $\nat[x_1, x_2, ... , x_n]/(x_1=p_1(x_1),..., x_n=p_n(x_n))$, hereafter referred to as the rig $\nat[\bar{x}]/(\bar{x}=\bar{p}(\bar{x}))$.  In order to do this, we need two additional theorems.  We need to know that high polynomials in such a rig form a ring, and we need to know which polynomials in this rig are high.

Fiore and Leinster note that the high polynomials of a rig form a ring \cite{Fiore2005}.  To prove this, we make use of the following lemmas which together prove the zero annihilation property holds of our ring of high polynomials.

\begin{lemma}\label{lemma:zz=z}
    Consider the additive identity $z$ of a rig's underlying additive semigroup as described in Lemma \ref{lemma:HAabeliangroup}. The following equality holds:
    \[ z = zz \]
\end{lemma}
\begin{proof}
    First, note that 
    \[ (1 + z) + (-(1 + z)) = z \]
    That is, by Lemma \ref{lemma:HAabeliangroup}, $(1 + z)$ has an additive inverse and together they add to the additive identity $z$. If we multiply the above equation by $z$, we get
    \[ (z + zz) + z(-(1 + z)) = z + (zz + z(-1(1 + z))) = zz\]
    We conclude from this equation that $z \le zz$ with $(zz + z(-1(1 + z)))$ as the witness.  Because $\le$ is transitive and $z$ is high, $zz$ is a high element in any rig containing additively high elements.

    Now consider the equation
    \[ z + zz = zz = z(z + z) = zz + zz \]
    Because $zz$ is high, we can add its inverse to both sides of this equation to obtain that
    \[ z = zz  \qedhere \]
\end{proof}

\begin{lemma}\label{lemma:az=z}
    For all additively high elements $a$ of a rig, $az = z$.
\end{lemma}
\begin{proof}
    Consider the following equation which holds by Lemma \ref{lemma:zz=z}.
    \[ az + az = a(z + z) = az = (a + z)z = az + zz = az + z\]
    By Lemma \ref{lemma:ha_closure}, the high elements are closed under addition by all elements of a rig's underlying additive semigroup.  Because, by the above equation, $az$ can be expressed as the high element $z$ added to another element $az$, $az$ is high.

    Because $az$ is high, by Lemma \ref{lemma:HAabeliangroup} we can add its additive inverse to both sides of the above equation to obtain that $az = z$.
\end{proof}

\begin{theorem}\label{thm:highring}
    High polynomials in the rig $\nat[\bar{x}]/(\bar{x}=\bar{p}(\bar{x}))$ together with their inherited binary operations form a ring.
\end{theorem}

\begin{proof}
    We have seen that the high elements of $\nat[\bar{x}]/(\bar{x}=\bar{p}(\bar{x}))$ form an abelian group with $z \in H(\nat[\bar{x}]/(\bar{x}=\bar{p}(\bar{x})))$ as its additive identity as described in Lemma \ref{lemma:HAabeliangroup}. This ring's multiplication and its related distributive property are inherited from $\nat[\bar{x}]/(\bar{x}=\bar{p}(\bar{x}))$.  It remains to show that there exists a multiplicative identity.
    \Needspace{6\baselineskip}
    We can show that the multiplicative unit of the ring we are defining is $z + 1$ as follows:
    Take any high polynomial $p$. By Lemma \ref{lemma:az=z}, the following equation holds:
    \[ p(z + 1) = pz + 1p = pz + p = z + p \]
    By Lemma \ref{lemma:HAabeliangroup}, we know for all high $a$, $a + z = a$, therefore $z + p = p$, and we have that $z + 1$ is a multiplicative unit for the high polynomials.
\end{proof}

Now that we have that the high polynomials of $\nat[\bar{x}]/(\bar{x}=\bar{p}(\bar{x}))$ form a ring, it is useful to determine what these high polynomials are.

First, note that the relation $\le$ respects both addition and, as proved in Lemma \ref{lemma:le_mul}, multiplication.

We proceed by proving the following lemma which is used to prove our main theorem about high polynomials in $\nat[\bar{x}]/(\bar{x}=\bar{p}(\bar{x}))$.

\begin{lemma}\label{lemma:mvmono}
    In the rig $\nat[\bar{x}]/(\bar{x}=\bar{p}(\bar{x}))$, if for all $i$, $p_i$ has a non-zero constant term and degree at least 2, then for any monomial $f = x_1^{k_1}x_2^{k_2}...x_n^{k_n}$ (including those where for one or more $i$, $k_i = 0$) with a leading coefficient of 1, $f(\bar{x}) \le x_1x_2...x_n$.
\end{lemma}

\begin{proof}
    By Lemma \ref{lemma:incdeg} together with Lemma \ref{lemma:xnx}, for each $i$, for all $k$, $x_i^k \le x_i$.

    By Lemma \ref{lemma:le_mul}, for any monomial $f$ with a leading coefficient of 1, $f \le x_1x_2...x_n$.
\end{proof}

\begin{theorem}
    In the rig $\nat[\bar{x}]/(\bar{x}=\bar{p}(\bar{x}))$, if for all $i$, $p_i$ contains a non-zero constant term and has degree at least 2, then every polynomial containing at least one monomial with every variable is high.
\end{theorem}

\begin{proof}
    By Lemma \ref{lemma:le_mul} together with Lemma \ref{lemma:kxx}, we have

    \[ x_1x_2...x_n \ge kx_1x_2...x_n = x_1x_2...x_n + x_1x_2...x_n + ... + x_1x_2...x_n \]

    The above relation and equation, when considered together with Lemma \ref{lemma:mvmono}, proves that $x_1x_2...x_n$ is high.

    Consider a polynomial $f$ containing at least one monomial with every variable.  Let $f = x_1^{k_1}x_2^{k_2}...x_n^{k_n} + f'(x)$ where all $k_i \ge 1$. Then by Lemmas \ref{lemma:le_mul} and \ref{lemma:incdeg}, 
    \[ f = x_1^{k_1}x_2^{k_2}...x_n^{k_n} + f'(x) \ge x_1^{k_1}x_2^{k_2}...x_n^{k_n} \ge x_1x_2...x_n \]
    \noindent
    Thus, $f$ is high.
\end{proof}

The following subsubsection proves that the ring of high polynomials in $\nat[\bar{x}]/(\bar{x}=\bar{p}(\bar{x}))$ is isomorphic to the ring $\ints[\bar{x}]/(p_1(x_1)-x_1, p_2(x_2)-x_2, ..., p_n(x_n)-x_n)$, which will henceforth be referred to as the ring $\ints[\bar{x}]/(\bar{p}(\bar{x})-\bar{x})$.

Next, we establish that the ring of high polynomials is, in fact, the ring of integer polynomials in the same variables.

\subsection{The ring of high polynomials}

\begin{notation}
For this subsubsection, we shall adopt the following notation:
\begin{itemize}[nosep]
\item Let $I$ be the ideal in $\ints[\bar{x}]$ generated by the polynomials $p_1(x_1)-x_1,..., p_n(x_n)-x_n$, where each $p_i$ has a nonzero constant term and degree at least 2.

\item Let $C$ be the least congruence in $\nat[\bar{x}]$ generated by the equations $x_1=p_1(x_1), ..., x_n=p_n(x_n)$.  

\item Let $H_C(\nat[\bar{x}])$ consist of the polynomials in $\nat[\bar{x}]$ whose equivalence classes are high in the ring $\nat[\bar{x}]/C$.  

\item Let $[s]_C$ be the equivalence class of the polynomial $s$ in the rig $H_C(\nat[\bar{x}])/C$.  

\item Let $[s]_I$ be the equivalence class of the polynomial $s$ in the ring $\ints[\bar{x}]/I$.
\end{itemize}
\end{notation}

Consider the following rig homomorphism $h^+$:
\begin{align*}
    h^+ &: H_C(\nat[\bar{x}]) \to \ints[\bar{x}]/I \\
    h^+(r(\bar{x})) &=  [r(\bar{x})]_I
\end{align*}

Observe that the congruence in $\ints[\bar{x}]$ generated by $I$ equals the congruence in $\ints[\bar{x}]$ generated by $C$, since $p_i(x_i) - x_i = 0 \implies p_i(x_i) = x_i$. Thus, if $q_1, q_2 \in \ints[\bar{x}]$, then
\[ (q_1(\bar{x}), q_2(\bar{x})) \in C \implies (q_1(\bar{x}), q_2(\bar{x})) \in I \]
and $q_1 = q_2$ in $\ints[x]/I$.

By the universal property of quotients of $\Sigma$-algebras, there exists a unique rig homomorphism $\tilde{h}^+ : H_C(\nat[\bar{x}])/C \to \ints[\bar{x}]/I$ such that 
\[\forall r \in H_C(\nat[\bar{x}]).\ \tilde{h}^+([r]_C) = h^+(r)\]
Because the domain and codomain of $\tilde{h}^+$ are each rings, by Proposition \ref{prop:rigringhom}, $\tilde{h}^+$ is a ring homomorphism.

To get a morphism from $\ints[\bar{x}]/I \to H_C(\nat[\bar{x}])/C$, first recall that $\ints[\bar{x}]$ is the free $E_{Ring}$-model over the set $\bar{x}$.  Next, note that by Theorem \ref{thm:highring}, $H_C(\nat[\bar{x}])/C$ satisfies the equational theory of a ring, i.e., is a ring.  

Let $z_C$ be the additive identity of the ring $H_C(\nat[\bar{x}])/C$ as described in Theorem \ref{thm:highring}. Let $\theta(x_i) = x_i + z_C$.  By the universal property of the free ring $\ints[\bar{x}]$, there exists a unique ring homomorphism $h^-$ satisfying the following:

\begin{align*}
    h^- &: \ints[\bar{x}] \to H_C(\nat[\bar{x}])/C \\
    h^-(x_i) &=\theta(x_i) = [x_i + z_C]_C && (x_i \in \bar{x}) \\
    h^-(f_1(\bar{x})+f_2(\bar{x})) &= h^-(f_1(\bar{x})) + h^-(f_2(\bar{x})) \\
    h^-(f_1(\bar{x})\times f_2(\bar{x})) &= h^-(f_1(\bar{x})) \times h^-(f_2(\bar{x})) \\
    h^-(-f(\bar{x})) &= -h^-(f(\bar{x})) \\
    h^-(0) &= [z_C]_C \\
    h^-(1) &= [1 + z_C]_C \\
\end{align*}

The action of the ring homomorphism $h^-$ simplifies using the following lemma when it acts on polynomials with only positive integer coefficients, i.e., on elements of $\nat[\bar{x}]$.

\begin{lemma}\label{lemma:h-z_C}
    If $f(\bar{x}) \in \nat[\bar{x}]$, then $h^-(f(\bar{x})) = [f(\bar{x}) + z_C]_C$.
\end{lemma}
\begin{proof}
    We proceed by induction on $f$.

    The following are our base cases:

    \begin{itemize}[nosep]
        \item Let $f(\bar{x}) = x_i$.  Then by definition, $h^-(f(\bar{x})) = [x_i + z_C]_C = [f(\bar{x}) + z_C]_C$.
        \item Let $f(\bar{x}) = 0$. Then $h^-(f(\bar{x})) = [z_C]_C = [0 + z_C]_C = [f(\bar{x}) + z_C]_C$.
        \item Let $f(\bar{x}) = 1$. Then by definition, $h^-(f(\bar{x})) = [1 + z_C]_C = [f(\bar{x}) + z_C]_C$.
    \end{itemize}

    The following are the inductive cases.  Assume that $h^-(f_1(\bar{x})) = [f_1(\bar{x}) + z_C]_C$ and $h^-(f_2(\bar{x})) = [f_2(\bar{x}) + z_C]_C$.  

    \begin{itemize}[nosep]
        \item Let $f(\bar{x}) = f_1(\bar{x}) + f_2(\bar{x})$.  Then 
        \begin{align*}
            h^-(f(\bar{x})) &= h^-(f_1(\bar{x})) + h^-(f_2(\bar{x})) \\
            &= [f_1(\bar{x}) + z_C]_C + [f_2(\bar{x}) + z_C]_C \\
            &= [f_1(\bar{x}) + z_C + f_2(\bar{x}) + z_C]_C \\
            &= [f_1(\bar{x}) + f_2(\bar{x}) + z_C + z_C]_C \\
            &= [f_1(\bar{x}) + f_2(\bar{x}) + z_C]_C + [z_C]_C \\
            &= [f_1(\bar{x}) + f_2(\bar{x}) + z_C]_C
        \end{align*}
        \item Let $f(\bar{x}) = f_1(\bar{x}) \times f_2(\bar{x})$. Then
        \begin{align*}
            h^-(f(\bar{x})) &= h^-(f_1(\bar{x})) \times h^-(f_2(\bar{x})) \\
            &= [f_1(\bar{x}) + z_C]_C \times [f_2(\bar{x}) + z_C]_C \\
            &= [(f_1(\bar{x}) + z_C) \times (f_2(\bar{x}) + z_C)]_C \\
            &= [f_1(\bar{x})\times f_2(\bar{x}) + f_1(\bar{x})\times z_C + f_2(\bar{x}) \times z_C + z_C \times z_C]_C \\
            &= [f_1(\bar{x})\times f_2(\bar{x}) + z_C]_C + [f_1(\bar{x})]_C \times [z_C]_C + [f_2(\bar{x})]_C \times [z_C]_C + [z_C]_C \times [z_C]_C \\
            &= [f_1(\bar{x})\times f_2(\bar{x}) + z_C]_C + [z_C]_C + [z_C]_C + [z_C]_C \;\;\;\;\;\;\;\;\;\;\;\;\;\;\;\;(\textsf{by Theorem \ref{thm:highring}}) \\
            &= [f_1(\bar{x})\times f_2(\bar{x}) + z_C]_C
        \end{align*}
        \item The case where $f(\bar{x}) = -f_1(\bar{x})$ need not be considered because by assumption, $f$ has only positive integer coefficients.
    \end{itemize}
\end{proof}

To obtain a ring homomorphism from the quotient $\ints[\bar{x}]/I$ to $H_C(\nat[\bar{x}])/C$, we first must prove the following lemma:

\begin{lemma}
    For all $(s, t) \in I$, $h^-(s) = h^-(t)$.
\end{lemma}
\begin{proof}
    It suffices to show that for all $i$, $h^-(p_i(x_i) - x_i) = [z_C]_C$.  This is shown as follows:

    \begin{align*}
        h^-(p_i(x_i) - x_i) &= h^-(p_i(x_i)) - h^-(x_i) \\
        &= [p_i(x_i) + z_C]_C - [x_i + z_C]_C && \textsf{(by Lemma \ref{lemma:h-z_C})} \\
        &= [x_i + z_C]_C - [x_i + z_C]_C && \textsf{(Because $x_i=p_i(x_i)$ is a generator of C)} \\
        &= [z_C]_C
    \end{align*}
\end{proof}

Now we have by the universal property of the quotient $\ints[\bar{x}]/I$, that there exists a unique ring homomorphism $\tilde{h}^- : \ints[\bar{x}]/I \to H_C(\nat[\bar{x}])/C$ such that
\[\tilde{h}^-([f(\bar{x})]_C) = h^-(f(\bar{x}))\]

\begin{theorem}
    $\tilde{h}^+$ and $\tilde{h}^-$ form a ring isomorphism between the rings $H_C(\nat[\bar{x}])/C$ and $\ints[\bar{x}]/I$.
\end{theorem}
\begin{proof}
    First, we show that $\tilde{h}^- \circ \tilde{h}^+ = id_{H_C(\nat[\bar{x}])/C}$ as follows:

    Let $f(\bar{x}) \in H_C(\nat[\bar{x}])$.  Then

    \begin{align*}
        \tilde{h}^- (\tilde{h}^+([f(\bar{x})]_C)) &= \tilde{h}^-(h^+(f(\bar{x}))) \\
        &= \tilde{h}^-([f(\bar{x})]_I) \\
        &= h^-(f(\bar{x})) \\
        &= [f(\bar{x}) + z_C]_C && (\textsf{By Lemma \ref{lemma:h-z_C}}) \\
        &= [f(\bar{x})]_C + [z_C]_C && (\textsf{Because $f(\bar{x}) \in H_C(\nat[\bar{x}])$}) \\
        &= [f(\bar{x})]_C \\
        &= id_{H_C(\nat[\bar{x}])/C}([f(\bar{x})]_C)
    \end{align*}

    We now show that $\tilde{h}^+ \circ \tilde{h}^- = id_{\ints[\bar{x}]/I}$ as follows: let $f(\bar{x}) \in \ints[\bar{x}]$. Without loss of generality, we can say that $f(\bar{x}) = f_1(\bar{x}) - f_2(\bar{x})$ where $f_1$ and $f_2$ each have only positive integer coefficients.  Then we have

    \begin{align*}
        \tilde{h}^+ (\tilde{h}^-([f(\bar{x})]_I)) &= \tilde{h}^+ (\tilde{h}^-([f_1(\bar{x}) - f_2(\bar{x})]_I)) \\
        &= \tilde{h}^+ (h^-(f_1(\bar{x}) - f_2(\bar{x}))) \\
        &= \tilde{h}^+(h^-(f_1(\bar{x})) - h^-(f_2(\bar{x}))) \\
        &= \tilde{h}^+([f_1(\bar{x}) + z_C]_C - [f_2(\bar{x}) + z_C]_C) && (\textsf{By Lemma \ref{lemma:h-z_C}}) \\
        &= \tilde{h}^+([f_1(\bar{x}) + z_C]_C) - \tilde{h}^+([f_2(\bar{x}) + z_C]_C) && (\textsf{Because $\tilde{h}^+$ is a ring homomorphism}) \\
        &= h^+(f_1(\bar{x}) + z_C) - h^+(f_2(\bar{x}) + z_C) \\
        &= [f_1(\bar{x}) + z_C]_I - [f_2(\bar{x}) + z_C]_I \\
        &= [f_1(\bar{x}) + z_C - (f_2(\bar{x}) + z_C)]_I \\
        &= [f_1(\bar{x}) - f_2(\bar{x})]_I \\
        &= [f(\bar{x})]_I \\
        &= id_{\ints[\bar{x}]/I}([f(\bar{x})]_I)
    \end{align*}

    and $\tilde{h}^+$ and $\tilde{h}^-$ are mutual inverses.  Indeed, $\tilde{h}^+$ and $\tilde{h}^-$ are isomorphisms.
\end{proof}

We can now show that equalities in $\ints[\bar{x}]/I$ are directly reflected as equalities in $H_C(\nat[\bar{x}])/C$.

\begin{corollary}\label{cor:equalityZtoN}
    Let $q_1(\bar{x}), q_2(\bar{x}) \in H_C(\nat[\bar{x}])$.  If $[q_1(\bar{x})]_I = [q_2(\bar{x})]_I$ in the ring $\ints[\bar{x}]/I$, then $[q_1(\bar{x})]_C = [q_2(\bar{x})]_C$ in the ring $H_C(\nat[\bar{x}])/C$.
\end{corollary}
\begin{proof}
    By assumption, we have that 
    \[ [q_1(\bar{x})]_I = [q_2(\bar{x})]_I \]
    By Definition of $\tilde{h}^+$, 
    \[ \tilde{h}^+([q_1(\bar{x})]_C) = \tilde{h}^+([q_2(\bar{x})]_C)\]
    $\tilde{h}^+$ is an isomorphism, so it is injective.  From the above, we arrive at
    \[ [q_1(\bar{x})]_C = [q_2(\bar{x})]_C \]
    in the domain of $\tilde{h}^+$, $H_C(\nat[\bar{x}])/C$.
\end{proof}

The above corollary also extends to equality in $\nat[\bar{x}]/C$, because $H_C(\nat[\bar{x}])$ is a subset of $\nat[\bar{x}]$.

\begin{corollary}\label{cor:equalityHNtoN}
    Let $q_1(\bar{x}), q_2(\bar{x}) \in H_C(\nat[\bar{x}])$.  If $[q_1(\bar{x})]_I = [q_2(\bar{x})]_I$ in the ring $\ints[\bar{x}]/I$, then $[q_1(\bar{x})]_C = [q_2(\bar{x})]_C$ in the rig $\nat[\bar{x}]/C$.
\end{corollary}
\begin{proof}
    By Corollary \ref{cor:equalityZtoN}, $[q_1(\bar{x})]_C  = [q_2(\bar{x})]_C$ in the ring $H_C(\nat[\bar{x}])/C$.  Because every ring is a rig, and $\nat[\bar{x}]/C$ is a subrig of $H_C(\nat[\bar{x}])/C$, there exists an inclusion homomorphism from $\nat[\bar{x}]/C$ to $H_C(\nat[\bar{x}])/C$.  By the injectivity of this homomorphism, $[q_1(\bar{x})]_C  = [q_2(\bar{x})]_C$ in the rig $\nat[\bar{x}]/C$.
\end{proof}

\subsection{Gr\"{o}bner bases and multivariate polynomial division}

In this subsubsection, we develop multivariate polynomial division and establish several of its key properties. Having shown that the high polynomials of a quotient of a free rig form a ring, we may now work within a ring of polynomials rather than a rig.

The purpose of this subsubsection is to introduce Gr\"obner bases and multivariate division in a form that can later be related back to the semiring setting. In particular, multivariate division provides a criterion for determining when a polynomial is in a particular ideal and, equivalently, when two polynomials are equal in a quotient of a free ring. Combined with the previous results, this allows us to decide when those polynomials represent the same equivalence class under a corresponding congruence in a rig.

Let $\bar{x} = \{x_1,...,x_n\}$ denote a fixed finite set of variables throughout this subsubsection.

In this subsubsection, we make use of the lexicographic ordering on monomials in a ring of polynomials.  This ordering is defined as follows:

\begin{definition}\label{def:lex}
The \emph{lexicographic ordering on monomials} is a strict total order on monomials defined as follows.  

Given two monomials $m_1 = c_1x_1^{i_1}x_2^{i_2}...x_n^{i_n}$ and $m_2 = c_2x_1^{j_1}x_2^{j_2}...x_n^{j_n}$, we consider their exponent vectors $v_1 = (i_1, i_2, ..., i_n)$ and $v_2 = (j_1, j_2, ... j_n)$.  If the first non-zero component of $v_1 - v_2$ is positive, then $m_1 > m_2$.  If the first non-zero component of $v_1 - v_2$ is negative, then $m_2 > m_1$.
\end{definition}

In this section, we think of a polynomial as a sum of monomials where terms with the same exponent vectors are grouped together and the monomials that form this sum are ordered in lexicographically descending order.\

\begin{definition} \hfill 
\begin{enumerate}[nosep]
\item    The \textit{leading term} of a polynomial is its lexicographically greatest monomial.\\  We denote the leading term of a polynomial $p$ as $lt(p)$.  
\item The \textit{ideal of leading terms} of an ideal $I$ is the ideal generated by $\{lt(q)\ |\ q \in I\}$.\\  We denote the ideal of leading terms of an ideal $I$ as $LT(I)$.
\end{enumerate}
\end{definition}

\begin{remark}
    Note that the lexicographic ordering of monomials can be extended to polynomials by considering only the leading term of each polynomial.
\end{remark}

\begin{definition}
    A \emph{monomial ideal} is an ideal generated by monomials.
\end{definition}

\begin{lemma}\label{lemma:monideal}
    A polynomial is in a monomial ideal $I$ if and only if each of its monomials is a multiple of a generator of $I$.
\end{lemma}
\begin{proof}
    We prove that $f \in I$ implies that each of $f$'s monomials is a multiple of a generator of $I$.  The reverse implication is trivial.
    
    Let $I = (m_1, m_2, ..., m_n)$.  Since $f \in I$, it can be written as a linear combination of generators of $I$:
    \[ f = \sum_i (g_im_i)\]
    Each $g_i$ can be written as a finite sum of monomials
    \[ g_i = \sum_k (g_{i_k}) \]
    Now $f$ can be written as
    \[ f = \sum_i \sum_k (g_{i_k}m_i)\]
    This is a sum of monomials each of which is a multiple of some $m_i$.
    This remains the case when like terms are grouped together.
%    Each component of the above sum is a monomial that is a multiple of a generator $m_i \in I$.  Therefore, every monomial in $f$ is in $I$.
\end{proof}


A monomial ordering is necessary to perform division on a multivariate polynomial because to divide a polynomial $f$ by a polynomial $g$, one must say in what sense some polynomial $q$ is the quotient.  Given our monomial ordering from Definition \ref{def:lex}, we say that $q$ is the quotient of $f$ divided by $g$ when $f = g\times q + r$ and no monomial in $r$ is divisible by the leading term of $g$.  In general, we divide a multivariate polynomial $f$ by a set $G$ of polynomials, where $G = \{g_1, ..., g_n\}$.  The set $\{q_1, ..., q_n\}$ is the quotient of $f$ divided by $G$ when $f = (\sum_{i=1}^n g_i \times q_i) + r$ and no monomial in $r$ is divisible by the leading term of any $g_i \in G$.  The remainder $r$ of the division of $f$ by $G$ is a representative of $f$ in the ring of polynomials to which $f$ and each $g_i$ belong modulo the ideal generated by $G$. Ideally, if $f$ is an element of the ideal generated by $G$, then $f = (\sum_{i=1}^n g_i \times q_i)$, i.e., $r = 0$.  Given an arbitrary set of polynomials $\{g_1, ..., g_n\}$, this is not necessarily the case.  In order to determine if $f \in (g_1, ..., g_n)$, we need to divide $f$ not by $G$, but by a Gr\"{o}bner basis for the ideal generated by $G$.
\Needspace{5\baselineskip}
\begin{definition}\label{def:gb}
    A \emph{Gr\"{o}bner basis} for an ideal $I$ is a set of polynomials $G=\{g_1, ..., g_n\}$ such that
    \begin{enumerate}[label=(\alph*),nosep]
        \item Each $g_i \in I$.
        \item The leading terms of $G$ generate the ideal of leading terms of $I$. That is, $LT(I) = (lt(g_1), ...,lt(g_n))$.
    \end{enumerate}
\end{definition}

\begin{prop}\label{prop:gbisi}
    Let $G$ be a Gr\"obner basis for an ideal $I$. The ideal generated by $G$ is $I$.
\end{prop}

\begin{proof}
    Let $\langle G \rangle$ be the ideal generated by $G$. Let $\langle LT(G) \rangle$ be the ideal generated by the leading terms of $G$, that is, 
    $\langle LT(G) \rangle = (lt(G_1), ...,lt(G_n))$.
    Condition (a) of Definition \ref{def:gb} shows that $\langle G \rangle \subseteq I$.  It remains to show that $I \subseteq \langle G \rangle$.

    Let $f \in I$.  We will show that $f \in \langle G \rangle$.  First, consider that $lt(f) \in LT(I) = \langle LT(G) \rangle$.  $\langle LT(G) \rangle$ is a monomial ideal.  By Lemma \ref{lemma:monideal}, the leading term of $I$ is in $\langle LT(G) \rangle$.  Therefore, there exists some $g_k \in G$ such that $lt(g_k)\ |\ lt(f)$, allowing us to cancel the leading term of $f$ by defining $f_1$ as follows:
    \[ f_1 = f - cx^\alpha g_k\]
    for some monomial $cx^\alpha$ such that $lt(f) = lt(cx^\alpha g_k)$.

    Since both $f$ and $g_k$ are in $I$, it follows that $f_1 \in I$.  Since $lt(f) = lt(cx^\alpha g_k)$, the leading term of $f_1$ is strictly lexicographically less than the leading term of $f$.  We can repeat this process.  Because the lexicographic ordering is well-founded, this process must terminate, at which point we derive that
    \[ f = (\sum h_ig_i) + r\]
    with each $g_i \in G$ and no term in $r$ is divisible by the leading term of any $g_k \in G$.

    Now, if we can show that the remainder $r$ is $0$, then we have determined that $f \in \langle G \rangle$, because it is a linear combination of elements of $G$.

    Assume $r \neq 0$.  $r \in I$ because it is a difference of elements of $I$, so $lt(r) \in LT(I) = \langle LT(G) \rangle$. So, there exists some $g_k \in G$ such that $lt(g_k)\ |\ lt(r)$.  This contradicts how $r$ was constructed, so $r = 0$.
\end{proof}

The proof above gives a blueprint for the multivariate polynomial division algorithm. We now define explicitly the multivariate polynomial division algorithm as follows.

\begin{definition}
The \emph{multivariate polynomial long division algorithm} is defined as follows.

Let $G = \{g_1, ..., g_n\}$ be a Gr\"obner basis for a polynomial ideal $I$.  Let $f$ be a polynomial. 

Let $f_0= f$, $r_0 = 0$. Let $i = 0$
\begin{enumerate}[nosep]
    \item If there is no $g \in G$ such that  $lt(g)\ |\ lt(f_i)$, then
    \begin{enumerate}[nosep]
        \item Let $f_{i+1} = f_i - lt(f_i)$
        \item Let $r_{i+1} = r_i + lt(f_i)$
    \end{enumerate}
    \item If there exists a $g \in G$ such that $lt(g)\ |\ lt(f_i)$, then pick $g_{k_i}$ to be a lexicographically maximal polynomial with this property.  If there are multiple polynomials in $G$ satisfying this requirement, pick any of them.
    \begin{enumerate}[nosep]
        \item Let $cx^\alpha$ be a monomial such that $lt(cx^\alpha g_{k_i}) = lt(f_i)$.
        \item Let $f_{i+1} = f_i - cx^\alpha g_{k_i}$
        \item Let $r_{i+1} = r_i$ 
    \end{enumerate} 
\end{enumerate}
Repeat this process until $f_n = 0$.  We call $r_n$ the remainder of dividing $f$ by $G$.

Since both cases result in the leading term of $f_i$ being subtracted from $f_i$, the leading term of $f_{i+1}$  is strictly less than that of $f_{i}$.  By well-foundedness of the lexicographic ordering, this process must terminate.
\end{definition}

\begin{notation}
    We denote the remainder of dividing a polynomial $f$ by a set of polynomials $G$ as $\overline{f}^G$
\end{notation}

\begin{remark}
    Note that in the above algorithm, we only ever divide by the leading term of an element of the Gr\"obner basis $G$.  If each $g_i \in G$ is monic, the only integer we ever divide $f_i$ by is $1$. This ensures that we can perform the entire division procedure in $\ints[\bar{x}]$ for a Gr\"obner basis consisting entirely of monic polynomials.
\end{remark}

\begin{lemma}\label{lemma:remmoninlti}
    Let $G$ be a Gr\"obner basis for an ideal $I$. Let $f$ be a polynomial.  No monomial of $\overline{f}^G$ belongs to $LT(I)$.
\end{lemma}
\begin{proof}
    Note that by construction, no monomial in $\overline{f}^G$ is divisible by any $lt(g_k)$ for any $g_k \in G$.  By application Lemma \ref{lemma:monideal} at each monomial in $\overline{f}^G$, no monomial in $\overline{f}^G$ is in the ideal generated by $\{lt(g_i)\ |\ g_i \in G\}$.  Since this ideal is $LT(I)$ by the definition of a Gr\"obner basis, no monomial in $\overline{f}^G$ is in $LT(I)$.
\end{proof}

The proof of Proposition \ref{prop:gbisi} establishes  that the division of $f \in I$ by a Gr\"obner basis $G$ for $I$ is $0$.  Next, we show that given any polynomial $f$, there exists a unique remainder $\overline{f}^G$ of dividing $f$ by $G$ using the multivariate polynomial division algorithm.  To show that, we will need the following lemma.

\begin{lemma}\label{lemma:remini}
    Given any polynomial $f$, and any Gr\"obner basis $G$ for an ideal $I$, $f - \overline{f}^G \in I$.
\end{lemma}
\begin{proof}
    We will prove the following, which is sufficient to prove this lemma. At every step $i$ of the division algorithm,
    \[ f - r_i \in f_i + I \]

    We proceed by induction on $i$.  When $i$ reaches $n$, and $f_i = f_n = 0$, we have that $f - \overline{f}^G \in I$.

    \begin{description}
        \item[Base Case ($i = 0$):] Here $f_0 = f$ and $r_0 = 0$, hence 
        \[ f-r_0 = f \in f + I = f_0 + I \]
        \item[Inductive Case: ] Assume 
        \[ f - r_i \in f_i + I \]
        We will show
        \[ f - r_{i+1} \in  f_{i+1} + I \]
        \begin{itemize}[nosep]
            \item Case 1: Assume there is no $g \in G$ such that $lt(g)\ |\ lt(f_i)$.
            Then $f_{i+1} = f_i - lt(f_i)$
            and $r_{i+1} = r_i + lt(f_i)$.
            We want to show
            \[ f - r_i - lt(f_i) \in f_i - lt(f_i) + I\]
            it is sufficient to show
            \[ f - r_{i} \in  f_{i} + I \]
            which is precisely our inductive hypothesis.
            \item Case 2: Assume there exists a $g \in G$ such that $lt(g)\ |\ lt(f_i)$, and let $g_{k_i}$ be a lexicographically maximal polynomial with this property. 
            
            Then $f_{i+1} = f_i - cx^\alpha g_{k_i}$
            and $r_{i+1} = r_i$.
            We want to show
            \begin{equation}\label{eq:diffinI}
                f - r_i \in f_i - cx^\alpha g_{k_i} + I
            \end{equation}
            By condition a) of Definition \ref{def:gb}, $g_{k_i} \in I$, so $cx^\alpha g_{k_i} \in I$.  We can reduce \eqref{eq:diffinI} to
            \[ f - r_i \in f_i + I \]
            which, again, is precisely our inductive hypothesis. \qedhere
        \end{itemize}
    \end{description}
\end{proof}
\Needspace{5\baselineskip}
\begin{prop}\label{prop:uniquer}
    Let $f$ be a polynomial.  Let $G$ be a Gr\"obner basis for an ideal $I$.  There exists a unique $r$ such that
    \begin{enumerate}[label=\alph*),nosep]
        \item $f - r \in I$
        \item No monomial in $r$ belongs to $LT(I)$
    \end{enumerate}
\end{prop}

\begin{proof}
    Suppose $r$ and $r'$ both satisfy a) and b) as above.
    Then $r - r' = (f - r') - (f - r)$.  $f - r' \in I$ and $f - r \in I$, so $r - r' \in I$.  If $r - r' \neq 0$, then $lt(r - r') \in LT(I)$.  No monomials in $r$ nor $r'$ belong to $LT(I)$, so this is a contradiction.  Therefore, $r - r'$ is the $0$ polynomial, and $r = r'$.
\end{proof}

\begin{lemma}\label{lemma:auxgb}
    Given any Gr\"obner basis $G$ for an ideal $I$, and any polynomial $f$, $\overline{f}^G$ satisfies conditions a) and b) of Proposition \ref{prop:uniquer}.
\end{lemma}
\begin{proof}
    $\overline{f}^G$ satisfies requirement a) of Proposition \ref{prop:uniquer} by Lemma \ref{lemma:remini}. It satisfies requirement b) of Proposition \ref{prop:uniquer} by Lemma \ref{lemma:remmoninlti}.
\end{proof}

\begin{corollary}\label{cor:uniquer}
    The multivariate polynomial division algorithm gives a unique remainder when dividing by a Gr\"obner basis.
\end{corollary}
\begin{proof}
    Assemble Proposition \ref{prop:uniquer} and Lemma \ref{lemma:auxgb}.
\end{proof}

The above corollary shows that the multivariate polynomial division algorithm computes a unique representative of a given polynomial $f$ in $\ints[\bar{x}]/I$ given that we divide $f$ by a Gr\"obner basis $I$.  It will be particularly important to us in the case where the remainder of a division is $0$.

Before addressing the existence of Gr\"obner bases in the context of ADT isomorphisms, we give an example of the application of the multivariate polynomial division algorithm.  

\begin{example}
Let $G = \{g_1,g_2\} = \{x - y, y^2 + 1\}$.  $G$ is a Gr\"obner basis for the ideal it generates, and $lt(g_1) = x$ and $lt(g_2) = y^2$.  Let $f = x^2 + y$.  We can now proceed with the algorithm in the following steps.

\begin{enumerate}[nosep]
    \item Pick $g_{k_0}$ to be $g_1$ because $lt(x-y)\ |\ lt(x^2+y)$
    \item Multiply $g_{1}$ by $x$: \[ xg_{1} = x^2 - xy \]
    \item Compute $f_1$:
    \[ f_1 = f - xg_{1} = x^2 + y - (x^2 - xy) = xy + y\]
    \item Compute $r_1$:
        \[ r_1 = r_0 \]
    \item Pick $g_{k_1}$ to be $g_1$ because $lt(x - y)\ |\ lt(xy + y)$
    \item Multiply $g_{1}$ by $y$: \[ yg_{1} = xy - y^2 \]
    \item Compute $f_2$:
    \[ f_2 = f_1 - yg_{1} = xy + y - (xy - y^2) = y^2 + y \]
    \item Compute $r_2$:
        \[ r_2 = r_1 \]
    \item Pick $g_{k_2}$ to be $g_2$ because $lt(y^2 + 1)\ |\ lt(y^2 + y)$
    \item Multiply $g_{2}$ by $1$
    \item Compute $f_3$:
    \[ f_3 = f_2 - g_{2} = y^2 + y - (y^2 + 1) = y - 1 \]
    \item Compute $r_3$:
        \[ r_3 = r_2 \]
    \item The leading term of $f_3$ is not divisible by the leading term of any $g \in G$, so
    \begin{align*}
        f_4 &= f_3 - y = 1 \\
        r_4 &= r_3 + y = y
    \end{align*}
    \item The leading term of $f_4$ is not divisible by the leading term of any $g \in G$, so
    \begin{align*}
        f_5 &= f_4 + 1 = 0 \\
        r_5 &= r_4 - 1 = y - 1
    \end{align*}
    \item $f_5$ is now $0$, so the division has terminated, and
    \[ \overline{f}^G = r_5 = y - 1 \]
\end{enumerate}

From this division, we have explicitly that
\[ f = xg_1 + yg_1 + g_2 + (y - 1)\]
and that
\[ f - \overline{f}^G = (x+y)g_1 + g_2 \in (g_1, g_2) \]
\end{example}

One necessary and sufficient condition to determine whether a set of polynomials is a Gr\"obner basis for the ideal it generates is {\em Buchberger's Criterion}.  To state Buchberger’s Criterion, we first introduce two definitions.

\begin{definition}
The \emph{least common multiple} ($LCM$) of monomials is defined by
\[ LCM(c_1x_1^{\alpha_1}\dots x_n^{\alpha_n}, c_2x_1^{\beta_1}\dots x_n^{\beta_n}) = lcm(c_1, c_2)x_1^{max(\alpha_1, \beta_1)} \cdots x_n^{max(\alpha_n, \beta_n)} \]
where $lcm$ is the least common multiple on the integers.
\end{definition}

\begin{definition}
    Two monomials $m_1$ and $m_2$ are said to be \emph{coprime} if
    \[ LCM(m_1, m_2) = m_1m_2 \]
\end{definition}


\begin{definition}
The \emph{$S$ polynomial} of two polynomials $f$ and $g$ is defined by
\[ S(f,g) = \frac{LCM(lt(f), lt(g))}{lt(f)}f - \frac{LCM(lt(f), lt(g))}{lt(g)}g \]
\end{definition}

We can now give a precise statement of Buchberger's Criterion in the following theorem.

\begin{theorem}[Buchberger's Criterion]
    A set of polynomials $G=\{g_1, ..., g_n\}$ is a Gr\"obner basis if and only if $\overline{S(g_i, g_j)}^G = 0$ for all $i \neq j$.
\end{theorem}
For a proof of this classic result, see Chapter 9, Proposition 26 of Dummit and Foote \cite{dummitfoote2003}.

The following lemma is useful for proving our main theorem about the relationship between complex numbers and rig isomorphisms.

\begin{lemma}\label{lemma:monicS}
    Given two monic polynomials $f$ and $g$, if $lt(f)$ and $lt(g)$ are coprime, then $\overline{S(f,g)}^{\{f,g\}} = 0$.
\end{lemma}
\begin{proof}
    Let $f = x^\alpha + p_1$ and $g = x^\beta + p_2$ for some exponent vectors $\alpha$ and $\beta$ where $lt(f) = x^\alpha$ and $lt(g) = x^\beta$. Then
     \begin{align*}
         S(f,g) &= \frac{LCM(lt(f), lt(g))}{lt(f)}f - \frac{LCM(lt(f), lt(g))}{lt(g)}g\\
         &= \frac{lt(f)lt(g)}{lt(f)}f - \frac{lt(f)lt(g)}{lt(g)}g \\
         &= lt(g)f - lt(f)g \\
         &= x^\beta (x^\alpha + p_1) - x^\alpha (x^\beta + p_2) \\
         &= x^\beta p_1 - x^\alpha p_2
     \end{align*}

     Consider the expression $p_1g - p_2f$.
     \[ p_1g - p_2f = p_1x^\beta + p_1p_2 - (p_2x^\alpha + p_2p_1) = x^\beta p_1 - x^\alpha p_2 \]
     By the above equation, $p_1g - p_2f$ is precisely equal to $S(f, g)$.  Because we have found a linear combination of $f$ and $g$ that is equal to $S(f, g)$ and with coefficients ($p_1$ and $-p_2$) that are lexicographically less than $S(f, g)$, we know that $\{f, g\}$ cleanly divides $S(f, g)$, i.e., $\overline{S(f,g)}^{\{f,g\}} = 0$. 
\end{proof}

\begin{corollary}\label{cor:coprimeGB}
    Any set of monic polynomials with pairwise coprime leading terms form a Gr\"obner basis for the ideal they generate.
\end{corollary}
\begin{proof}
    Let $G=\{g_1, ..., g_n\}$ be a set of polynomials with pairwise coprime leading terms.  By Lemma \ref{lemma:monicS}, for all $i, j$ where $i \neq j$, $\overline{S(g_i,g_j)}^{\{g_i,g_j\}} = 0$.  If the set $\{g_i, g_j\}$ cleanly divides $S(g_i,g_j)$, then any superset of $\{g_i, g_j\}$ also cleanly divides $S(g_i,g_j)$.  Therefore, $\overline{S(g_i,g_j)}^{G} = 0$.  By Buchberger's Criterion, $G$ is a Gr\"obner basis for the ideal generated by $G$.
\end{proof}

The final classic fact from abstract algebra that we will need is the following.

\begin{prop}\label{prop:existsgb}
    Given any set of equations $E$ over $\ints[\bar{x}]$, there exist a Gr\"obner basis for $\sim_E$ considered as an ideal.
\end{prop}
\begin{proof}
    For this proof, we refer to ``Ideals, Varieties, and Algorithms'', by Cox, Little, and O'Shea \cite{cox2010}.  By Hilbert Basis Theorem (Theorem 5.4 in \cite{cox2010}), there exists a finite set of polynomials generating the ideal $\sim_E$.  By Theorem 7.2 (loc. cit.), there exists a Gr\"obner basis for any ideal generated by a finite set of polynomials.
\end{proof}

In the next subsubsection, we prove our main theorem about multivariate rig isomorphisms.

\subsection{The main theorem for the multivariate case}

The last missing piece in proving our main theorem about multivariable rig isomorphisms is determining when two polynomials $q_1$ and $q_2$ are equal in a quotient ring.  Specifically, we consider the quotient ring $\ints[x_1, ..., x_n]/I$, where $I$ is an ideal generated by a Gr\"obner basis $G = \{g_1, \dots, g_n\}$, and each $g_i$ is univariate in $x_i$ and has no repeated complex roots.  Identically, we can ask when is $q_1 - q_2$ an element of an ideal such as $I$. We now know that we can divide $q_1-q_2$ by the Gr\"obner basis $G$ (see Corollary \ref{cor:coprimeGB}) to obtain that

\[ q_1-q_2 = \left(\sum_{i=0}^n g_i(x_i)t_i(\bar{x})\right) + r\]

Our question now becomes, when is the remainder $r$ of this division procedure the zero polynomial?  The following lemma used in proving the Combinatorial Nullstellensatz answers exactly this.

\begin{lemma}\label{lemma:nullstellensatzlemma}
Let $r \in \mathbb{C}[x_1, ..., x_n]$ be a polynomial vanishing on a grid $S_1 \times \dots \times S_N$, where $S_i \subseteq \mbC$, and $\deg_{x_i}(r) < |S_i|$ for all $i$. Then $r$ must be the zero polynomial.
\end{lemma}
\begin{proof}
    We proceed by induction on the number of variables, $n$.
    \begin{description}
        \item[Base case ($n = 1$):]  We have that $r \in \mathbb{C}[x_1]$, that $r$ vanishes on the grid $S_1$, and that $\deg_{x_1}(r) < |S_1|$, i.e., $r$ has more roots than its degree.  By the fundamental theorem of algebra, $r$ must be the zero polynomial.
        \item[Inductive Step:] Assume that the theorem holds for a polynomial in $n$ variables.

        We have that $r \in \mathbb{C}[x_1, ..., x_n, x_{n+1}]$, that $r$ vanishes on the grid $S_1 \times ... \times S_{n+1}$, and that $\deg_{x_i}(r) < |S_i|$ for all $i$.  We may view $r$ as a polynomial in the variable $x_{n+1}$ with coefficients from $\mathbb{C}[x_1, ..., x_n]$.
        Let $r = \sum_{i=0}^{d_N}(c_i(x_1, ..., x_n)x_{n+1}^i$, where $d_N = deg_{x_{n+1}}(r)$.  Pick an arbitrary element $s_i$ from each set $S_i$ where $0 \le i \le n$. Consider the polynomial defined by 
        \[ p(x_{n+1}) = \sum_{i=0}^{d_N}(c_i(s_1, ..., s_n)x_{n+1}^i) \]
        Because $\deg_{x_{n+1}}(p) < S_{n+1}$, by the base case, $p$ must be the zero polynomial.  If $p$ is the zero polynomial, then all of its coefficients $c_i(s_1, ..., s_n)$ must be 0.  Because our choice of $s_1$ through $s_n$ was done by picking arbitrary elements of $S_1, ..., S_n$, we know that $c_i(x_1, ... x_n)$ vanishes on the grid $S_1 \times ... \times S_n$.  By hypothesis, we have that $\deg_{x_i}(r) < |S_i|$ for all $i$.  We can now apply the inductive hypothesis to conclude that $c_i(x_1, ..., x_n)$ is the zero polynomial for all $i$. \qedhere
    \end{description}
\end{proof}

Before stating our main theorem, we refine one of our earlier corollaries from Chapter \ref{sec:eqtheory} to the current context.

\begin{lemma}
    Let $E$ be a set of equations in $\ints[x_1,..., x_n]$.  Let $q_1, q_2 \in \ints[x_1, ..., x_n]$.
    Then a Gr\"obner basis for $\sim_E$ exists, and all of the following are equivalent:
    \begin{enumerate}[label=(\arabic*),nosep]
        \item $E \models_{Ring} q_1 = q_2$
        \item $[q_1]_{\sim_{E}} = [q_2]_{\sim_{E}}$ in the quotient ring $\ints[x_1, ..., x_n]/\sim_{E}$.
        \item Any Gr\"obner basis for $\sim_E$ divides $q_1 - q_2$ with remainder $0$. \qedhere
    \end{enumerate}
\end{lemma}
\begin{proof}
    The existence of Gr\"obner basis follows from Proposition \ref{prop:existsgb}.

    The equivalence of the three conditions can now be established as follows.
    \begin{description}
        \item[$(1) \iff (2)$] By Corollary \ref{cor:rigmodels}.
        \item[$(2) \implies (3)$] Let $G$ be any Gr\"obner basis for $\sim_E$.  We know that if $[q_1]_{\sim_E} = [q_2]_{\sim_E}$, then $[q_1 - q_2]_{\sim_E} = 0$.  That is, $q_1 - q_2$ is in the ideal $\sim_E$.  Proposition \ref{prop:uniquer} says that there is a unique polynomial that simultaneously satisfies conditions a) and b).  Thus, if $0$ satisfies conditions a) and b) then, because $\overline{f}^G$ satisfies a) and b) by Lemma \ref{lemma:auxgb}, the remainder $\overline{f}^G$ is $0$.  By hypothesis, $0$ satisfies a).  Condition b) is vacuously true since $0$ is the empty sum of monomials.  Therefore, $\overline{f}^G = 0$.
        \item[$(3) \implies (2)$] Let $G$ be a Gr\"obner basis for $\sim_E$.  Let $\overline{q_1 - q_2}^G = 0$.  Then by Lemma \ref{lemma:remini}, $q_1 - q_2$ is in the ideal $\sim_E$, and $[q_1]_{\sim_E} = [q_2]_{\sim_E}$. \qedhere
    \end{description}
\end{proof}

We now have all we need to prove our main theorem as stated below.

\begin{theorem}\label{thm:mvringtorig}
    Let $p_1, \dots, p_n \in \nat[x_n]$ be polynomials each of which has a non-zero constant term and degree at least 2.  Let $C = (x_1 = p_1(x_1), ..., x_n = p_n(x_n))$ be the least congruence on $\nat[x_1, ..., x_n]$ generated by the given relations.

    Let $q_1, q_2 \in H_C(\nat[x_1, ..., x_n])$ be polynomials.

    Then 
    \begin{align*}
    &x_1 = p_1(x_1), ..., x_n = p_n(x_n) \models_{Ring} q_1(x_1, ..., x_n) = q_2(x_1, ..., x_n) \\
    &\implies 
    x_1 = p_1(x_1), ..., x_n = p_n(x_n) \models_{Rig} q_1(x_1, ..., x_n) = q_2(x_1, ..., x_n)
    \end{align*}
\end{theorem}

\begin{proof}
    Assemble Corollary \ref{cor:rigmodels} and Corollary \ref{cor:equalityHNtoN}.
\end{proof}

\begin{theorem}\label{thm:mvrigiso}
    Let $p_1 \in \nat[x_1]$, ..., $p_n \in \nat[x_n]$ be polynomials with a non-zero constant term and degree at least 2.  Additionally, suppose that for all $i$, $p_i(x_i) - x_i$ is monic and has no repeated complex roots.  Let $C = (x_1 = p_1(x_1), ..., x_n = p_n(x_n))$ be the least congruence on $\nat[x_1, ..., x_n]$ generated by the given relations.
    
    Let $q_1, q_2 \in H_C(\nat[x_1, ... x_n])$ be polynomials such that for all complex numbers $a_1, ..., a_n$ satisfying $p_i(a_i) = a_i$ for all $i$, $q_1(a_1, ..., a_n) - q_2(a_1, ..., a_n) = 0$, i.e., $q_1 - q_2$ vanishes on the roots of $p_1(x_1) - x_1, ..., p_n(x_n) - x_n$.

    Then $\{x_i = p_i(x_i)\ |\ 1 \le i \le n\} \models_{Rig} q_1(\bar{x}) = q_2(\bar{x})$.
\end{theorem}

\begin{proof}
    Let $p_1, ..., p_n$ and $q_1, q_2$ be polynomials as described in the statement of Theorem \ref{thm:mvrigiso}.

    Let $f_i(x_i) = p_i(x_i) - x_i$ for all $i$.
    
    By Corollary \ref{cor:coprimeGB}, $\{f_1, ..., f_n\}$ forms a Gr\"obner Basis for the ideal $(f_1, ..., f_n)$ of $\ints[x_1, ..., x_n]$.  Because each $f_i$ is monic, $q_1(\bar{x}) - q_2(\bar{x})$ can be divided by $\{f_1, ..., f_n\}$ completely in $\ints$ as integer division is only ever performed with the leading coefficient of an element of the Gr\"obner basis as a divisor.  Below is the result of this division:

    \[ q_1(\bar{x}) - q_2(\bar{x}) = \left(\sum_{i=1}^nf_i(\bar{x})t_i(x_i)\right) + r\]
    
    The multivariate polynomial division algorithm ensures that no monomial in $r$ is divisible by the leading term of any $f_i$.  Because each $f_i$ contains only the variable $x_i$, this means that for all $i$, $deg_{x_i}(r) < deg_{x_i}(f_i)$.

    Let $S_i$ be the set of all complex roots of $f_i$ for each $i$.  By hypothesis, $q_1(\bar{x}) - q_2(\bar{x})$ vanishes on $S_1 \times ... \times S_n$.  Because each $f_i$ has no repeated roots, $\deg_{x_i}(r) < |S_i| = deg_{x_i}(f_i)$.  We can now apply Lemma \ref{lemma:nullstellensatzlemma} to conclude that $r$ is the zero polynomial.

    Because $r$ is the zero polynomial, $q_1(\bar{x}) - q_2(\bar{x}) \in (f_1, ..., f_n)$, and $q_1(\bar{x}) = q_2(\bar{x})$ in $\ints[\bar{x}]/(f_1, ..., f_n)$.  By Corollary \ref{cor:equalityHNtoN}, $q_1(\bar{x}) = q_2(\bar{x})$ in $\nat[\bar{x}]/(x_1 = p_1(x_1), ..., x_n=p_n(x_n))$.  The ideal $(x_1 = p_1(x_1), ..., x_n=p_n(x_n))$ quotients the polynomial ring $\nat[\bar{x}]$ by the equations from our hypothesis: $\forall i. x_i=p_i(x_i)$, so this is exactly what we wanted to show.
\end{proof}

\subsection{An example derivation of isomorphism from equality in the multivariate case}

In the remainder of this chapter, we use Theorem \ref{thm:mvrigiso} to derive an explicit nontrivial isomorphism between ADTs involving multiple distinct closed recursive polynomial types.  The isomorphism is nontrivial in the sense that it cannot be derived using Fiore and Leinster's results alone.  In generating this isomorphism, we use only the lemmas and propositions that are necessary for our case in order to simplify the presentation of the example.  We reduce to the univariate case where possible, so that certain isomorphisms can be generated with Fiore and Leinster's results alone.

Let $K, L : ADT_\emptyset$ be defined by
\begin{align*}
    K &= \mu X.X^2 + X + \mb{2}  \\
    L &= \mu X.X^2 + 2X + \mb{1}
\end{align*}

The fixed point isomorphisms of $K$ and $L$ tell us
\begin{align*}
    K &\simeq K^2 + K + \mb{2} \\
    L &\simeq L^2 + 2L + \mb{1}
\end{align*}

If we consider the above isomorphisms as polynomial equations, we get
\begin{align*}
    x_1 &= x_1^2 + x_1 + 2 = p_1(x_1) \\
    x_2 &= x_2^2 + 2x_2 + 1 = p_2(x_2)
\end{align*}
The roots of the first equation, which corresponds to the data type $K$, are $\pm i\sqrt{2}$.  The roots of the second equation, which corresponds to the data type $L$, are $\pm e^{2 \pi i/3}$.  Note that $x_1^2+x_1+2$ and $x_2^2+2x_2+1$, when they are considered as $p_1$ and $p_2$, respectively, for Theorem \ref{thm:mvrigiso}, indeed satisfy the theorem's hypotheses that each polynomial has a non-zero constant term and degree at least $2$, and that $p_1(x_1) - x_1$ and $p_2(x_2) - x_2$ are monic and have no repeated roots.  That tells us that given any two high polynomials $q_1$ and $q_2$ such that $q_1(x_1, x_2) = q_2(x_1,x_2)$ for all roots $x_1$ of $x_1 = x_1^2 + x_1 + 2$ and $x_2$ of $x_2 = x_2^2 + 2x_2 + 1$, there there exists an isomorphism
\[ q_1(K, L) \simeq q_2(K, L)\]

Consider that $x_1 = \pm i\sqrt{2}$, and $x_2 = \pm e^{2 \pi i/3}$, so 
\[ x_1^4x_2^3 = 4(1) = 2(-2)(-1) = 2x_1^2(x_2^2 + x_2) = 2x_1^2x_2^2 + 2x_1^2x_2\]
By the above equation, the polynomials $q_1$ and $q_2$ defined by
\begin{align*}
    q_1(x_1,x_2) &= x_1^4x_2^3 \\
    q_2(x_1,x_2) &= 2x_1^2x_2^2 + 2x_1^2x_2
\end{align*}
satisfy our requirements; they contain a monomial with every variable, so they are high.  We can use Theorem \ref{thm:mvrigiso} to generate an isomorphism
\[ q_1(K,L) = K^4L^3 \simeq 2K^2L^2 + 2K^2L = q_2(K,L)\]

Since the proof is constructive, we can use the proof as an algorithm to generate an explicit isomorphism. We now proceed to do that, omitting extraneous detail where possible.

We begin by dividing $q_1(x_1, x_2) - q_2(x_1,x_2)$ by the Gr\"obner basis $\{p_1(x_1) - x_1, p_2(x_2) - x_2\}$ (see Corollary \ref{cor:coprimeGB}) to find that
\[ x_1^4x_2^3 - (2x_1^2x_2^2 + 2x_1^2x_2) = ((x_1^2 + x_1 + 2) - x_1)(x_1^2x_2^3 - 2x_2^3 - 2x_2^2 - 2x_2) + ((x_2^2 + 2x_2 + 1) - x_2)(4x_2) \]
The division algorithm derives an equality in $\ints[x_1,x_2]$, so both sides of the above equation are the same polynomial.  Adding the negatives to both sides of the equation to obtain an equation entirely in $\nat[x_1,x_2]$ maintains this fact, and we find
\begin{align*}
    &x_1^4x_2^3 + x_1(x_1^2x_2^3) + (x_1^2 + x_1 + 2)(2x_2^3 + 2x_2^2 + 2x_2) + 4x_2^2 \\
    &= 2x_1^2x_2^2 + 2x_1^2x_2 + x_1^2x_2^3(x_1^2 + x_1 + 2) + x_1(2x_2^2 + 2x_2^2 + 2x_2) + 4x_2(x_2^2 + 2x_2 + 1)
\end{align*}
Because both sides of the above equation are the same polynomial, we can replace $x_1$ with $K$ and $x_2$ with $L$ and there must be the following isomorphism using only the isomorphisms corresponding to the rig axioms.
\begin{align*}
    &K^4L^3 + K(K^2L^3) + (K^2 + K + \mb{2})(\mb{2}L^3 + \mb{2}L^2 + \mb{2}L) + \mb{4}L^2 \\
    &\simeq \mb{2}K^2L^2 + \mb{2}K^2L + K^2L^3(K^2 + K + \mb{2}) + K(\mb{2}L^2 + \mb{2}L^2 + \mb{2}L) + \mb{4}L(L^2 + \mb{2}L + \mb{1})
\end{align*}
Applying the fixed point isomorphisms for $K$ and $L$ where applicable, we obtain
\begin{align*}
    &K^4L^3 + K(K^2L^3) + (K)(\mb{2}L^3 + \mb{2}L^2 + \mb{2}L) + \mb{4}L^2 \\
    &\simeq \mb{2}K^2L^2 + \mb{2}K^2L + K^2L^3(K) + K(\mb{2}L^2 + \mb{2}L^2 + \mb{2}L) + \mb{4}L(L)
\end{align*}
Simplifying using the rig axioms, we get
\begin{equation}\label{eq:mvexwiths}
\begin{aligned}
    &K^4L^3 + K^3L^3 + \mb{2}KL^3 + \mb{2}KL^2 + \mb{2}KL + \mb{4}L^2 \\
    &\simeq \mb{2}K^2L^2 + \mb{2}K^2L + K^3L^3 + \mb{2}KL^2 + \mb{2}KL^2 + \mb{2}KL + \mb{4}L^2
\end{aligned}
\end{equation}
Let $S = K^3L^3 + \mb{2}KL^3 + \mb{2}KL^2 + \mb{2}KL + \mb{4}L^2$.  It remains to ``cancel off'' $S$ from both sides of the isomorphism 
\[ K^4L^3 + S \simeq \mb{2}K^2L^2 + \mb{2}K^2L + S \]
To do this, we will instantiate several lemmas from this chapter to this specific example.

By Corollary \ref{cor:uniquez}, the zero element of the ring of high polynomials, which we will call $z$, is defined by
\[ z = (K^2 + \mb{2})(L^2 + L + \mb{1}) \]
because, by $L$'s fixed point isomorphism, 
\[ (K^2 + \mb{2})L + z \simeq (K^2 + \mb{2})(L^2 + \mb{2}L + \mb{1}) \simeq (K^2 + \mb{2})L \]
That is, $z$ witnesses $(K^2+\mb{2})L \le (K^2+\mb{2})L$.

Note that $z = z_Kz_L$ where $z_K$ is the zero of the ring of high polynomials in $K$ and $z_L$ is the zero of the ring of high polynomials in $L$.  This fact will be important to us.

To cancel $S$ from the isomorphism given in \eqref{eq:mvexwiths}, we first turn to Lemma \ref{lemma:az=z}.  This lemma establishes that in a semiring of high elements, multiplying an element by zero yields zero.  We can instantiate this lemma to our particular case using the following lemmas.

\begin{lemma}\label{lemma:z+z=z}
For all $n \in \nat$,
\[ \mb{n}z \simeq z\]
\end{lemma}
\begin{proof}
    We prove that $\mb{2}z \simeq z$. The desired result then follows by induction on $n$.
    $\mb{2}z \simeq z$ can be proven by first simplifying to $z_Kz_L + z_Kz_L$, after which the result reduces to the univariate case:
    \[ \mb{2}z \simeq z_Kz_L + z_Kz_L \simeq (z_K + z_K)z_L \simeq z_Kz_L \simeq z \qedhere\]
\end{proof}

\begin{lemma}\label{lemma:mz=z}
    Multiplying any monomial in $\nat[K,L]$ by $z$ yields $z$.
\end{lemma}
\begin{proof}
    For this proof, we will show
    \begin{itemize}[nosep]
        \item $Kz \simeq z$
        \item $Lz \simeq z$
        \item $\mb{n}z \simeq z$
    \end{itemize}
    After showing the above, we can use induction to finish the proof.
    
    $Kz \simeq z$ follows from considering, in the univariate case, that $Kz_K \simeq z_K$.  From this, using the congruence of $\simeq$, we can derive that 
    \[ Kz \simeq Kz_Kz_L \simeq z_Kz_L \simeq z \]
    The proof that $Lz \simeq z$ is similar.
    $\mb{n}z \simeq z$ is proved by Lemma \ref{lemma:z+z=z}.
    

    Now, consider a monomial $\mb{n}K^iL^j$.  By induction on $j$, we can prove $\mb{n}K^iL^jz \simeq \mb{n}K^iz$. By induction on $i$, we can prove $\mb{n}K^iz \simeq \mb{n}z$.  By Lemma \ref{lemma:z+z=z}, $\mb{n}z = z$.  Therefore, for any monomial $m$, $mz \simeq z$.
\end{proof}

We next must show that given any high polynomial $R$, $R + z \simeq z$. Because every high polynomial (including $q_1(K, L)$ and $q_2(K, L)$) contains at least one high monomial, it suffices to show that given any high monomial $m$, $m + z \simeq z$.  For this proof, we make use of the following lemma.

\begin{lemma}\label{lemma:auxz1}
    The following isomorphisms exist in $\nat[K,L]$.
    \[ z \simeq  Kz_L\]
    and
    \[ z \simeq z_KL\]
\end{lemma}
\begin{proof}
    We derive an isomorphism for $z \simeq Kz_L$ only.  The derivation for $z \simeq z_KL$ is similar.
    
    First, we can show that
    \begin{equation}\label{eq:auxzkzl1}
        z = z_Kz_L \simeq z_Kz_L + z_L
    \end{equation}
    by the following chain of isomorphisms:
    \begin{equation*}
    \begin{aligned}
        z_Kz_L 
        &= (K^2 + \mb{2})z_L \\
        &\simeq (K^2 + \mb{1})z_L + z_L \\
        &\simeq (K^2 + \mb{1})z_L + z_L + z_L && \textsf{(by the univariate case)} \\
        &\simeq (K^2 + \mb{2})z_L + z_L \\
        &= z_Kz_L + z_L
    \end{aligned}
    \end{equation*}
    Notice that to reduce to the univariate case, we need only factor out $z_L$ from $z_Kz_L$.

    We now derive the desired isomorphism below.
    \begin{align*}
        z 
        &= z_Kz_L \\
        &\simeq Kz_Kz_L && \textsf{(by Lemma \ref{lemma:mz=z})} \\
        &\simeq K(z_Kz_L + z_L)  && \textsf{(by equation \eqref{eq:auxzkzl1})} \\
        &\simeq Kz_Kz_L + Kz_L \\
        &\simeq z_Kz_L + Kz_L && \textsf{(by Lemma \ref{lemma:mz=z})} \\
        &\simeq (z_K + K)z_L \\
        &\simeq Kz_L && \textsf{(by the univariate case)}
    \end{align*}
\end{proof}

The previous lemma lets us derive isomorphisms between $z + m$ and $m$ for any high monomial $m$.  This is done in the following lemma.

\begin{lemma}\label{lemma:m+z=z}
    Let $m$ be a high monomial in $\nat[K,L]$.  Then $m + z \simeq z$.
\end{lemma}
\begin{proof}
    Let $m = \mb{n}K^iL^j$ with $i, j \ge 1$.  We can derive the desired isomorphism using the following chain of isomorphisms.
    \begin{align*}
        m + z 
        &= \mb{n}K^iL^j + z \\
        &\simeq \mb{n}K^iL^j + zL^{j-1} && \textsf{(by Lemma \ref{lemma:mz=z})} \\
        &\simeq \mb{n}K^iL^j + z_KL^j && \textsf{(by Lemma \ref{lemma:auxz1})} \\
        &\simeq (\mb{n}K^i + z_K)L^j \\
        &\simeq \mb{n}K^iL^j &&\textsf{(by the univariate case)} \\
        &= m
    \end{align*}
\end{proof}

We now have all we need to derive the explicit isomorphism
\[ q_1(K,L) \simeq q_2(K,L) \]
starting only from the isomorphism
\[ q_1(K,L) + S \simeq q_2(K,L) + S \]
of equation \ref{eq:mvexwiths}, where $S = K^3L^3 + 2KL^3 + 2KL^2 + 2KL + 4L^2$.

First, consider that we can obtain the additive inverse of $S$ in $H(\nat[K,L])$ by subtracting $1$ from $z$ and multiplying by $S$. $z$ is high, so every element is related to it by $\le$.  In particular, $\mb{1} \le z$.  We can expand $z$ in the following way.
\begin{equation}\label{eq:expandz}
z = z_Kz_L = (K^2 + \mb{2})(L^2 + L + \mb{1}) \simeq K^2L^2 + K^2L + K^2 + \mb{2}L^2 + \mb{2}L + \mb{2}
\end{equation}
Clearly, then, if we add $1$ to 
\[ -\mb{1}_H = K^2L^2 + K^2L + K^2 + \mb{2}L^2 + \mb{2}L + \mb{1} \]
we get back $z$.  Therefore equation \eqref{eq:expandz} becomes an isomorphism between $-\mb{1}_H + \mb{1}$ and $z$.

The isomorphism that we must sum with the isomorphism of equation \ref{eq:mvexwiths} is therefore the identity isomorphism for $-\mb{1}_H S$.

The desired isomorphism is now given as follows.
\begin{align*}
    q_1(K,L)
    &= K^4L^3 \\ 
    &\simeq K^4L^3 + z && \textsf{(by Lemma \ref{lemma:m+z=z})} \\
    &\simeq K^4L^3 + z + z + z + z + z && \textsf{(by Lemma \ref{lemma:z+z=z})} \\
    &\simeq K^4L^3 + K^3L^3z + \mb{2}KL^3z + \mb{2}KL^2z + \mb{2}KLz + \mb{4}L^2z &&\textsf{(by Lemma \ref{lemma:mz=z})} \\
    &\simeq K^4L^3 + Sz && \textsf{(by distributivity)} \\
    &\simeq K^4L^3 + S + S(-\mb{1}_H) && \textsf{(by equation \eqref{eq:expandz})} \\
    &\simeq \mb{2}K^2L^2 + \mb{2}K^2L + S + S(-\mb{1}_H) && \textsf{(by equation \eqref{eq:mvexwiths})} \\
    &\simeq \mb{2}K^2L^2 + \mb{2}K^2L + Sz && \textsf{(by equation \eqref{eq:expandz})} \\
    &\simeq \mb{2}K^2L^2 + \mb{2}K^2L + K^3L^3z + \mb{2}KL^3z + \mb{2}KL^2z + \mb{2}KLz + \mb{4}L^2z && \textsf{(by distributivity)} \\
    &\simeq \mb{2}K^2L^2 + \mb{2}K^2L + z + z + z + z + z &&\textsf{(by Lemma \ref{lemma:mz=z})} \\
    &\simeq \mb{2}K^2L^2 + \mb{2}K^2L + z && \textsf{(by Lemma \ref{lemma:z+z=z})} \\
    &\simeq \mb{2}K^2L^2 + \mb{2}K^2L && \textsf{(by Lemma \ref{lemma:m+z=z})} \\
    &= q_2(K,L)
\end{align*}

While we used particular instances of the theorems and lemmas from this chapter to generate the isomorphism from the above example, the methodology can be extended to any isomorphism specified by Theorem \ref{thm:mvrigiso}.

\chapter{Mu Rules}\label{sec:murules}

In the previous chapter, we introduced the concept of ``rig isomorphisms''.  These isomorphisms capture the fact that algebraic data types form a rig, when they are considered modulo certain natural isomorphisms.  $ADT_\vars$ form a $\Sigma_{Rig}$-algebra with $\times$ and $+$ as the two binary operators and $\mb{1}$ and $\mb{0}$ as the two nullary operators.  When we quotient $ADT_\vars$ by the natural isomorphisms obtained from the fact that we interpret them in a rig category, we get a semiring, $\ADTrig$.  ADTs have an additional constructor, $\mu$.  To have a full account of the isomorphisms between ADTs would require identification of those isomorphisms in $\mcC$ involving interpretations of ADTs constructed using $\mu$ that hold universally in any category equipped with the necessary structure to interpret ADTs.  Here, such a category is an $\omega$-cocomplete rig category. 
%\hi{Refer to $\mu$-complete categories?}

Also in the previous chapter, we introduced the notion of ``fixed point isomorphisms'' for recursively defined ADTs.  The interpretation of recursive ADTs as least fixed points of functors, which are the object components of initial algebras, guarantees that for any ADT $A$ interpreted as $LFP(a)$,
\[ \semof{A}_\Gamma \simeq a(\semof{A}_\Gamma)\]
in $\mcC$.  Lifting the above isomorphism to a syntactic isomorphism, that is, an equality in $\ADTrig$, gives us the by now familiar rig isomorphism presented in Figure \ref{fig:adt_isos}: for all ADTs $A \in ADT_{\vars, X}$,
\[ \mu X.A \simeq A[\mu X.A/X] \]
The above isomorphism is the most obvious of the syntactic isomorphisms we call mu rules in the definition below.

We have seen that we can use the tools of algebra to reason about the isomorphisms induced by fixed point isomorphisms.  In the case of closed recursive polynomial types, if $p \in (\Sigma_{Rig})^*_X$, then
\[ \mu X.p(X) \simeq p(\mu X.p(X)) \]
That is, we have an equality $A = p(A)$ for the appropriate element $A$ of the semiring $\ADTrig$.  From here we can use semiring theory (and, as we have seen, ring theory) to reason about type isomorphisms.  In this chapter, we address several new questions: what other isomorphisms exist between interpretations of ADTs in $\mcC$?  Phrased differently, what mu rules exist?  Can we use the tools of algebra to reason about the type isomorphisms induced by further mu rules?

\begin{definition}
    A \emph{mu rule} is an inference rule of the form 
    \[ \inferrule{P_1 \quad \cdots \quad P_k}{L \simeq R}\]
    with the following properties, for any category $\mcC$ interpreting ADTs:
    \begin{enumerate}[nosep]
        \item $L,R \in ADT_\vars$ for some set $\vars$.
        \item $P_i$ is either an equation between ADTs, an ADT, or another parameter to the rule (e.g., $X \in \vars$).
        \item An equation $A \simeq B$ is interpreted as asserting the existence of a natural isomorphism between the functors $\semof{A}_{(-)}$
        and $\semof{B}_{(-)}$ of type $\mcC^\vars \to \mcC$.
        \item If the premises above are true, then the conclusion is also true.
        \item The conclusion is not derivable from the isomorphisms corresponding to the rig axioms alone.
    \end{enumerate}
\end{definition}

In the remainder of this chapter, we introduce several mu rules.  The propositions in this chapter prove the validity of mu rules that are given as inference rules immediately following them. (However, we eschew proofs of naturality except when it is needed elsewhere).  The proofs of these propositions arise from the results given in Chapter \ref{sec:catprelim}.  We give an example of each mu rule's in practice in deriving a categorical isomorphism between interpretations of ADTs.  We conclude this chapter with a discussion of how algebra may be used to decide the existence of type isomorphisms arising from the mu rules we introduce.

\section{Substitution Lemmas}
We will need the following lemmas for several of the propositions in this chapter in which we prove the soundness of mu rules.

% We will need the following notation and lemmas for several of the propositions in this chapter in which we prove mu rules.

% \begin{notation}
%     Let $\theta : \vars \to ADT_\wars$.  Let $A : ADT_\vars$. Then we write $A[\theta]$ to mean the substitution of all variables in $A$ for their image under $\theta$.  That is, if $\vars = \{v_1, ..., v_n\}$, then
%     \[ A[\theta] = A[\theta(v_1)/v_1]...[\theta(v_n)/v_n]\]
%     Note that $A[\theta] \in \wars$.
% \end{notation}
% \begin{lemma}
%     Given any ADT $A \in ADT_{\vars}$ and $\theta : \vars \to ADT_{\wars}$
%     Then there exists an isomorphism natural in $\Gamma : \wars \to \mcC$ 
%     \[ \semof{A}_{\semof{\theta(-)}_\Gamma} \simeq \semof{A[\theta]}_\Gamma\]
% \end{lemma}

% \begin{proof}
%     \begin{description}
%         \item[Least fixed point]
%         Suppose $A = \mu X. A'$. We can assume $X \notin \wars$, performing $\alpha$-conversion if necessary.

%         Then $A[\theta] = \mu X. A'[\theta]$.

%          By functoriality of $LFP$, it is sufficient to construct an isomorphism, natural in $Z$.
%          \[ \semof{A'}\]
         
%     \end{description}
% \end{proof}

\begin{lemma} \label{lemma:semofsub}
    Given any ADTs $A \in ADT_{\vars,X}$ and $B \in ADT_{\vars,X}$ and any environment $\Gamma : \vars \cup \{X\} \to \mcC$,
    there exists an isomorphism, natural in $\Gamma$
    \[ \semof{A}_{\Gamma[X:=\semof{B}_\Gamma]} \simeq \semof{A[B/X]}_\Gamma \]
\end{lemma}

\begin{proof}
    By induction on the structure of $A \in ADT_{\vars,X}$, we prove the following statement:
    \begin{equation}
    \label{eq:sublemma}
        \forall B.\ \forall \Gamma.\  \semof{A}_{\Gamma[X:=\semof{B}_\Gamma]} \simeq \semof{A[B/X]}_\Gamma
    \end{equation}
    and each isomorphism is natural in $\Gamma$.
    \begin{description}
    \item[$A = \mb{1}$, $A = \mb{0}$.] The bases cases where $A = \mb{1}$ or $A = \mb{0}$ are trivial.
    \item[$A \in \vars$.]  If $A \neq X$, then
    \[ \semof{A}_{\Gamma[X:=\semof{B}_\Gamma]} = \Gamma A = \semof{A[B/X]}_\Gamma \]
    If $A = X$, then
    \[ \semof{A}_{\Gamma[X:=\semof{B}_\Gamma]} = \semof{B}_\Gamma = \semof{A[B/X]}_\Gamma \]
    In both cases, functoriality comes trivially from the fact that $\semof{D}_-$ is a functor for all $D$ and both isomorphisms are identity isomorphisms.
    \item[$A = A_1 \times A_2$.]  Assume \eqref{eq:sublemma} holds for $A_1$ and $A_2$.
    Then the following holds by the congruence of $\simeq$ with respect to $\times$ together with the inductive hypothesis.
    \[ \semof{A}_{\Gamma[X:=\semof{B}_\Gamma]} = \semof{A_1}_{\Gamma[X:=\semof{B}_\Gamma]} 
    \times \semof{A_2}_{\Gamma[X:=\semof{B}_\Gamma]} \simeq \semof{A_1[B/X]}_\Gamma \times \semof{A_2[B/X]}_\Gamma = \semof{A[B/X]}_\Gamma \]
    Naturality can be derived by considering the above as products of $\semof{A_1}_{(-)}$ and $\semof{A_2}_{(-)}$ in the functor category.

    Being a bifunctor on the functor category, $\times$ acts on natural transformations to produce natural transformations.
    \item[$A = A_1 + A_2$.]  Assume \eqref{eq:sublemma} holds for $A_1$ and $A_2$.
    Then the following holds by the congruence of $\simeq$ with respect to $+$ together with the inductive hypothesis.
    \[ \semof{A}_{\Gamma[X:=\semof{B}_\Gamma]} = \semof{A_1}_{\Gamma[X:=\semof{B}_\Gamma]} 
    + \semof{A_2}_{\Gamma[X:=\semof{B}_\Gamma]} \simeq \semof{A_1[B/X]}_\Gamma + \semof{A_2[B/X]}_\Gamma = \semof{A[B/X]}_\Gamma \]
    Similar to the previous case, naturality is derived from considering $+$ as a bifunctor on the functor category.
    
    \item[$A = \mu Y.A'$, $A' \in ADT_{\vars,X,Y}$.]  Assume that
    \[ \forall B'.\ \forall \Gamma'.\ \semof{A'}_{\Gamma'[X:=\semof{B'}_{\Gamma'}]} \simeq \semof{A'[B'/X]}_{\Gamma'} \]
    Given $B$ and $\Gamma$, we want to show
    \[ \semof{\mu Y.A'}_{\Gamma[X:=\semof{B}_\Gamma]} \simeq \semof{\mu Y.A'[B/X]}_\Gamma \]
    or
    \[ LFP(\lambda Z.\semof{A'}_{\Gamma[X:=\semof{B}_\Gamma][Y:= Z]}) \simeq LFP(\lambda Z.\semof{A'[B/X]}_{\Gamma[Y:= Z]}) \]
    By the functoriality of $LFP$, it is sufficient to show that the following isomorphism exists and is natural in $Z$.
    \begin{equation}
        \label{eq:subst-lemma-ih}
    \semof{A'}_{\Gamma[X:=\semof{B}_\Gamma][Y:= Z]} \simeq \semof{A'[B/X]}_{\Gamma[Y:= Z]} 
    \end{equation} 
    Apply the induction hypothesis with $B' = B$ and $\Gamma' = \Gamma[Y:= Z]$ to obtain the following isomorphism which is natural in $\Gamma'$
    \[ \semof{A'}_{\Gamma[Y:= Z][X:=\semof{B}_\Gamma]} \simeq \semof{A'[B/X]}_{\Gamma[Y:=Z]} \]
    Because $Y \neq X$, we can commute the updating of $Y$ and of $X$ in $\Gamma$ to obtain exactly equation \eqref{eq:subst-lemma-ih}.
    The above isomorphism is natural in $\Gamma'$, so it is natural in $Z$.  Applying $LFP$, we attain the desired result.  Naturality follows from the functoriality of $LFP$. Specifically, the naturality square for the constructed isomorphism is exactly the image, under the functorial action of $LFP$, of the naturality square for the isomorphism obtained by the inductive hypothesis in \eqref{eq:subst-lemma-ih}. \qedhere
    \end{description}
\end{proof}
\begin{corollary}\label{cor:semofsub2}

    Given any ADTs $A, B \in ADT_{\vars,X}$, any environment $\Gamma : \vars \to \mcC$, and any object $D \in \mcC$
    \[ \semof{A}_{\Gamma[X:=\semof{B}_{\Gamma[X := D]}]} \simeq \semof{A[B/X]}_{\Gamma[X := D]} \]
\end{corollary}
\begin{proof}
    Let $\Gamma' = \Gamma[X:= D]$.
    
    Observe that
    \[ \Gamma[X := \semof{B}_{\Gamma'}] = \Gamma'[X:= \semof{B}_{\Gamma'}] \]
    By Lemma \ref{lemma:semofsub} with $\Gamma = \Gamma'$, we have our desired result.
\end{proof}

\begin{lemma}\label{lemma:doublesub}
    Let $A, B, C \in ADT_{\vars, X}$. Then
    \[ A[B/X][C/X] = A[B[C/X]/X] \]
\end{lemma}
\begin{proof}
    We proceed by induction on $A$.
    The base cases where $A = \mb{0}$ or $A = \mb{1}$ are trivial.
    Let $A \in \vars \cup \{X\}$.
    If $A = X$, then
    \[ A[B/X][C/X] = B[C/X] = A[B[C/X]/X] \]
    If $A \neq X$ then
    \[ A[B/X][C/X] = A = A[B[C/X]/X] \]
    Let $A = A_1 + A_2$.  Assume 
    \[ A_i[B/X][C/X] = A_i[B[C/X]/X] \]
    Then
    \begin{align*}
        A[B/X][C/X] &= A_1[B/X][C/X] + A_2[B/X][C/X] \\
        &= A_1[B[C/X]/X] + A_2[B[C/X]/X] \\
        &= A[B[C/X]/X]
    \end{align*}
    Let $A = A_1 \times A_2$.  Assume 
    \[ A_i[B/X][C/X] = A_i[B[C/X]/X] \]
    Then
    \begin{align*}
        A[B/X][C/X] &= A_1[B/X][C/X] \times A_2[B/X][C/X] \\
        &= A_1[B[C/X]/X] \times A_2[B[C/X]/X] \\
        &= A[B[C/X]/X]
    \end{align*}
    Let $A = \mu Y.A'$.  Assume
    \[ A'[B/X][C/X] = A'[B[C/X]/X] \]
    Then
    \[ A[B/X][C/X] = \mu Y.A'[B/X][C/X] = \mu Y.A'[B[C/X]/X] = A[B[C/X]/X] \qedhere \]
\end{proof}

\section{Iteration Rule}

From here forward, to state several of the following mu rules, we make use of the following notation.

\begin{notation}
    Let $A, B \in ADT_{\vars, X}$ be ADTs.  We write
    $A[B/X]^n$ to denote $n$-fold repeated substitution of $B$ for $X$ in $A$.  That is,
    \[ A[B/X]^n = A\underbrace{[B/X][B/X]...[B/X]}_{\text{$n$ times}} \]
\end{notation}

The first mu rule directly corresponds to Theorem \ref{theorem:cofinal}. We can now state the first of the mu rules.

\begin{prop}\label{prop:sofs}
    Given any ADT $S \in ADT_{\vars,X}$ and any environment $\Gamma : \vars \to \mcC$, for all $n$,
    \[ \semof{\mu X.S}_\Gamma \simeq \semof{\mu X. S[S/X]^n}_\Gamma \]
\end{prop}
\begin{proof}
    By the definition of $\semof{-}$, the above equation simplifies to
    \[ LFP(\lambda W. \semof{S}_{\Gamma[X:= W]}) \simeq LFP(\lambda W. \semof{S[S/X]^n}_{\Gamma[X:=W]}) \]
    By $n$-fold repeated application of Lemma \ref{lemma:doublesub}, we get
    \begin{equation}\label{eq:sofs1}
        \mu X.S[S/X]^n = \mu X.S\underbrace{[S[S[...]/X]/X]}_{\text{$n$ times}}
    \end{equation}
    By Corollary \ref{cor:semofsub2}, we have
    \begin{equation*}
        LFP(\lambda W. \semof{S\underbrace{[S[S[...]/X]/X]}_{\text{$n$ times}}}_{\Gamma[X:=W]}) \simeq LFP(\lambda W. \semof{S}_{\Gamma[X:=\semof{\underbrace{ S[S[S[...]/X]/X]}_{\text{$n-1$  times}}}_{\Gamma[X:=W]}]})
    \end{equation*}
    Let $s : \mcC \to \mcC$ be defined by
    \[ s(Z) = \semof{S}_{\Gamma[X := Z]} \]
    Then repeated application of Corollary \ref{cor:semofsub2} derives
    \begin{equation}\label{eq:sofs2}
        LFP(\lambda W. \semof{S \underbrace{[S[S[...]/X]/X]}_{\text{$n$ times}}}_{\Gamma[X:=W]}) \simeq LFP(s^n) 
    \end{equation}
    By Corollary \ref{cor:omegacofinal} with $f$ defined by $f(x) = nx$, 
    \begin{equation}\label{eq:sofs3}
        LFP(s^n) \simeq LFP(s) = LFP(\lambda W.\semof{S}_{\Gamma[X := W]}) = \semof{\mu X.S}_\Gamma
    \end{equation}
    Combining equations \eqref{eq:sofs1},  \eqref{eq:sofs2}, and \eqref{eq:sofs3}, we obtain the desired result.
\end{proof}

Proposition \ref{prop:sofs} gives rise to the following mu rule, which is valid for ADTs but is not generated by isomorphisms corresponding to the rig axioms.
We adjoin the following inference rule to the definition of equality in $\ADTrig$.
\begin{equation}\label{eq:sofsinfer}
    \inferrule{S \in ADT_{\vars,X} \qquad n \in \nat}{\mu X.S[S/X]^n \simeq \mu X.S}
\end{equation}

\begin{example}
    Consider the ADT $B$ defined by
    \[ B = \mu X.\mb{1} + X^2 \]
    Let $B' = \mb{1} + X^2$.
    Proposition \ref{prop:sofs} tells us that $B \simeq \mu X.B'[B'/X]^n$.  If, e.g., $n = 1$, then we have that
    \[ B = \mu X. \mb{1} + X^2 \simeq \mu X.\mb{1} + (\mb{1} + X^2)^2\]
    If $n = 2$, then we have that
    \[ B = \mu X. \mb{1} + X^2 \simeq \mu X.\mb{1} + (\mb{1} + (\mb{1} + X^2)^2)^2\] 
\end{example}

\section{Diagonal Rule}

The second mu rule directly corresponds to Proposition \ref{prop:iteratedlfp} from Chapter \ref{sec:catprelim}.  It is stated as follows.

\begin{prop}\label{prop:muxmuy}
    Given any ADT $S \in ADT_{\vars, X, Y}$ and any environment $\Gamma : \vars \to \mcC$,
    \[ \semof{\mu X.\mu Y.S}_{\Gamma} \simeq \semof{\mu X. S[X/Y]}_\Gamma \]
\end{prop}
\begin{proof}
    Simplifying the above equation using the definition of $\semof{-}$, we find that it suffices to show
    \begin{equation*}\label{eq:muxmuy1}
        LFP(\lambda W.LFP(\lambda Z.\semof{S}_{\Gamma[X:=W][Y:=Z]})) \simeq LFP(\lambda W.\semof{S[X/Y]}_{\Gamma[X:=W]})
    \end{equation*}
    Combining Corollary \ref{cor:semofsub2} with Proposition \ref{prop:iteratedlfp}, we have our desired result. \qedhere
\end{proof}

We lift Proposition \ref{prop:muxmuy} to the syntax of ADTs by adjoining the following inference rule to the definition of equality in $\ADTrig$.
\begin{equation}\label{eq:muxmuysyntactic}
    \inferrule{S \in ADT_{\vars, X, Y}}{\mu X.\mu Y.S \simeq \mu X.S[X/Y]}
\end{equation}

\begin{example}
Consider the ADT $A$ defined by
\[ A = \mu X. \mu Y. \mb{1} + Y + X^2 \]
$A$ can be thought of as a type $B$ defined as
\[ B = \mu Y. \mb{1} + Y + A^2 \]
B can be thought of as the type
\[ \mb{1} + B + A^2 \]
That is, $B$ is either $\mb{1}$ or $B$ or $A^2$.
Using the inference rule given in \eqref{eq:muxmuysyntactic}, we get
\[ A \simeq \mu X. \mb{1} + X + X^2 \]
and $A$ becomes a much easier type to reason about.
\end{example}

\section{Unfolding rule}

Next, we derive a useful generalization of Proposition \ref{prop:sofs} that utilizes both the previously given mu rules in its proof.

\begin{prop}\label{prop:arbitraryunfoldings}
    Let $S \in ADT_{\vars,X,Y}$ be an ADT and let $\Gamma : \vars \to \mcC$ be an environment. Then
    \[ \semof{\mu X.S[X/Y]}_\Gamma \simeq \semof{\mu X.S[S/Y]^n[X/Y]}_{\Gamma}  \]
\end{prop}
\begin{proof}
    We assemble Propositions \ref{prop:muxmuy} and \ref{prop:sofs} in the following chain of isomorphisms.
    \begin{align*}
        \semof{\mu X.S[X/Y]}_\Gamma
        &\simeq \semof{\mu X.\mu Y.S}_\Gamma && \textsf{(By Proposition \ref{prop:muxmuy})}\\
        &\simeq \semof{\mu X.\mu Y.S[S/Y]^n}_\Gamma && \textsf{(By Proposition \ref{prop:sofs})}\\
        &\simeq \semof{\mu X.S[S/Y]^n[X/Y]}_\Gamma  && \textsf{(By Proposition \ref{prop:muxmuy})}
    \end{align*}
\end{proof}

We lift Proposition \ref{prop:arbitraryunfoldings} to the syntax of ADTs by adjoining the following inference rule to the definition of equality in $\ADTrig$.
\begin{equation}\label{eq:arbitraryunfoldings}
    \inferrule{S \in ADT_{\vars, X, Y} \qquad n \in \nat}{\mu X.S[X/Y] \simeq \mu X.S[S/Y]^n[X/Y]}
\end{equation}

The above proposition tells us that in any recursively defined type, we can ``look inside'' of any recursive variable to find the same type.  This shows us that substituting a recursive type variable inside a $\mu$ constructor with the body of the recursive type itself results in an isomorphic type.  We demonstrate this with the following example.

\begin{example}
    Consider the type $M$ defined by
    \[ M = \mu X.\mb{1} + X + X^2 \]
    Observe that $M = \mu X.M'[X/Y]$ when $M' = 1 + Y + X^2$.
    Applying Proposition \ref{prop:arbitraryunfoldings} with $n=1$ results in an isomorphism
    \[ M = \mu X.M'[X/Y] \simeq \mu X.M'[M'[X/Y]/Y] = \mu X.\mb{1} + (\mb{1} + X + X^2) + X^2\]
\end{example}

Note that using Proposition \ref{prop:arbitraryunfoldings}, we can choose arbitrary instances of the recursive variable $X$ in the definition of $M$ and replace them with the body of $M$'s $\mu$ constructor.  For example, using the same process, we can replace both the first-degree $X$ and one of the two occurrences of $X$ in $X^2$ in the definition of $M$ with $\mb{1} + X + X^2$ to obtain an isomorphic type.  To do this, we repeat the same process with $M' = \mb{1} + Y + X \times Y$.

Proposition \ref{prop:arbitraryunfoldings} is particularly useful in extending the utility of Theorem \ref{thm:mvrigiso} in the context of ADTs.  The following definition facilitates this.

\begin{definition}
    Let $p \in (\Sigma_{Rig})^*_{\{x\}}$ be an element of the free $\Sigma_{Rig}$-algebra over $\{x\}$.  The \emph{representative polynomial} of the closed recursive polynomial type $\mu X.p(X)$ is $\semof{p}_{id}^{\ints[x]}~-~[x]_{E_{Ring}}$.
\end{definition}

Observe that the representative polynomial of a closed recursive polynomial type is an element of $\ints[x]$ and not an element of $\nat[x]$.

Notice that for each $p_i$ in the set of polynomials $p_1, ..., p_n$ to which Theorem \ref{thm:mvrigiso} is applied, $p_i - x_i$ must be monic.  Going forward, we extend the concept of monicity to closed recursive polynomial types.

\begin{definition}
    A closed recursive polynomial type $\mu X.p(X)$ is \emph{monic} if its representative polynomial is monic.
\end{definition}

For ADTs, Theorem \ref{thm:mvrigiso} is restricted to generating isomorphisms between closed recursive polynomial types $\mu X.p_i(X)$ whose representative polynomial has a leading coefficient of $1$.  Proposition \ref{prop:arbitraryunfoldings} can generate isomorphisms between certain closed recursive polynomial types that are not monic, and ones that are.  We begin by introducing the following definition.

\begin{definition}
    The \emph{content} of a polynomial in $\ints[x]$ is the greatest common divisor of its coefficients.  The \emph{content} of a closed recursive polynomial type is the content of its representative polynomial.
\end{definition}

Observe that we can divide a polynomial $p \in \ints[x]$ by its content to get a result in $\ints[x]$ because every coefficient in $p$ has the content of $p$ as a factor.  Some polynomials become monic when they are divided by their content.  For example,
\[ \frac{3x^3 + 6x^2 + 24x + 9}{3} = x^3 + 2x^2 + 8x + 3\]
If we have a set of closed recursive polynomial types whose representative polynomials become monic when they are divided by their content, then we can apply \ref{prop:arbitraryunfoldings} to ``reduce their contents'' such that they become monic.  Once we do this, we can then apply Theorem \ref{thm:mvrigiso} to decide type isomorphisms between them.  We provide an illustrative example below.
\begin{example}
    Let 
    \begin{align*}
        T_2 &= \mu X.\mb{2} + X + \mb{2}X^3 \\
        T &= \mu X.\mb{1} + X + X^3 \\
        M_3 &= \mu X.\mb{3} + X + \mb{3}X^2 \\
        M &= \mu X. \mb{1} + X + X^2
    \end{align*}

    
    The content of $T_2$ is $2$ and the content of $M_3$ is $3$.  Let $t$ be the representative polynomial of $T_2$ and $m$ be the representative polynomial of $M_3$.  Proposition \ref{prop:arbitraryunfoldings} lets us generate an isomorphism between $T_2$ and $T$ and between $M_3$ and $M$, that is, it lets us reduce the content of $T_2$ and of $M_3$ to $1$.

    Using the isomorphisms corresponding to the rig axioms and the congruence law for $\mu$, we can generate the isomorphism
    \[ T_2 \simeq \mu X. \mb{1} + (\mb{1} + X + X^3) + X^3\]

    Using Proposition \ref{prop:arbitraryunfoldings} with $S = \mb{1} + Y + X^3$ and $n = 1$ together with the isomorphism above, we can generate the isomorphism
    \[ T_2 \simeq \mu X.\mb{1} + X + X^3\]

    Similarly, using the rig axioms and Proposition \ref{prop:arbitraryunfoldings} with $S = \mb{1} + Y + X^3$ and $n = 2$, we can generate the isomorphism
    \[ M_3 \simeq \mu X. \mb{1} + X + X^2\]

    Now that we have reduced the content of $T_2$ and of $M_3$, we can use Theorem \ref{thm:mvrigiso} to decide the existence of type isomorphisms between them.  For example, we can generate an isomorphism
    \[ T_2M_3 \simeq TM \simeq T^4M^3 \simeq T_2^4M^3_3 \]
    because when we let $r_t$ be an arbitrary root of $t$ and $r_m$ be an arbitrary root of $m$, $r_tr_m = r_t^4r_m^3$.
\end{example}

Note that the above process for reducing the content of a closed recursive polynomial type to $1$ generalizes to any closed recursive polynomial type with content greater than or equal to $1$.

\section{Rolling Rule}

The last mu rule we will discuss in this chapter directly corresponds to Theorem \ref{theorem:FGcolimit} is given as follows.  

\begin{prop}\label{prop:adtrolling}
Let $S,T \in ADT_{\vars, X}$.  Then for all environments $\Gamma : \vars \to \mcC$
\[ \semof{\mu X.S[T/X]}_\Gamma \simeq \semof{S[\mu X.T[S/X]/X]}_\Gamma \]
\end{prop}
\begin{proof}
    Let $s$ be the functor defined by
    \[ s(Y) = \semof{S}_{\Gamma[X:= Y]}\]
    and $t$ be the functor defined by
    \[ t(Y) = \semof{T}_{\Gamma[X:= Y]}\]
    We can now prove the proposition using the following chain of isomorphisms
    \begin{align*}
        \semof{\mu X.S[T/X]}_\Gamma 
        &= LFP(\lambda W. \semof{S[T/X]}_{\Gamma[X := W]}) && \textsf{(By definition of $\semof{-}$)}\\
        &\simeq LFP(\lambda W. \semof{S}_{\Gamma[X:= \semof{T}_{\Gamma[X:= W]}]} && \textsf{(By Corollary \ref{cor:semofsub2})} \\
        &= LFP(s \circ t)  && \textsf{(By the definitions of $s$ and of $t$)} \\
        &\simeq s(LFP(t \circ s)) && \textsf{(by Theorem \ref{theorem:FGcolimit})} \\
        &= \semof{S}_{\Gamma[X := LFP(\lambda W.\semof{T}_{\Gamma[X:= \semof{S}_{\Gamma[X := W]}]})]} && \textsf{(By the definitions of $s$ and of $t$)} \\
        &\simeq \semof{S}_{\Gamma[X := LFP(\lambda W.\semof{T[S/X]}_{\Gamma[X := W]})]} && \textsf{(By Corollary \ref{cor:semofsub2})} \\
        &\simeq \semof{S}_{\Gamma[X := \semof{\mu X.T[S/X]}_{\Gamma}]} && \textsf{(By definition of $\semof{-}$)} \\
        &\simeq \semof{S[\mu X.T[S/X]/X]}_\Gamma && \textsf{(By Lemma \ref{lemma:semofsub})} 
    \end{align*}\qedhere
\end{proof}

We lift Proposition \ref{prop:adtrolling} to the equalities in $\ADTrig$ given by the following inference rule.
\begin{equation}\label{eq:adtrolling}
    \inferrule{S,T \in ADT_{\vars, X}}{\mu X.S[T/X] \simeq S[\mu X.T[S/X]/X]}
\end{equation}

The following example illustrates the use of the mu rule defined by Proposition \ref{prop:adtrolling}.

\begin{example}
    Let $S = \mb{1} + X$ and $T = X^2$.  Applying the inference rule given in equation \eqref{eq:adtrolling}, we get
    \[ \mu X.S[T/X] = \mu X.\mb{1} + X^2 \simeq \mb{1} + \mu X.(\mb{1} + X)^2 = S[\mu X.T[S/X]/X]\]
    Further simplifying using the rig axioms, we get
    \[ \mu X.\mb{1} + X^2 \simeq \mb{1} + \mu X.\mb{1} + \mb{2}X + X^2\]
    Notice that we can also apply equation \eqref{eq:adtrolling} with $S = T$ and $T = S$ to obtain
    \[ \mu X.(1 + X)^2 \simeq (\mu X.\mb{1} + X^2)^2 \]
    Again simplifying using the rig axioms, we get
    \[ \mu X.\mb{1} + \mb{2}X + X^2 \simeq (\mu X.\mb{1} + X^2)^2\]
\end{example}

We give a complete account of the mu rules introduced in this chapter in Figure \ref{fig:murules}.

\begin{figure}
\fbox{
  \begin{minipage}{0.9\linewidth}
    \begin{mathpar}
      \inferrule{S \in ADT_{\vars,X} \qquad n \in \nat}{\mu X.S[S/X]^n \simeq \mu X.S}\eqref{eq:sofsinfer} \qquad
        \inferrule{S \in ADT_{\vars, X, Y}}{\mu X.\mu Y.S \simeq \mu X.S[X/Y]}\eqref{eq:muxmuysyntactic} \\
        \inferrule{S \in ADT_{\vars, X, Y} \qquad n \in \nat}{\mu X.S[X/Y] \simeq \mu X.S[S/Y]^n[X/Y]}\eqref{eq:arbitraryunfoldings} \qquad
        \inferrule{S,T \in ADT_{\vars, X}}{\mu X.S[T/X] \simeq S[\mu X.T[S/X]/X]}\eqref{eq:adtrolling} \\
    \end{mathpar}
  \end{minipage}
}
\caption{The inference rules that correspond to the mu rules introduced in Chapter \ref{sec:murules}}\label{fig:murules}
\end{figure}


\section{The algebra of mu rules}
Recursively defined ADTs' fixed point isomorphisms induce many additional isomorphisms in the semiring $\ADTrig$.  As we saw in Chapter \ref{sec:ADTs}, determining when such isomorphisms exist is far from trivial, requiring tools from algebra, including both ring theory and semiring theory.  Much like the isomorphisms induced by fixed point isomorphisms, the equalities in $\ADTrig$ introduced in this chapter also generate a rich collection of further equalities.  Many of these can be effectively studied using algebraic techniques, particularly in the case of closed recursive polynomial types.  

As an example, Proposition \ref{prop:sofs} and its associated inference rule \eqref{eq:sofsinfer} show that for any $p \in (\Sigma_{Rig})^*_{\{x\}}$, the type $P$ defined as $\mu X.p(X)$ is not only isomorphic to $p(P)$, but also to each type $P_n$ defined as $\mu X.p^n(X)$ for $n \ge 1$.  Each $P_n$ is, in turn, isomorphic to $p^n(P_n)$.  While this rule may appear unremarkable in one of its directions because $P \simeq p(P)$ already implies that $P \simeq p^n(P)$, it can be used to discover type isomorphisms when applied in the opposite direction.  For instance, it is not immediately evident that a type like $A$ defined by $A = \mu X. \mb{4} + \mb{8} X + \mb{8} X^2 + \mb{4} X^3 + X^4$ is isomorphic to $\mu X.\mb{1} + \mb{2}X + X^2$.  As in Chapter \ref{sec:ADTs}, such phenomena can be analyzed using algebra. We can study when a closed recursive polynomial type admits a simpler presentation using equation \eqref{eq:sofsinfer} using the algebraic concept of Chebyshev polynomials.  In fact, the only polynomials that can be written as $g (\circ g)^n$ for $n \ge 1$ are Chebyshev polynomials of the first kind and certain monomials, the latter corresponding to uninteresting closed recursive polynomial types.

By contrast, the isomorphisms arising from Proposition \ref{prop:muxmuy} (equation \eqref{eq:muxmuysyntactic}) are algebraically straightforward.  The quotient ring induced by the two fixed point isomorphisms of $\mu X.\mu Y.S$ coincides with that induced by the fixed point isomorphism for $\mu X.S[X/Y]$.  In other words, collapsing the nested $\mu$-constructors does not introduce any new algebraic structure.

The mu rule arising from Proposition \ref{prop:arbitraryunfoldings} (equation \ref{eq:arbitraryunfoldings}) is similar to that of Proposition \ref{prop:sofs}. It enables the construction of isomorphisms between closed recursive polynomial types such as $\mu X.p(X,X)$ and $\mu X.p(X,p(X, p(X, X)))$, where $p \in (\Sigma_{Rig})^*_{\{x,y\}}$.  To study this rule, we are led to the following question.  Let $q \in \nat[x]$, and for all polynomials $p \in \nat[x,y]$, define $p_x$ by $p_x(y) = p(x,y)$. Under what conditions can $q(x)$ be expressed as $p_x^n(x)$ for some $p$?

The last mu rule in this chapter, introduced in Proposition \ref{prop:adtrolling} (equation \ref{eq:adtrolling}), is perhaps the most algebraically interesting.  Restricting again to closed recursive polynomial types, this mu rule says that given any $s, t \in (\Sigma_{Rig})^*_X$, there exist types $A$ defined as $\mu X.s(t(X))$ and $B$ defined as $\mu X.t(s(X))$ such that $t(A) \simeq B$ and $s(B) \simeq A$.  The resulting isomorphisms can be studied using methods similar to those of Chapter \ref{sec:ADTs}.  In particular, we can consider the quotient ring $\nat[x,y]/(y = t(x), x = s(y))$.  There may be certain properties of this semiring that can be exploited to generate proofs of equality which correspond to isomorphisms in $\ADTrig$ similar to how high polynomials and their properties were used in Chapter \ref{sec:ADTs}.

Taken together, the mu rules introduced in this chapter highlight a range of algebraic phenomena that arise in $\ADTrig$, particularly in the setting of closed recursive polynomial types.  These rules suggest a rich and as of yet undiscovered algebraic structure in the isomorphisms they generate.  Determining exactly which isomorphisms they induce remains a nontrivial problem.  


\chapter{Conclusion and Directions for Future Work}

This thesis developed a framework for understanding type isomorphisms between algebraic data types (ADTs) through the combined use of category theory and algebra.  Recursive types are interpreted as least fixed points and, following Marcelo Fiore and Tom Leinster, their behavior is analyzed in quotient semirings.  In this setting, type equivalence is treated as an algebraic problem and a wide range of tools become available to us.  By treating types simultaneously as categorical objects and as elements of algebraic structures, algebraic identities can be systematically translated into concrete type isomorphisms, bridging symbolic reasoning and semantic equivalence.

The central contribution of this work is a substantial extension of the results of Marcelo Fiore and Tom Leinster.  Whereas their work addresses type isomorphisms involving a single closed recursive polynomial type, Theorems \ref{thm:mvrigiso} and \ref{thm:mvringtorig} generalize this to \textit{multiple distinct recursive types}.  These results show that, under certain conditions, equalities between polynomials interpreted over $\mathbb{C}$ can be systematically lifted to isomorphisms between ADTs.  Theorems \ref{thm:mvrigiso} and \ref{thm:mvringtorig} provide a mechanism for generating nontrivial isomorphisms that involve interactions between several recursive types simultaneously.  In doing so, this thesis identifies a previously unknown structure inherent in the theory of ADTs and demonstrates that this structure can be used to lift numeric equalities to type isomorphisms in a principled way.

Beyond its theoretical significance, this work has direct implications for practical settings in programming languages and type theory.  Type isomorphisms are important in program transformations, data representation changes, and interpreter optimizations: two isomorphic types can be substituted without affecting the behavior of a program. The results presented here expand the class of isomorphisms that can be systematically derived to the ones described in Chapter \ref{sec:ADTs} involving multiple different recursive data types.  Moreover, by expanding the class of types that have isomorphisms derivable using algebra, we have strengthened the prospects for automating reasoning about programs.  Algebraic tools such as ri(n)g solvers could identify equivalences in this wider setting, and the constructions developed here could then be used to synthesize the corresponding isomorphisms.  

These results leave several problems open.  The mu rules introduced in this thesis enlarge the space of provable isomorphisms, but do not yet generate a complete calculus.  It remains unclear whether a full characterization of type isomorphisms for recursive ADTs is possible, or whether the space of such isomorphisms is inherently too rich to admit a finite presentation.  Similarly, while the quotient structures induced by these rules clearly encode meaningful algebraic information, their properties are not yet well understood and call for further investigation.  Given that many of the mu rules can be studied in the context of polynomial rings under a quotient, further progress will likely depend on deeper connections with ring and semiring theory.

Looking forward, two directions of research appear especially promising.  First, studying the free category generated by the structure required to interpret ADTs could provide a universal setting in which such isomorphisms arise, which would allow general results (such as Theorem \ref{thm:mvrigiso}) to be transported to concrete models.  Second, formalization in proof assistants such as Agda or Lean would eliminate the need for intricate manual constructions.  The complexity of results such as Theorem \ref{thm:mvrigiso} makes explicitly deriving the corresponding isomorphisms extremely difficult in practice.  A formalization of this result would instead allow these constructions to be generated automatically.

Ultimately, this thesis builds on the perspective introduced by Marcelo Fiore and Tom Leinster by reinforcing the view that type isomorphisms of recursive data types are best understood as manifestations of underlying algebraic structure.  The results of this thesis show that this perspective scales beyond the single-type setting to encompass interactions between multiple recursive types, revealing a yet richer algebraic structure governing their equivalences.  In this way, algebra is not merely a language for describing types, but a powerful method of systematic reasoning about the construction of their isomorphisms.

\bibliography{citations}

\begin{vita}
Ben Lenox was born in Raleigh, North Carolina.  He graduated from Panther Creek High School in June 2019.  He enrolled in Wake Technical Community College in Raleigh, North Carolina in January 2020, where he received an Associate of Applied Sciences Degree in Software Development in May 2022.  He then transferred to Appalachian State University in Boone, North Carolina, where he received the degree of Bachelor of Science in Computer Science in December 2024.  During his undergraduate studies, he worked as a research assistant with Dr. Andrew Polonsky, who later advised him in his Master's studies.  Ben earned a Master of Science in Computer Science from Appalachian State University in 2026.  Afterward, he went on to study at Karlsruhe Institute of Technology in Germany under the supervision of Professor Dr. Sebastian Erdweg in pursuit of a PhD.
\end{vita}
%%%%%%%%%%%%%%%%%%%%%%%%%%%%%%%%%%%%%%%%%%%

\end{document}  

